

\chapter{Models for symmetric monoidal \texorpdfstring{$(\oo,1)$}-categories}
\label{ch2:: title}

The goal of this section is to introduce a model for $(\oo,1)$-categories that complements the existing literature. The motivation here is geometric: most models of $(\oo,1)$-categories are not necessarily tailored for the construction of geometric moduli spaces. The rare exception is the recent work of Ayala--Francis--Rozenblyum \cite{AFR1}, where the authors present \emph{striation sheaves} — a very bespoke class of constructible sheaves on stratified spaces, designed exactly with the intent of constructing $(\oo,1)$-categories of a moduli-flavour using point-set data. Our stated intent is to construct an $(\oo,1)$-category of stratified cobordisms, and to sketch its $(\oo,n)$-version; doing so with complete Segal spaces proves to be fraught with complications.


In \Cref{ch2:: sec: complete_segal_spaces} we introduce the construction of $(\oo,1)$-categories using simplicial spaces i.e., presheaves of simplicial sets, as a way of motivating our construction using open intervals — namely we introduce \emph{complete Segal spaces}.  In \Cref{ch2:: subsec: interval_spaces} we introduce \emph{interval spaces} and provide some intuition for how to think about them. In \Cref{ch2:: subsec: interval_spaces_model_infinity_cats} we prove that they offer a model of $(\oo,1)$-categories, by proving the existence of an explicit Quillen equivalence of model categories, and we discuss important homotopy-categorical invariants associated to interval Segal spaces, like the homotopy category. Lastly, in \Cref{ch2:: sec: symmetric monoidal structure} we discuss symmetric monoidal structures for $(\oo,1)$-categories and construct a model structure whose fibrant objects are symmetric monoidal Segal interval spaces. The ideas in this chapter, namely the reformulation of complete Segal spaces using the interval category, are based on unpublished work by Manuel Araújo and Pedro Boavida de Brito.

\section{Preliminaries on simplicial homotopy theory}
\label{ch2:: sec: preliminaries on simplicial homotopy theory}

We will introduce fundamental objects, definitions, and results and a few important notation conventions which we will stick to throughout the remainder of the thesis.


\subsection{The simplicial indexing category}

Throughout the thesis we will use simplicial sets as avatars or models for spaces, so we begin by introducing them.

\begin{definition}\label{ch2:: def: simplex category}
    We define the \textbf{simplex category} $\bDelta$ as the category whose:
    \begin{itemize}
        \item Objects are finite linearly ordered sets of the form 
        \begin{equation}
        [n] = \{0 < 1 < \cdots < n\}
        \end{equation} for $n\geq 0$.
        \item A morphism from $[n]$ to $[m]$ is a weakly order-preserving function
        \begin{equation}
            \alpha \maps [n] \longra [m].
        \end{equation} That is, a function such that 
        \begin{equation}
            \alpha(i) \leq \alpha(j)
        \end{equation}
        for every $0\leq i < j \leq n$.
    \end{itemize}
\end{definition}

\begin{observation}
    Less frequently, the simplex category is defined as the category whose objects include every finite non-empty linearly ordered set. This category is bigger and manifestly equivalent, as it is its skeleton. More can be said, since it is canonically equivalent: given a finite nonempty linearly ordered set $I$, there is a \emph{unique} weakly order preserving function $I \xra{\cong} [n]$, for some $n\geq 0$.  
\end{observation}

\begin{definition}\label{ch2:: def: join}
    The simplex category carries a (strict) non-unital monoidal structure, the \textbf{join}
    \begin{equation}
        \star \maps \bDelta \times \bDelta \longra \bDelta, 
        \qquad ([m],[n]) \longmapsto [m+n+1].
    \end{equation}
    For morphisms $\alpha \maps [m] \ra [n]$ and $\beta \maps [p] \ra [q]$, their join is the morphism
    \begin{equation}
        \alpha \star \beta \maps [m]\star [p] \longra [n] \star [q], \quad
        k \longmapsto
        \begin{cases}
            \alpha(k), & \text{if } k = 0,\ldots, m, \\
            n + \beta(k-m-1), & \text{if } k = m + 1, \ldots, m + p + 1.
        \end{cases}
    \end{equation}
    If we adjoin an initial object $[-1] = \emptyset$ to obtain the augmented simplex category $\bDelta_{+}$, then this non-unital monoidal structure extends to a unital monoidal structure with unit $[-1]$.
\end{definition}


\begin{proposition}\cite[Section~5, Proposition~1]{Mac}
    The augmented simplex category $(\bDelta_{+},\star,[-1])$ is the free monoidal category generated by a single object.
\end{proposition}


\begin{observation}
    Each $[n]$ is a finite linearly ordered set, hence a partially ordered set, hence a small category, and weakly order-preserving maps are exactly functors between these categories. Thus $\bDelta$ may be viewed as a full subcategory
    \begin{equation}
        \bDelta \longhookra \Cat.
    \end{equation}
\end{observation}

\begin{definition}\label{ch2:: def: simplicial objects}
    Let $\sC$ be a category, a \textbf{simplicial object} in $\sC$ is a functor 
    \begin{equation}
        K_\bullet \maps \bDelta^{\op} \longra \sC.
    \end{equation}
    The analog functors whose source is $\bDelta$ are called \textbf{cosimplicial objects.}
\end{definition}
\begin{notation}
    The $\bullet$ in the subscript of a simplicial object is meant as a placeholder. The image of $K$ on $[n]\in \bDelta^{\op}$ will be denoted $K_n$. We will not always be consistent with the notation $K_\bullet$ for simplicial objects, at some point we will drop it and use it only under specific circumstances.
\end{notation}
\begin{notation}
    Whenever we want to talk about the category of simplicial objects in $\cF$, we will write 
    \begin{equation}
        \mathsf{s}\sC := \fun(\bDelta^{\op},\sC).
    \end{equation}
    for the functor category from $\bDelta^{\op}$ to $\sC$. In particular, the case of simplicial sets $\sset$ will show up very frequently. When we want to emphasize that the target is already $\sset$, we will write
    \begin{equation}
        \mathsf{s}\sset := \fun(\bDelta^{\op},\sset)
    \end{equation}
    and refer to its objects as \textbf{simplicial spaces}.
\end{notation}



\begin{definition}\label{ch2:: def: face and deg maps}
    The simplex category $\bDelta$ has two important sets of distinguished morphisms. Let $n>0$, the first class $\{d^{i}\maps [n-1] \longra [n]\}_{0\leq i \leq n}$ is defined as 
    \begin{equation}
        d^{i}(j) =  
        \begin{cases}
        j & \text{ if } j < i \\
        j + 1 & \text{ if } j \geq i,    
        \end{cases}
    \end{equation} 
    and we call these the \textbf{face maps}.

    The second distinguished class of maps $\{s^{i}\maps [n+1] \longra [n]\}_{0\leq i \leq n}$ are defined as 
    \begin{equation}
        s^{i}(j) = 
        \begin{cases}
            j & \text{ if } j \leq i \\
            j - 1 & \text{ if } j > i,
        \end{cases}
    \end{equation}
    and we call these \textbf{degeneracy maps.}
\end{definition}

Note that we abuse notation when we write $d^{i}$ and $s^{i}$ without explicit reference to the source and target objects. We do this for mere notational parsimony and with the understanding that it is unlikely to hamper readability. The face and degeneracy maps satisfy some long but easy-to-check relations, which we will now enumerate for completion, together with a characterization result.

\begin{proposition} 
    The simplex category can be presented as the category whose objects are the nonempty finite totally ordered sets $[n]$, and the morphisms are generated by the face and degeneracy maps $d^{i},s^{i}$, together with the following set of relations:
    \begin{align*}
        d^{i}\circ d^{j} &= d^{j+1} \circ d^{i} \quad \text{for } i\leq j, \\
        s^{j}\circ s^{i} &= s^{i} \circ s^{j+1} \quad \text{for } i\leq j, \\
    \end{align*}
    and
    \begin{equation}
        s^{j}\circ d^{i} = 
        \begin{cases}
           d^{i} \circ s^{j-1} & \text{ if }  i < j \\
           \id_n, & \text{ if } i - j \leq 1 \\
            d^{i-1} \circ s^{j} & \text{ if } j + 1 < i,
        \end{cases}
    \end{equation}
    which we will call the \textbf{simplicial relations or identities}.
\end{proposition}

\begin{notation}
We can describe maps in $\bDelta$ using a more verbose and explicit notation. Let $f\maps [n] \longra [m]$ be an order-preserving map which sends $k \mapsto i_k$, for $0\leq k \leq n$ and $0\leq i_k \leq n$. We must have that $i_0 \leq i_1 \leq \cdots \leq i_n$, so we define 
\begin{equation}
    f \, \overset{\sf def}{=} \, \langle i_0, i_1, \cdots, i_n \rangle.
\end{equation}  
As an example, the face map $d_1 \maps [1] \longra [2]$ is written as $\langle 0,2 \rangle$ and the degeneracy map $s \maps [2] \longra [1]$ is written as $\langle 001 \rangle$, which in our opinion often improve readability. 
\end{notation}

\begin{remark}\label{ch2:: rmk: inert active}
    There is a useful alternative factorization system on $\bDelta$ — common-knowledge for users of $\oo$-operads and alike $\oo$-structures, but less so for more pragmatic users of $\oo$-categories. The two classes of morphisms are called \textbf{inert} and \textbf{active}. \emph{Inert} morphisms $\alpha \maps [m] \hookra [n]$ are those that satisfy $\alpha(i+1) = \alpha(i) + 1$ for all $i = 0,\ldots, m-1$, i.e., they are like embeddings of linear subposets. \emph{Active} morphisms $\beta \maps [m] \ra [n]$ are those that preserve extrema $\beta(0) = 0$ and $\beta(m) = n$.
\end{remark}


\subsection{Simplicial sets}

We now recall the basic comparison between simplicial sets and topological spaces. We introduce the Kan--Quillen model structure on $\sset$ and the adjunction relating it to the usual homotopy theory of spaces.

\begin{definition}\label{ch2:: def: simplicial set}
    A \textbf{simplicial set} is a simplicial object in $\set$, i.e.\ a functor
    \begin{equation}
        K_\bullet\maps \bDelta^{\op} \longra \set.
    \end{equation} 
\end{definition}

Simplicial sets are ubiquitous in modern homotopy theory. A set $X$ is a simplicial set via the \emph{constant} construction: we set $X_n = X$ for all $n$ and take all face and degeneracy maps to be the identity on $X$. Graphs can also be realised as simplicial sets, with $0$-simplices given by vertices, $1$-simplices encoding edges, and appropriate face maps recording the endpoints.

\begin{definition}\label{ch2:: def: standard simplex}
    Let $n\geq 0$. The \textbf{standard $n$-simplex}, denoted $\Delta^{n}$, is the representable functor
    \begin{equation}
        \Hom_{\bDelta}(-, [n])\maps \bDelta^{\op} \longra \set.
    \end{equation} 
    Its $k$-simplices are given by 
    \begin{equation}
        (\Delta^{n})_k = \Hom_{\bDelta}([k],[n]).
    \end{equation}
\end{definition}

\begin{observation}
    The Yoneda embedding gives a fully faithful functor
    \begin{equation}
        \bDelta \hookra \fun(\bDelta^{\op},\set)
    \end{equation} 
    sending $[n]$ to $\Delta^n$ for every $n\geq 0$, and an order-preserving $f\maps [n] \longra [m]$ in $\bDelta$ to the natural transformation $f^\ast\colon \Delta^{n}\to\Delta^{m}$, which acts on $k$-simplices by 
    \begin{align*}
        \Delta^n &\longra \Delta^m \\
        (\alpha\maps [k] \longra [n]) &\longmapsto ( f\circ \alpha \maps [k] \longra [m]).
    \end{align*}
    In particular, there is a natural isomorphism 
    \begin{equation}
        \Hom_{\sset}(\Delta^n, K) \xra{\,\,\cong\,\,} K_n
    \end{equation} 
    for every simplicial set $K$ and every $n\geq 0$.
\end{observation}

\begin{observation}
    There is a fully faithful inclusion 
    \begin{equation}\label{ch2:: def: set in ssets}
        \set \hookra \sset 
    \end{equation}
    sending a set $S$ to the constant simplicial set $([n] \mapsto S)$.
\end{observation}

\begin{definition}\label{ch2:: def: sset-mapping-space}
    The category $\sset$ is cartesian closed. Given simplicial sets $K,L\in\sset$, we denote by
    \begin{equation}
        \Map_{\sset}(K,L)
    \end{equation}
    the \textbf{simplicial mapping space} whose $k$-simplices are
    \begin{equation}
        \Map_{\sset}(K,L)_k \;\overset{\sf def}{=}\;
        \Hom_{\sset}\big(K \times \Delta^k,\, L\big),
    \end{equation}
    where $\Delta^k$ is the standard $k$-simplex. This construction makes $\sset$ enriched, tensored, and cotensored over itself.
\end{definition}

There are a few additional simplicial sets derived from the standard simplices that will be important for us.

\begin{definition}\label{ch2:: def: boundary n-simplex}
    Let $n\geq 0$. The \textbf{boundary of} $\Delta^{n}$ is the simplicial set $\partial \Delta^{n}$ defined by
    \begin{equation}
        (\partial \Delta^n)([m]) = \{\alpha \in \Hom_{\bDelta}([m],[n])\,\maps\, \alpha \text{ is not surjective}\}.
    \end{equation}
\end{definition}

\begin{definition}\label{ch2:: def: k-horn}
    Let $n > 0$ and let $0\leq k \leq n$. The \textbf{$k$-th horn of the standard $n$-simplex} is the simplicial set $\Lambda_{k}^{n}$ defined by
    \begin{equation}
        (\Lambda_{k}^{n})([m]) = \{\alpha\in \Hom_{\bDelta}([m],[n])\, : \, ([n]\setminus \{k\}) \not\subseteq \alpha([m])\}.
    \end{equation} 
    If $0<k<n$ we call $\Lambda_k^n$ an \textbf{inner horn}, and if $k\in \{0,n\}$ we call it an \textbf{outer horn}. 
\end{definition}

Both the boundary and the $k$-th horn are sub-simplicial sets of the standard $n$-simplex. In particular, the corresponding inclusions of sets give rise to monomorphisms
\begin{equation}
    \partial \Delta^n \hookra \Delta^n,
    \qquad
    \Lambda^{n}_k \hookra \Delta^n.
\end{equation}

As mentioned above, simplicial sets arise as a combinatorial model for topological spaces. We now make this comparison explicit by constructing a topological analogue of the standard $n$-simplex.

\begin{definition}\label{ch2:: def: topological n-simplex}
    Let $n\geq 0$. The \textbf{standard topological $n$-simplex} is the subset
    \begin{equation}
        \Delta^n_{t} = \Big\{(t_0,\ldots,t_n)\in \RR^{n+1}\,\maps \, t_{i} \geq 0,\ \sum_{i=0}^n t_i = 1\Big\}
    \end{equation}
    endowed with the subspace topology from $\RR^{n+1}$.
\end{definition}

\begin{proposition}
    There is a cosimplicial object in topological spaces
    \begin{equation}
        \Delta^{\bullet}_{t} \maps \bDelta \longra \Top
    \end{equation}
    which sends $[n]$ to the standard topological $n$-simplex $\Delta^n_{t}$ and sends an order-preserving map $\alpha\maps [m] \longra [n]$ to the linear map
    \begin{equation}
        \alpha(t_0,\ldots,t_m)_i = 
        \begin{cases}
            0, & \text{if } \alpha^{-1}(i) = \emptyset, \\
            \displaystyle\sum_{j\in \alpha^{-1}(i)} t_j, & \text{if } \alpha^{-1}(i) \neq \emptyset,    
        \end{cases}
    \end{equation}
    for $i=0,\ldots,n$.
\end{proposition}

\begin{definition}
    Given a simplicial set $K\maps \bDelta^{\op} \longra \set$, its \textbf{category of simplices} $\bDelta / K$ is its category of elements.
\end{definition}

Recall from \Cref{ap:: def: cat of els} that the category of simplices associated to $K$ has objects pairs $([n],x)$ with $[n]\in \bDelta^{\op}$ and $x\in K([n])$, and morphisms $([n],x)\to([m],y)$ given by order-preserving maps $\alpha \maps [n] \longra [m]$ in $\bDelta^{\op}$ such that $\alpha^\ast(y) = x$. Crucially, knowing the category of simplices of $K$ is enough to reconstruct $K$.

\begin{lemma}
    Given a simplicial set $K$, there is a natural isomorphism 
    \begin{equation}
        K \,\cong\, \mathsf{colim}\big( \bDelta / K \longra \bDelta \xra{\,\,\Delta^\bullet\,\,} \sset \big).
    \end{equation}
\end{lemma}

The previous lemma states that every simplicial set can be constructed by gluing together copies of standard simplices. This suggests how to construct a topological space from a simplicial set. The next definition makes this precise and provides the bridge between the homotopy theory of spaces and that of simplicial sets.

\begin{definition}
    Let $K$ be a simplicial set. We define its \textbf{geometric realization} $|K|$ as the following colimit of topological spaces:
    \begin{equation}
        |K| = \mathsf{colim}\big(\bDelta / K \longra \bDelta \xra{\,\,\Delta^\bullet_{t}\,\,} \Top\big).
    \end{equation}
    The construction is manifestly functorial, and we write
    \begin{equation}\label{ch2:: eq: geometric real}
        |-| \maps \sset \longra \Top
    \end{equation}
    for the associated functor.
\end{definition}

\begin{observation}
    The category of simplices $\bDelta/\Delta^{n}$ is canonically equivalent to $\bDelta/[n]$ via the Yoneda embedding. Moreover, $\bDelta/[n]$ has a terminal object, namely $[n] \xra{\id} [n]$. Together, these observations imply that 
    \begin{equation}
        |\Delta^n| \simeq \Delta^n_{t}.
    \end{equation} 
\end{observation}

Geometric realization constructs a topological space from a simplicial set. Conversely, we often wish to associate a simplicial set to a given space $X$ that captures its homotopy type. For this we use the singular complex.

\begin{definition}\label{ch2:: def: singular complex}
    Let $X$ be a topological space. We define a simplicial set $\sing(X)$ by
    \begin{equation}
        \sing(X)_n = \Hom_{\Top}(\Delta^{n}_t, X),
    \end{equation}
    and for a morphism $f\maps [m]\longra [n]$ in $\bDelta$ we set
    \begin{equation}
        \sing(X)(f) \colon \sing(X)_n = \Hom_{\Top}(\Delta^{n}_{t},X)
        \longra \Hom_{\Top}(\Delta^{m}_{t},X) = \sing(X)_m
    \end{equation}
    to be precomposition with the induced map $\Delta^m_t \to \Delta^n_t$.
    We call $\sing(X)$ the \textbf{singular complex of} $X$, and this defines a functor
    \begin{equation}
        \sing \maps \Top \longra \sset.
    \end{equation}
\end{definition}

\begin{proposition}
    The geometric realization functor is left adjoint to the singular complex functor:
    \begin{equation}
        \begin{tikzcd}[ampersand replacement=\&]
            \sset \&\& {\Top}
            \arrow[""{name=0, anchor=center, inner sep=0}, "{|-|}", curve={height=-6pt}, from=1-1, to=1-3]
            \arrow[""{name=1, anchor=center, inner sep=0}, "{\sing}", curve={height=-6pt}, from=1-3, to=1-1]
            \arrow["\dashv"{anchor=center, rotate=-90}, draw=none, from=0, to=1]
        \end{tikzcd}
    \end{equation} 
\end{proposition}

\begin{proof}
    The proof is an instance of abstract nonsense. If $K$ is a simplicial set and $X$ is a topological space, then
    \begin{align}
        \Hom_{\Top}(|K|,X)\,\, 
        &\cong \,\, \Hom_{\Top}\Big(\mathsf{colim}_{\bDelta/ K}|\Delta^{n}|, X\Big) \\
        &\cong\,\, \Limit[\bDelta / K]\Hom_{\Top}\big(|\Delta^n|,X\big) \\
        &\cong\,\, \Limit[\bDelta / K]\Hom_{\sset}\big(\Delta^n,\sing(X)\big) \\
        &\cong \,\, \Hom_{\sset}\big(K, \sing(X)\big),
    \end{align}
    where every isomorphism is natural. We have used the fact that $\Hom$ out of a colimit is a limit of $\Hom$-sets, the identification $|\Delta^n|\simeq\Delta^n_t$, and the Yoneda lemma to pass from the third line to the fourth.
\end{proof}


%%% end of subsection 

\begin{observation}\label{ch2:: obs: pi_0 simp sets}
    Just as spaces admit a functor $\pi_0 \maps \Top \longra \set$ which records its path components, so do simplicial sets:
    \begin{equation}
        \pi_0 \maps \sset \xra{\,\,\, |-|\,\,\,} \Top \xra{\,\,\, \pi_0 \,\,\,} \set.
    \end{equation}
    For $X\in \sset$, we have that $[x] = [y]$ in $\pi_0(X)$ if and only if there exists a 1-simplex $\gamma \maps \Delta^1 \ra X$ such that its two 0-faces 
    \begin{equation}
    \partial \gamma = \gamma_{|\partial \Delta^1} = \{x,y\}
    \end{equation}
    connecting the two.
\end{observation}


\begin{observation}
    If $S$ and $T$ are discrete simplicial sets, then they are weakly equivalent if and only if $\pi_k(|S|) \cong \pi_k(|T|)$ for every $k \geq 0$. All of these are equivalent to $\pi_0(|S|) \cong \pi_0(|T|)$ being bijective i.e., $S \cong T$.
\end{observation}


\subsection{A primer on model categories}
\label{subsec:model-cats-primer}

The language of model categories provides an axiomatic framework for doing
homotopy theory in settings more general than topological spaces.  Our main
references are Hirschhorn's book \cite{Hir14} and Hovey's monograph on model
categories.  In this subsection we collect the basic definitions and results
we will use later, together with a few guiding examples.



\subsubsection{Model categories and homotopy categories}


\begin{definition}
  A \emph{model category} is a complete and cocomplete category
  $\cM$ equipped with three distinguished classes of morphisms:
  \begin{itemize}
    \item \define{weak equivalences} $\cW$,
    \item \define{cofibrations} $\cC\mathsf{of}$,
    \item \define{fibrations} $\cF\mathsf{ib}$,
  \end{itemize}
  such that the following axioms hold:
  \begin{enumerate}
    \item[(MC1)] (\emph{2-out-of-3}) If $f$ and $g$ are composable morphisms
      in $\cM$ and two of the three maps $f$, $g$ and $g \circ f$ are weak
      equivalences, then so is the third.
    \item[(MC2)] (\emph{Retracts}) The classes $\cW$, $\cC\mathsf{of}$ and
      $\cF\mathsf{ib}$ are closed under retracts.
    \item[(MC3)] (\emph{Lifting}) Cofibrations have the left lifting property
      with respect to acyclic fibrations, and fibrations have the right
      lifting property with respect to acyclic cofibrations, where
      \begin{equation}
        \cC\mathsf{of} \cap \cW = \{\text{acyclic cofibrations}\}, \qquad
        \cF\mathsf{ib} \cap \cW = \{\text{acyclic fibrations}\}.
      \end{equation}
      Concretely, $i\colon A \to B$ is a cofibration if and only if for every
      trivial fibration $p\colon X \to Y$, any solid commutative square
      \begin{equation}
        \begin{tikzcd}[ampersand replacement=\&]
            A \&\& X \\
            \\
            B \&\& Y
            \arrow[from=1-1, to=1-3]
            \arrow["i"', from=1-1, to=3-1]
            \arrow["p", from=1-3, to=3-3]
            \arrow[from=3-1, to=3-3]
            \arrow[dashed, from=3-1, to=1-3]
        \end{tikzcd}
      \end{equation}
      admits a dashed diagonal lift.  Dually for trivial cofibrations and
      fibrations.
    \item[(MC4)] (\emph{Factorization}) Every morphism $f \colon X \to Y$ in
      $\cM$ admits two functorial factorizations
      \begin{equation}
        X \xrightarrow{i} Z \xrightarrow{p} Y,
      \end{equation}
      with $i$ a cofibration and $p$ an acyclic fibration, and
      \begin{equation}
        X \xrightarrow{j} Z' \xrightarrow{q} Y,
      \end{equation}
      with $j$ an acyclic cofibration and $q$ a fibration.
  \end{enumerate}
  Since $\cM$ is complete and cocomplete it has an initial and a terminal
  object.  An object $X \in \cM$ is called \define{cofibrant} if the unique map
  $\varnothing \to X$ is a cofibration, and \define{fibrant} if the unique map
  $X \to *$ is a fibration.
\end{definition}

\begin{lemma}\label{ch2:: lem: fibrant-cofibrant}
    Let $\cM$ be a model category. For every object $X$ there is
    \begin{itemize}
      \item a cofibrant object $\bQ X$ together with a weak equivalence
        $\bQ X \xra{\sim} X$;
      \item a fibrant object $\bR X$ together with a weak equivalence
        $X \xra{\sim} \bR X$.
    \end{itemize}
    We call these \define{cofibrant} and \define{fibrant} replacements of $X$.
\end{lemma}

\begin{proof}
    Factor the unique map $\varnothing \to X$ as
    \begin{equation}
        \varnothing \xra{i} \bQ X \xra{p} X
    \end{equation}
    as in~\textup{(MC4)}, with $i$ a cofibration and $p$ an acyclic fibration.
    Then $\bQ X$ is cofibrant and $p$ is a weak equivalence.  The construction
    of $\bR X$ is dual, factoring $X \to *$ as an acyclic cofibration followed
    by a fibration.
\end{proof}

\begin{definition}[Cofibrantly generated model category]
  A model category $\cM$ is \define{cofibrantly generated} if there exist sets
  of maps
  \begin{equation}
    I, J \subseteq \Mor(\cM)
  \end{equation}
  such that:
  \begin{enumerate}
    \item A map is an acyclic fibration if and only if it has the right
          lifting property with respect to every map in \(I\).
    \item A map is a fibration if and only if it has the right lifting
          property with respect to every map in \(J\).
    \item The domains of maps in \(I\) are small relative to \(I\)-cell
          complexes, and the domains of maps in \(J\) are small relative to
          \(J\)-cell complexes (see \cite[Def.~10.1.2]{Hir14} for the precise
          smallness condition).
  \end{enumerate}
  We call \(I\) a \define{set of generating cofibrations} and \(J\) a
  \define{set of generating acyclic cofibrations}.
\end{definition}

\begin{definition}[Combinatorial model category]
  A model category $\cM$ is \define{combinatorial} if it is cofibrantly
  generated and its underlying category is locally presentable (equivalently,
  is cocomplete and generated under filtered colimits by a set of objects).
\end{definition}

\begin{definition}[Left and right properness]\label{ch2:: def: proper}
  A model category $\cM$ is
  \begin{itemize}
    \item \define{left proper} if every pushout of a weak equivalence along a
          cofibration is again a weak equivalence;
    \item \define{right proper} if every pullback of a weak equivalence along a
          fibration is again a weak equivalence;
    \item \define{proper} if it is both left and right proper.
  \end{itemize}
\end{definition}

\begin{remark}
  In practice we will almost always work with \emph{cofibrantly generated} and
  \emph{combinatorial} model categories, because these admit good existence
  results for localizations and diagram model structures; see
  \cite[Ch.~11, Ch.~13]{Hir14}.  Properness will be important when we want
  homotopy (co)limits to behave well with respect to weak equivalences.
\end{remark}

\begin{theorem}[Quillen model structure on topological spaces]
  \label{ch2:: thm: Quillen model structure}
  There is a proper, cofibrantly generated model structure on the category
  $\Top$ of topological spaces such that:
  \begin{enumerate}
    \item weak equivalences are the \emph{weak homotopy equivalences}, i.e.\
      maps $f\colon X \to Y$ inducing isomorphisms
      \begin{equation}
        \pi_0(X) \xrightarrow{\;\cong\;} \pi_0(Y),
        \qquad
        \pi_n(X,x_0) \xrightarrow{\;\cong\;} \pi_n(Y,f(x_0))
      \end{equation}
      for all $n\geq 1$ and all basepoints $x_0\in X$;
    \item fibrations are the Serre fibrations;
    \item cofibrations are the retracts of relative cell complexes.
  \end{enumerate}
  A set of generating cofibrations is given by
  \begin{equation}
    I = \{S^{n-1} \hookra D^{n}\}_{n\geq 0},
  \end{equation}
  where $S^{n-1}$ is the $(n-1)$-sphere and $D^n$ the $n$-disk, and a set of
  generating acyclic cofibrations is
  \begin{equation}
    J = \{D^{n}\times \{0\} \hookra D^{n} \times [0,1]\}_{n \geq 0}.
  \end{equation}
  We denote this model structure by $\Top_{\mathrm Q}$ and refer to it as the
  \textbf{Quillen model structure}.
\end{theorem}

\begin{theorem}[Kan--Quillen model structure on simplicial sets]
  \label{ch2:: thm: Kan-quillen model}
  There is a proper, cofibrantly generated model structure on the category
  $\sset$ of simplicial sets such that:
  \begin{enumerate}
    \item weak equivalences are those maps $f\colon K \to L$ whose geometric
      realization $|f|\colon |K| \to |L|$ is a weak homotopy equivalence of
      topological spaces;
    \item fibrations are the Kan fibrations;
    \item cofibrations are the monomorphisms.
  \end{enumerate}
  In particular, every simplicial set is cofibrant.  A set of generating
  cofibrations is
  \begin{equation}
    I = \{\partial \Delta^n \hookra \Delta^{n}\}_{n\geq 0},
  \end{equation}
  and a set of generating acyclic cofibrations is
  \begin{equation}
    J = \{\Lambda^{n}_{k} \hookra \Delta^{n}\,:\, 0 \leq k \leq n\}_{n \geq 0}.
  \end{equation}
  We denote this model structure by $\sset_{\mathrm{QK}}$ and call it the
  \textbf{Kan--Quillen model structure}.
\end{theorem}

\begin{remark}
  Theorems~\ref{ch2:: thm: Quillen model structure} and
  \ref{ch2:: thm: Kan-quillen model} provide our basic examples of model
  categories.  We will see in \Cref{ch2:: thm: homotopy hyp} that
  $\Top_{\mathrm Q}$ and $\sset_{\mathrm{QK}}$ model the same homotopy theory.
\end{remark}

The homotopy theory encoded by a model category $\cM$ is captured by its
\emph{homotopy category}.

\begin{definition}
  Let $\cM$ be a model category.  Its \emph{homotopy category}
  $\ho(\cM)$ is obtained from $\cM$ by formally inverting the weak
  equivalences:
  \begin{equation}
    \ho(\cM) \,\overset{\sf def}{=} \, \cM[\cW^{-1}].
  \end{equation}
\end{definition}

This definition is entirely independent of the fact that $\cM$ is a model
category.  In particular, it makes no reference to the axioms.  In general,
localization (the process of formally inverting a class of maps) is
ill-behaved.  One way to think about model categories is precisely as a
framework where localization behaves well.

We can therefore understand $\ho(\cM)$ through a more concrete description
using homotopies.  Let $X$ be cofibrant and $Y$ fibrant.  A
\define{cylinder object} for $X$ is a factorization of the codiagonal
\begin{equation}
  X \sqcup X \longrightarrow X
\end{equation}
as an acyclic cofibration $X \sqcup X \to \mathsf{Cyl}(X)$ followed by a
fibration $\mathsf{Cyl}(X) \to X$.  A \emph{left homotopy} between two
morphisms $f,g \colon X \to Y$ is then a map
\begin{equation}
  H \colon \mathsf{Cyl}(X) \longrightarrow Y
\end{equation}
restricting to $f$ and $g$ on the two copies of $X$.  If $X$ is cofibrant and
$Y$ is fibrant, left homotopy induces an equivalence relation $\sim_l$ on
$\Hom_{\cM}(X,Y)$; see \cite[Proposition~4.7]{Hir14}.

\begin{theorem}\label{ch2:: thm: homotopy cat def}
  Let $\cM$ be a model category and let $X$ be cofibrant and $Y$ fibrant.
  Then the canonical map
  \begin{equation}
    \Hom_{\cM}(X,Y)/\sim_{l}
    \longrightarrow
    \ho(\cM)(X,Y)
  \end{equation}
  is a bijection.  In particular, to compute morphisms in $\ho(\cM)$ it
  suffices to replace objects by weakly equivalent cofibrant–fibrant objects.
\end{theorem}

\begin{cor}[Fundamental theorem of model categories]
  \label{ch2:: cor: fundamental theorem model categories}
  Let $\cM_{\mathrm{cf}} \subseteq \cM$ denote the full subcategory of
  cofibrant–fibrant objects, and let
  \begin{equation}
    \gamma \colon \cM \longra \ho(\cM)
  \end{equation}
  be the localization functor.  Then $\gamma$ induces an equivalence of
  categories
  \begin{equation}
    \cM_{\mathrm{cf}}/\sim_l \xra{\simeq} \ho(\cM),
  \end{equation}
  where $\cM_{\mathrm{cf}}/\sim_l$ has the same objects as
  $\cM_{\mathrm{cf}}$ and morphisms given by left-homotopy classes of maps.
\end{cor}

\begin{remark}
  \Cref{ch2:: cor: fundamental theorem model categories} gives the concrete
  description of the localization promised above.  In particular, for every
  pair of objects $X,Y\in \cM$ there is a natural bijection
  \begin{equation}
    \Hom_{\ho(\cM)}(X,Y)
    \;\cong\;
    \Hom_{\cM}(\bQ X,\bR Y)/\sim_l.
  \end{equation}
  Moreover, if $X$ and $Y$ are cofibrant–fibrant, then a morphism
  $f\colon X \to Y$ is a weak equivalence if and only if there exists a map
  $g\colon Y \to X$ such that
  \begin{equation}
    g f \sim_l \id_X,
    \qquad
    f g \sim_l \id_Y.
  \end{equation}
  In other words, weak equivalences between cofibrant–fibrant objects are
  precisely homotopy equivalences.
\end{remark}

\begin{example}
  The model category $\Top_{\mathrm Q}$ of
  \Cref{ch2:: thm: Quillen model structure} has a homotopy category equivalent
  to the classical homotopy category of CW-complexes and homotopy classes of
  maps.
\end{example}

\begin{example}
  The model category $\sset_{\mathrm{QK}}$ of
  \Cref{ch2:: thm: Kan-quillen model} has a homotopy category equivalent to
  $\ho(\Top_{\mathrm Q})$; this follows from the Quillen equivalence of
  \Cref{ch2:: thm: homotopy hyp} below.
\end{example}


\subsubsection{Quillen adjunctions and Quillen equivalences}

Model categories are meant to be compared by functors that preserve the
homotopy theory they encode.

\begin{definition}
  Let $\cM$ and $\cN$ be model categories.  An adjunction
  \[
    F \colon \cM \rightleftarrows \cN \colon G
  \]
  is called a \emph{Quillen adjunction} (or \emph{Quillen pair}) if the left
  adjoint $F$ preserves cofibrations and acyclic cofibrations.  Equivalently,
  the right adjoint $G$ preserves fibrations and acyclic fibrations.
\end{definition}

Given a Quillen adjunction, one can form its \emph{derived} adjunction on the
homotopy categories.

\begin{definition}
  Let $F \dashv G$ be a Quillen adjunction between model categories.  Choose a
  cofibrant replacement functor $Q \colon \cM \to \cM$ and a fibrant
  replacement functor $R \colon \cN \to \cN$.  The \emph{total left derived
  functor} and \emph{total right derived functor} are the functors
  \begin{equation}
    \LL F \colon \ho(\cM) \longrightarrow \ho(\cN),
    \qquad
    \RR G \colon \ho(\cN) \longrightarrow \ho(\cM),
  \end{equation}
  defined on objects by $\LL F(X) = F(\bQ X)$ and $\RR G(Y) = G(\bR Y)$, and on
  morphisms by applying $F$ or $G$ to representatives between cofibrant or
  fibrant replacements.  The pair $(\LL F,\RR G)$ is again an adjoint pair.
\end{definition}

\begin{definition}
  A Quillen adjunction $F \dashv G$ is called a \define{Quillen equivalence} if
  any (hence all) of the following equivalent conditions hold:
  \begin{enumerate}
    \item For every cofibrant object $X \in \cM$ and every fibrant object
      $Y \in \cN$, a morphism
      \begin{equation}
        F(X) \longrightarrow Y
      \end{equation}
      is a weak equivalence in $\cN$ if and only if its adjoint
      $X \to G(Y)$ is a weak equivalence in $\cM$.
    \item The derived adjunction
      \begin{equation}
        \LL F \colon \ho(\cM) \rightleftarrows \ho(\cN) \colon \RR G
      \end{equation}
      is an equivalence of categories.
    \item The left adjoint $F$ reflects weak equivalences between cofibrant
      objects and, for every fibrant $Y \in \cN$, the composite
      \begin{equation}
        \bQ F(G(Y)) \longrightarrow F(G(Y)) \xrightarrow{\;\varepsilon_Y\;} Y
      \end{equation}
      is a weak equivalence in $\cN$, where $\bQ$ is a cofibrant replacement
      and $\varepsilon$ is the counit.
  \end{enumerate}
\end{definition}

\begin{example}
  The adjoint pair
  \[
    |{-}| \colon \sset_{\mathrm{QK}}
    \rightleftarrows
    \Top_{\mathrm Q} \colon \Sing
\]
  of geometric realization and the singular complex is a Quillen adjunction,
  and indeed a Quillen equivalence (Theorem~\textup{2.1.30}).  Its derived
  adjunction identifies $\ho(\sset_{\mathrm{QK}})$ with $\ho(\Top_{\mathrm Q})$.
\end{example}

\begin{theorem}[Projective and injective model structures on functor categories, cf.\ {\cite[Ch.~11]{Hir14}}]
  \label{ch2:: thm: functor model structures}
  Let $\cF$ be a small category and let $\cM$ be a combinatorial model
  category. Then $\Fun(\cF,\cM)$ admits both the projective and the injective
  model structures, each with weak equivalences defined levelwise. Moreover, if
  $\cM$ is left proper or simplicial, then the corresponding functor model
  structures inherit the same property.
\end{theorem}

\begin{example}
  Let $u \colon \cF \to \cG$ be a functor between small categories, and let
  $\cM$ be a combinatorial model category.  Then precomposition
  \[
    u^{\ast} \colon \Fun(\cF,\cM) \longrightarrow \Fun(\cG,\cM)
  \]
  admits both a left and a right adjoint, given by left and right Kan
  extension along $u$, which we denote by $\lkan_u$ and $\rkan_u$:
    \begin{equation}
        \begin{tikzcd}[ampersand replacement=\&]
        {\fun(\cF,\cM)} \\
        \\
        {\fun(\cG,\cM)}
        \arrow[""{name=0, anchor=center, inner sep=0}, "{u^*}"{description}, from=1-1, to=3-1]
        \arrow[""{name=1, anchor=center, inner sep=0}, "{\lkan_{u}}", bend left=50, from=3-1, to=1-1]
        \arrow[""{name=2, anchor=center, inner sep=0}, "{\rkan_{u}}"', bend right=50, from=3-1, to=1-1]
        \arrow["\dashv"{anchor=center}, draw=none, from=0, to=2]
        \arrow["\dashv"{anchor=center}, draw=none, from=1, to=0]
        \end{tikzcd}
    \end{equation}
  For the \emph{projective} model structures on functor categories,
  $(\lkan_u,u^{\ast})$ is a Quillen pair, and for the \emph{injective}
  model structures $(u^{\ast},\rkan_u)$ is a Quillen pair; to see why, note that if $F\maps \cF'\longra \cM$ is a level-wise weak equivalence and $\alpha \maps c \longra d$ is a morphism in $\cF$, then $u^{*}F(\alpha) = F\big(u(\alpha)\big)$ will be a level-wise weak equivalence i.e., $u^*$ preserves weak equivalences. Moreover, if $F$ is a projective fibration, then $u^* F$ is a projective fibration by the same argument. Dually, $u^*$ preserves injective cofibrations. We will use this repeatedly when passing
  between model structures on diagram categories.
\end{example}

\begin{theorem}[Homotopy hypothesis for spaces vs.\ simplicial sets]
  \label{ch2:: thm: homotopy hyp}
  The adjoint pair
  \begin{equation}
    |{-}| \colon \sset_{\mathrm{QK}}
    \rightleftarrows
    \Top_{\mathrm Q} \colon \Sing
  \end{equation}
  of geometric realization and singular complex is a Quillen adjunction, and
  indeed a Quillen equivalence of model categories.
\end{theorem}

\begin{remark}
  As a consequence, the homotopy category of $\sset_{\mathrm{QK}}$ is
  equivalent to that of $\Top_{\mathrm Q}$.  In particular, the familiar
  homotopy theory of topological spaces can be studied equally well using
  simplicial sets with the notion of weak equivalence from
  \Cref{ch2:: thm: Kan-quillen model}.  From now on we will often refer to
  simplicial sets simply as \emph{spaces}, and we will only distinguish
  between topological spaces and simplicial sets when the context requires it.
\end{remark}



\subsubsection{Homotopy limits, colimits and left Bousfield localization}
\label{ch2:: subsubsec: homotopy-limits-bousfield}

We will need homotopy limits and colimits of diagrams valued in a simplicial
model category.

\begin{definition}[Homotopy (co)limits]\label{ch2:: def: homotopy-(co)limits}
  Let $\cM$ be a model category and let $\cF$ be a small category.
  We endow the diagram category $\Fun(\cF,\cM)$ with one of the standard
  projective or injective model structures from
  \Cref{ch2:: thm: functor model structures}, so that weak equivalences are
  levelwise in~$\cM$.

  The constant diagram functor
  \begin{equation}
    \const \colon \cM \longra \Fun(\cF,\cM)
  \end{equation}
  admits both a left adjoint and a right adjoint, namely the colimit and
  limit functors
  \begin{equation}
    \Colimit[\cF] \colon \Fun(\cF,\cM) \longra \cM,
    \qquad
    \Limit[\cF] \colon \Fun(\cF,\cM) \longra \cM,
  \end{equation}
  and for the chosen diagram model structure:
  \begin{enumerate}
    \item $\Colimit[\cF]$ is a left Quillen functor;
    \item $\Limit[\cF]$ is a right Quillen functor.
  \end{enumerate}

  Let $X \in \Fun(\cF,\cM)$.
  \begin{itemize}
    \item The \define{homotopy limit} of $X$ is the value of the total
      right derived functor of $\lim_{\cF}$:
      \begin{equation}
        \hoLimit[\cF] X
        \;\overset{\mathsf{def}}{=}\;
        \RR\!\lim_{\cF} X.
      \end{equation}
      Equivalently, if $\bR X$ is a fibrant replacement of $X$
      in $\Fun(\cF,\cM)$ for the chosen diagram model structure, then
      \begin{equation}
        \hoLimit[\cF] X
        \;\simeq\;
        \lim_{\cF} (\bR X).
      \end{equation}
    \item The \define{homotopy colimit} of $X$ is the value of the total
      left derived functor of $\colim_{\cF}$:
      \begin{equation}
        \hoColimit[\cF] X
        \;\overset{\mathsf{def}}{=}\;
        \LL\!\Colimit[\cF] X.
      \end{equation}
      Equivalently, if $\bQ X$ is a cofibrant replacement of $X$
      in $\Fun(\cF,\cM)$, then
      \begin{equation}
        \hoColimit[\cF] X
        \;\simeq\;
        \Colimit[\cF] (\bQ X).
      \end{equation}
  \end{itemize}
\end{definition}



When $\cM$ is simplicial, one can also describe homotopy (co)limits via
Bousfield--Kan formulas \cite[Ch.~18]{Hir14}. In many examples of interest
(e.g.\ $\cM=\sset$ and $\cF$ direct), the Reedy model structure on $\Fun(\cF,\cM)$ exists and coincides with the injective model structure,
but we will not use this explicitly.



\begin{definition}[Homotopy Kan extensions]\label{ch2:: def: homotopy-kan}
  Let $\cM$ be a model category, let $\cF,\cG$ be small categories, and let
  \[
    u \colon \cF \longra \cG
  \]
  be a functor.

  \begin{enumerate}
    \item \textbf{Homotopy left Kan extension.}
    Endow $\Fun(\cF,\cM)$ and $\Fun(\cG,\cM)$ with the \emph{projective}
    diagram model structures. Precomposition with $u$,
    \begin{equation}
      u^{\ast} \colon \Fun(\cG,\cM) \longra \Fun(\cF,\cM),
    \end{equation}
    has a left adjoint
    \begin{equation}
      \lkan_u \colon \Fun(\cF,\cM) \longra \Fun(\cG,\cM),
    \end{equation}
    the (ordinary) left Kan extension along $u$.
    In the projective model structure the functor $\lkan_u$ is left
    Quillen, so we can form its total left derived functor
    \begin{equation}
      \Llkan \colon
      \ho\big(\Fun(\cF,\cM)\big) \longra \ho\big(\Fun(\cG,\cM)\big).
    \end{equation}
    For $X\in \Fun(\cF,\cM)$ we call
    \begin{equation}
      \Llkan_u X
    \end{equation}
    the \define{homotopy left Kan extension} of $X$ along $u$.

    Equivalently, if $\bQ X$ is a cofibrant replacement of $X$ in
    $\Fun(\cF,\cM)$ (projective), then
    \begin{equation}
      \Llkan_u X \;\simeq\; \lkan_u(\bQ X)
    \end{equation}
    in $\ho\big(\Fun(\cG,\cM)\big)$.

    \item \textbf{Homotopy right Kan extension.}
    Endow $\Fun(\cF,\cM)$ and $\Fun(\cG,\cM)$ with the \emph{injective}
    diagram model structures. Then $u^{\ast}$ has a right adjoint
    \begin{equation}
      \rkan_{u} \colon \Fun(\cF,\cM) \longra \Fun(\cG,\cM),
    \end{equation}
    the (ordinary) right Kan extension along $u$.
    In the injective model structure the functor $\rkan_{u}$ is right
    Quillen, so we can form its total right derived functor
    \begin{equation}
      \Rrkan_{u} \colon
      \ho\big(\Fun(\cF,\cM)\big) \longra \ho\big(\Fun(\cG,\cM)\big).
    \end{equation}
    For $X\in \Fun(\cF,\cM)$ we call
    \begin{equation}
      \Rrkan_{u} X
    \end{equation}
    the \emph{homotopy right Kan extension} of $X$ along $u$.

    Equivalently, if $\bR X$ is a fibrant replacement of $X$ in
    $\Fun(\cF,\cM)$ (injective), then
    \begin{equation}
      \Rrkan_{u} X \;\simeq\; \rkan_{u}(\bR X)
    \end{equation}
    in $\ho\big(\Fun(\cG,\cM)\big)$.
  \end{enumerate}
\end{definition}

\begin{remark}\label{ch2:: rem: pointwise-kan}
 For $X\colon \cF\to\cM$ and
  $g\in \cG$ one has the usual formulas
  \begin{equation}
    (\lkan X)(g) \;\cong\;
    \mathsf{colim}\big(\cF/g \xra{\,\, \sf forget \,\,} \cG \xra{u} \cM \big),
    \qquad
    (\rkan_{u}X)(g) \;\cong\;
    \mathsf{lim}\big(g/\cF \xra{\,\, \sf forget \,\,} \cG \xra{u} \cM \big),
  \end{equation}
  where $(u/g)$ and $(g/u)$ are the usual comma categories.
  Passing to derived functors and using
  Definition~\textup{\ref{ch2:: def: homotopy-(co)limits}} we obtain
  natural weak equivalences
  \begin{equation}
    (\Llkan_u X)(g)
    \;\simeq\;
    \mathsf{hocolim}\big(\cF/g \xra{\,\, \sf forget \,\,} \cG \xra{u} \cM \big),
    \qquad
    (\Rrkan_{u} X)(g)
    \;\simeq\;
    \mathsf{holim}\big(g/\cF \xra{\,\, \sf forget \,\,} \cG \xra{u} \cM \big).
  \end{equation}
  The existence of this pointwise forumalae for homotopy Kan extensions is asserted and proved in \cite[Proposition~1.14]{Cis09}
\end{remark}

\begin{remark}\label{ch2:: rem: holim-hocolim-as-kan}
  Let $\cF$ be a small category and let
  \begin{equation}
    p \colon \cF \longra *
  \end{equation}
  be the unique functor to the terminal category.
  Then $\lkan_p$ is the colimit functor and $\rkan_{p}$ is the limit functor.
  In particular, for $X\in\Fun(\cF,\cM)$ there are canonical identifications
  \begin{equation}
    \hoColimit[\cF] X
    \;\simeq\;
    \big(\LL\mathsf{Kan}_p X\big)(*),
    \qquad
    \hoLimit[\cF] X
    \;\simeq\;
    \big(\Rrkan_{p} X\big)(*).
  \end{equation}
  Thus, homotopy (co)limits are precisely the (left) right homotopy Kan extensions along the terminal functor.
\end{remark}


\begin{proposition}\label{ch2:: prop: homotopy-limit-homotopically-constant}
  Let $\cM$ be a simplicial model category, let $\cF$ be a small category
  whose nerve is contractible, and let
  $X \colon \cF \to \cM$ be a homotopically constant diagram (i.e.~every map
  $X(f)$ is a weak equivalence) which is fibrant as an $\cF$-diagram.  Then
  for every $i \in \cF$ the canonical maps
  \[
    \holim_{\cF} X \longrightarrow X(i),
    \qquad
    X(i) \longrightarrow \hocolim_{\cF} X
  \]
  are weak equivalences in $\cM$.
\end{proposition}

\begin{example}
  Let $X \colon \Delta^{\op} \to \sset$ be an injectively fibrant simplicial object in $\sset$, which is levelwise weakly equivalent to a constant diagram on a
  Kan complex $K$.  Then
  $\holim_{\Delta^{\op}} X \simeq K \simeq \hocolim_{\Delta^{\op}} X$.  This
  is a typical situation when simplicial resolutions are used to compute
  derived mapping spaces.
\end{example}

\begin{definition}[Homotopy initial and homotopy final functors]
  Let $u \colon \cF \to \cG$ be a functor between small categories and let
  $\cM$ be a model category.

  \begin{enumerate}
    \item We say that $u$ is \define{homotopy final} if, for every diagram
      $X \colon \cG \to \cM$, the canonical map
      \[
        \hoColimit[\cF] (X \circ u)
        \longrightarrow
        \hoColimit[\cG] X
      \]
      is a weak equivalence in $\cM$, whenever these homotopy colimits are
      defined.

    \item We say that $u$ is \define{homotopy initial} if, for every diagram
      $X \colon \cG \to \cM$, the canonical map
      \[
        \hoLimit[\cG] X
        \longrightarrow
        \hoLimit[\cF] (X \circ u)
      \]
      is a weak equivalence in $\cM$, whenever these homotopy limits are
      defined.
  \end{enumerate}
\end{definition}

The preceding definition is conceptually clean but not very practical.  The
usual criterion is phrased in terms of comma categories and their nerves.

\begin{proposition}\label{ch2::prop:homotopy-final-initial}
  Let $u \colon \cF \to \cG$ be a functor between small categories.
  \begin{enumerate}
    \item The functor $u$ is homotopy final if and only if, for every object
      $g \in \cG$, the nerve of the over-category $(g / \cF)$ is
      weakly contractible:
      \begin{equation}
        \sN\bigl(g / \cF\bigr) \simeq *.
      \end{equation}
    \item The functor $u$ is homotopy initial if and only if, for every object
      $g \in \cG$, the nerve of the under-category $(\cF / g)$ is
      weakly contractible:
      \begin{equation}
        \sN\bigl(\cF / g\bigr) \simeq *.
      \end{equation}
  \end{enumerate}
  In particular, these conditions do not depend on the choice of model
  category $\cM$.
\end{proposition}

\begin{example}
  If $\cF$ has a terminal object $t$ and $u \colon \cF \to \cG$ sends $t$ to a
  terminal object of $\cG$, then $u$ is homotopy final.  Indeed, for any
  $g \in \cG$ the over-category $(g / \cF)$ has a terminal object (if
  it is nonempty), so its nerve is contractible.  Dually, if $\cF$ has an
  initial object that $u$ preserves, then $u$ is homotopy initial.
\end{example}



Finally, we recall left Bousfield localization, which will be used to
construct the model structures for (complete) Segal spaces and (complete) interval Segal spaces.

\begin{definition}
  Let $\cM$ be a model category and let $K$ be a set of maps in $\cM$.
  \begin{itemize}
    \item A fibrant object $Z \in \cM$ is called \define{$K$-local} if for every
      $f \colon A \to B$ in $K$, the induced map on derived mapping spaces
      \begin{equation}
        \Map^{h}_{\cM}(B,Z) \longrightarrow \Map^{h}_{\cM}(A,Z)
      \end{equation}
      is a weak equivalence of simplicial sets.
    \item A map $g \colon X \to Y$ in $\cM$ is a \define{$K$-local equivalence}
      if for every $K$-local $Z$ the induced map
      \begin{equation}
        \Map^{h}_{\cM}(Y,Z) \longrightarrow \Map^{h}_{\cM}(X,Z)
      \end{equation}
      is a weak equivalence.
  \end{itemize}
  A \define{left Bousfield localization} of $\cM$ with respect to $K$ is a model
  structure $L_K\cM$ on the same underlying category, with the same
  cofibrations, whose weak equivalences are the $K$-local equivalences and
  whose fibrant objects are the $K$-local fibrant objects.
\end{definition}

\begin{theorem}[Left Bousfield localization, {\cite[Thm.~3.3.1, Thm.~4.1.1]{Hir14}}]
  \label{ch2:: thm: left Bousfield localization}
  Let $\cM$ be a left proper, cellular or combinatorial model category, and
  let $K$ be a set of maps in $\cM$.  Then the left Bousfield localization
  $L_K\cM$ exists.  Moreover:
  \begin{itemize}
    \item $L_K\cM$ is left proper and cellular or combinatorial (accordingly);
    \item if $\cM$ is a simplicial model category, then so is $L_K\cM$.
  \end{itemize}
\end{theorem}

\begin{proposition}
  Let $\cM$ be a model category and let $K$ be a set of maps in $\cM$.  Then
  a map between $K$-local fibrant objects is a $K$-local weak equivalence if
  and only if it is a weak equivalence in $\cM$.
\end{proposition}




\begin{theorem}\label{ap:: thm: transfer of localizations}
    Suppose that $\cM$ and $\cN$ are model categories and $F\maps \cM \leftrightarrows \cN\maps G$ is a Quillen pair. Suppose that $T$ is a class of maps in $\cM$ such that the localized model structure $L_{T}\cM$ exists, and furthermore that the localized model structure $L_{\mathbb{L}F(T)}\cN$ also exists. Then:
    \begin{enumerate}
        \item The adjoint pair $F\dashv G$ still defines a Quillen pair
        \begin{equation}
            F\maps L_{T}\cM \leftrightarrows L_{\mathbb{L}F(T)}\cN \maps G.
        \end{equation}
        \item If the pair $F\dashv G$ defines a Quillen equivalence between $\cM$ and $\cN$, then it also defines a Quillen equivalence on the localized model structures.
    \end{enumerate}
\end{theorem}


%%% suggestions for the model category stuff end %%%%



\section{A primer on complete Segal spaces}
\label{ch2:: sec: complete_segal_spaces}

The main goal of the present chapter is to propose a variant model for $(\oo,1)$-categories, yet, so far, we have not made any progress into introducing what these are. This section not only addresses that missing exposition, but it introduces Segal spaces after which we will model our own variant of $(\oo,1)$-categories. The goal is not to give a precise account of the subject, which contains many intricate bells and whistles — details about them can be found in Rezk's original paper \cite{Rezk01} or Julie Bergner's comprehensive book on the homotopy theory of $(\oo,1)$-categories \cite{Ber18} — these are our main references for this section.

An $(\oo,1)$-category is a generalization of the usual notion of category, where the set of morphisms is endowed with more structure, namely that of a homotopy type modeled by a topological space or a simplicial set. Of course, this conceptual change comes with modifications to the way one compares and classifies $(\oo,1)$-categories, especially since there is an obvious expectation that the collection of $(\oo,1)$-categories forms itself an $(\oo,1)$-category. With these objects in mind, we finally turn to the work of Rezk.

\subsection{Simplicial spaces}

Segal spaces, despite their suggestive name, are not actually spaces, they are simplicial spaces. Before we delve into why these are the right objects to use, let us quickly go over what is a category again. category $\sC$ is a collection of objects $\sC_{0}$, together with a collection of morphisms $\sC_{1}$, an additional pair of maps 
\begin{equation}
    \sC_0 \xleftarrow{\,\,\,\mathsf{source}\,\,\,} \sC_{1} \xra{\,\,\,\mathsf{target}\,\,\,} \sC_{0},
\end{equation} 
together with a composition rule 
\begin{equation}
    \sC_{1} \times_{\sC_{0}} \sC_{1} \xra{\,\,\ \circ \,\,} \sC_{1} 
\end{equation}
satisfying an additional associativity condition, encoded by the commutativity of the square (\ref{ch2:: eq: assoc})
\begin{equation}\label{ch2:: eq: assoc}
    \begin{tikzcd}
        {\sC_1\times_{\sC_{0}} \sC_{1}\times_{\sC_{0}} \sC_{1}} && {\sC_1\times_{\sC_{0}} \sC_{1}} \\
        \\
        {\sC_1\times_{\sC_{0}} \sC_{1}} && {\sC_1}\, .
        \arrow["{ \id \times \circ}", from=1-1, to=1-3]
        \arrow["\circ", from=1-3, to=3-3]
        \arrow["{\circ \times \id}", from=3-1, to=1-1]
        \arrow["{ \circ}"', from=3-1, to=3-3]
    \end{tikzcd}
\end{equation}
Additionally, there is a map $\sC_{0} \xra{e} \sC_{1}$ sending $x$ to its corresponding identity morphism. 

Importantly, all of this data and structure could just as well be packaged as a simplicial set: 

\begin{equation}
    \sN(\sC) \maps \bDelta^{\op} \longra \mathsf{Set}.
\end{equation} 
Let $\sN(\sC)_0 = \sC_0$, $\sN(\sC)_{1} = \sC_1$, and 
\begin{equation}
\sN(\sC)_{n} = \underbrace{\sC_{1} \times_{\sC_{0}} \sC_{1} \times_{\sC_{0}} \cdots \times_{\sC_{0}} \sC_{1}}_{n \text{ times}}.
\end{equation} 

Moreover, let $\langle 0\rangle^{*},\langle 1\rangle^{*}$ get sent to the source and target maps, respectively. Let the unique map $\langle 00 \rangle^{*}$ get sent to $\sC_{0} \xra{e} \sC_{1}$. Immediately we note that the simplicial relations enforce a condition — the fact that the source and target of the identity morphism are the same. Composition is defined as the map $\langle 02 \rangle^{*}\maps \sN(\sC)_{2} \longra \sN(\sC)_{1}$, which means that associativity, i.e., a condition on the data, also gets enforced or encoded by the simplicial relations. 

This construction highlights the fact that categories naturally give rise to simplicial sets. More is true, since they do so in a functorial way.

\begin{proposition} \label{ch2:: prop: nerve}
    There exists a fully faithful functor, which we will call the \textbf{nerve}
    \begin{equation}
        \sN \maps \Cat \longra \sset.
    \end{equation}
\end{proposition}

An interesting question is whether we can characterize the essential image of the nerve construction. In other words, can we find a condition on a simplicial set $X_\bullet$ which ensures it is isomorphic to the nerve $\sN(\sC)$ of some category $\sC$. The answer is in the positive, and the condition is recorded in the theorem below.
\begin{theorem}\label{ch2:: thm: Kan--quillen}
    Let $\cC\maps \bDelta^{\op} \longra \set$ be a simplicial set satisfying the following condition: for every $n,m\geq 1$, the diagram
    \begin{equation}
        \begin{tikzcd}
            {\cC_{n+m}} && {\cC_{n}} \\
            \\
            {\cC_{m}} && {\cC_{0}}
            \arrow["{\langle 01\cdots n \rangle^{*} }", from=1-1, to=1-3]
            \arrow["{\langle n\cdots n+m\rangle^{*}}"', from=1-1, to=3-1]
            \arrow["\lrcorner"{anchor=center, pos=0.125}, draw=none, from=1-1, to=3-3]
            \arrow["{\langle n\rangle^{*}}", from=1-3, to=3-3]
            \arrow["{\langle m \rangle^{*}}"', from=3-1, to=3-3]
        \end{tikzcd}
    \end{equation}
    is a pullback square in $\mathsf{Set}$. We call this condition the \textbf{Segal condition}. A simplicial set $\cC_\bullet$ is isomorphic to the nerve of some category $\sC$ if and only if it satisfies the Segal condition.
\end{theorem}

\begin{remark}
 Note that we can consider an equivalent variant of the Segal condition where instead of decomposing $[n+m]$ into linear subposets $[n]$ and $[m]$, we decompose them into more linear subposets of length 1. We can consider the inert maps $[1] \xra{\langle i, i+1\rangle} [n]$ and re-define the Segal condition by demanding that for every $n\geq 2$ the map
 \begin{equation}
    \cC_{n} \longra \cC_{1} \times_{\cC_{0}} \cdots \times_{\cC_{0}} \cC_{1}
 \end{equation}
 is an isomorphism of sets.
\end{remark}

\begin{observation}
The Segal condition says that the $n$-simplices of $\cC$ are, up to isomorphism, strings of $n$-composable 1-simplices, which in turn are manifestly interpreted as the 1-morphisms of some category whose nerve is isomorphic to $\cC$.
\end{observation}

This theorem informs us that simplicial sets house the two ingredients which are important to build $(\oo,1)$-categories: 1) categories and 2) homotopy types. Even more explicitly, we want our classes of morphisms to become homotopy types and we want that associativity and unitality be modified accordingly — associativity need only hold up to a weak homotopy equivalence. Simplicial spaces, or bisimplicial sets, offer an appropriate receptacle for both structures at once as we are interested in going from sets to spaces, it makes sense that we study functors
\begin{equation}
    \bDelta^{\op} \longra \sset,
\end{equation}
which can equivalently be defined and studied as bisimplicial sets. The former perspective is the one we will take when we refer to simplicial spaces — with exceptions, all of which are duly flagged.

\begin{remark}
    The category $\fun(\bDelta^{\op},\sset)$ of simplicial spaces contains two fully faithful copies of $\sset$, and their respective images have very different interpretations. Given a simplicial set $K$, we define the \textbf{constant} simplicial space on $K$ by $([n] \mapsto K\in \sset)$ and every $[n] \ra [m]\in \bDelta^{\op}$ gets sent to the identity. Alternatively, we have the case that $K$ factors through the inclusion $\mathsf{Set} \hookra \sset$ so that $K_n$ is a discrete simplicial set i.e., a set for every $n\geq 0$. What makes this distinction so relevant is the fact that it is only one of the versions of $\Delta^{n}$ in $\fun(\bDelta^{\op},\sset)$ which is the stand-in for the standard $n$-simplex, namely the discrete one. We denote this simplicial space by $\Delta[n]$ in order to contrast with the simplicial set $\Delta^n$. In particular we have that for every simplicial space $X\maps \bDelta^{\op} \ra \sset$, there is a weak equivalence of simplicial sets
    \begin{equation}
        \Map(\Delta[n], X) \simeq X_n 
    \end{equation}
    for every $n\geq 0$.
\end{remark}

We will shift our focus to simplicial spaces now. 

\begin{observation}\label{ch2:: obs: reedy is injective}
In particular, the category $\fun(\bDelta^{\op},\sset)$ can be endowed with
either the injective model structure or the projective model structure. In fact
it can even be endowed with a third one, called the \emph{Reedy model
structure}. All of them are Quillen equivalent, and in this case the injective
and Reedy model structures agree. For simplicity, we will choose the injective
model structure.
\end{observation}

\begin{theorem}
    There exists a model structure on the category of simplicial spaces $\fun(\bDelta^{\op},\sset)$ which we will denote $\fun(\bDelta^{\op},\sset)_{\sf inj}$, with the following properties:
    \begin{enumerate}
        \item The weak equivalences are levelwise weak equivalences in $\sset_{\sf QK}$.
        \item The cofibrations are levelwise cofibrations in $\sset_{\sf QK}$.
        \item The model structure is left proper combinatorial and simplicial.
    \end{enumerate}
\end{theorem}

\begin{observation}
    \label{ch2:: obs: discreteness}
A pair $X$ and $Y$ of discrete simplicial spaces are weakly equivalent as injective simplicial spaces if and only if $X_{k} \simeq Y_{k}$ for every $k$. In other words, if the sets $X_k$ and $Y_k$ are bijective for all $k\geq 0$. 
\end{observation}


\subsection{Imposing the Segal condition}

We want a definition of $(\oo,1)$-category which resembles the nerve of an
ordinary category but relaxes the requirements of strict composition,
associativity, and unitality.

\begin{construction}
Consider the class of inert morphisms
\[
  \big\{[1] \xra{\langle i, i+1 \rangle} [n]\,:\, 0\leq i < n\big\}_{n \geq 1}
\]
in $\bDelta$. Let $\bDelta^{\sf int}\subset\bDelta$ denote the subcategory of
inert morphisms and let $\bDelta^{\sf int}_{\leq 1}\subset\bDelta^{\sf int}$ be
the full subcategory on $[0]$ and $[1]$. For each $[n]\in\bDelta$ we can form
the slice category $\bDelta^{\sf int}_{\leq 1}/[n]$. We define the
\emph{$n$-spine} to be the homotopy colimit
\begin{equation}
  G(n) \,\overset{\sf def}{=}\,
  \hocolim\Big(
    \bDelta^{\sf int}_{\leq 1}/[n]
      \longra \bDelta_{\leq 1}^{\sf int}
      \xra{\,\bDelta[\bullet]\,}
      \Fun(\bDelta^{\op},\sset)
  \Big),
\end{equation}
where $\bDelta[\bullet]$ is the Yoneda embedding.
\end{construction}

\begin{remark}
    The objects of $\bDelta^{\sf int}_{\leq 1}/[n]$ are precisely the inert
    maps $[1] \xra{\langle i,i+1\rangle} [n]$ for $0\leq i<n$.
\end{remark}

\begin{observation}
    The $2$-spine $G(2)$ happens to coincide with the inner horn
    $\Lambda_{1}[2]$, but this is no longer true for $n\geq 3$. For higher $n$,
    the $n$-spine is a string of $1$-simplices and does not contain higher
    non-degenerate simplices, unlike the higher inner horns.
\end{observation}

Let $\cF\colon \bDelta^{\op} \to \sset$ be a simplicial space, which we will
assume to be injectively fibrant. For $n\geq 2$ there is an inclusion of
simplicial spaces $G(n) \hookra \bDelta[n]$, which induces a map
\begin{equation}
    \Map(\bDelta[n], \cF) \longra \Map(G(n), \cF)
      \,\simeq\,
      \hoLimit{(\bDelta^{\sf int}_{\leq 1}/[n])}
      \Map(\bDelta[1], \cF).
\end{equation}
Using the standard identifications $\Map(\bDelta[n],\cF)\cong\cF_n$ and
$\Map(\bDelta[1],\cF)\cong\cF_1$, this can equivalently be recast as a map
\begin{equation}
    \varphi_{n} \colon
    \cF_{n} \longra
    \underbrace{\cF_1 \times_{\cF_{0}}^{h} \cdots
      \times_{\cF_{0}}^{h} \cF_{1}}_{n \text{ times}}.
\end{equation}
These maps $\varphi_n$ for all $n\geq 2$ are called the \textbf{Segal maps}.\footnote{Segal first defined
these maps in the reduced case $\cF_{0} = *$, as we will re-encounter when we
discuss symmetric monoidal $(\oo,1)$-categories.}

Recall that a simplicial set is (isomorphic to) the nerve of an ordinary
category if and only if its Segal maps are isomorphisms of (discrete) simplicial
sets. For a simplicial space we want $\cF_0$ to play the role of the space of
objects, $\cF_1$ the space of morphisms, and $\cF_n$ the space of $n$ composable
morphisms; asking the Segal maps to be weak equivalences is exactly the
homotopy-theoretic analogue of strict composition.

\begin{definition}\label{ch2:: def: segal space}
    Let $\cF \colon \bDelta^{\op} \to \sset$ be an injectively fibrant simplicial
    space. We say that $\cF$ is a \emph{Segal space} if for all $n\geq 2$ the
    Segal maps
    \begin{equation}
        \varphi_{n} \colon
        \cF_{n} \longra
        \underbrace{\cF_1 \times_{\cF_{0}}^{h} \cdots
          \times_{\cF_{0}}^{h} \cF_{1}}_{n \text{ times}}
    \end{equation}
    are weak equivalences of simplicial sets in the Quillen--Kan model
    structure, i.e.\ weak homotopy equivalences. Since $\cF$ is fibrant, the
    ordinary fibre products compute these homotopy fibre products, so we may
    omit the superscript $h$ when convenient.
\end{definition}


\begin{observation}
    Observe that we can also define Segal spaces to be the injectively fibrant simplicial spaces $\cF\maps \bDelta^{\op} \longra \sset$ such that their restriction $\cF_{|\bDelta^{\mathsf{int}, \op}}$ corresponds to the right Kan extension 
    \begin{equation}
    \begin{tikzcd}
        {\bDelta^{\mathsf{int},\op}_{\leq 1}} &&& \sset \\
        \\
        {\bDelta^{\mathsf{int},\op}}
        \arrow["{\cF_{|\bDelta^{\mathsf{int}, \op}_{\leq 1}}}", from=1-1, to=1-4]
        \arrow[hook, from=1-1, to=3-1]
        \arrow["{\cF_{|\bDelta^{\mathsf{int}, \op}}}"', curve={height=12pt}, dashed, from=3-1, to=1-4]
    \end{tikzcd}
    \end{equation}
    We find this a useful reframing of the Segal conditions, since it admits a more immediate generalization when passing from $(\oo,1)$- to $(\oo,n)$-categories via presheaves of simplicial sets.
\end{observation}

\begin{proposition}\label{ch2:: prop: Segal-model-structure}
    There exists a model structure on the category of simplicial spaces
    $\fun(\bDelta^{\op},\sset)$, which we denote by $\mathsf{SSp}$, obtained by
    left Bousfield localizing the injective model structure
    $\fun(\bDelta^{\op},\sset)_{\mathsf{inj}}$ at the set of spine inclusions
    \begin{equation}
        \{G(n)\hookrightarrow \Delta[n]\}_{n\geq 2}.
    \end{equation}
    It has the following properties:
    \begin{enumerate}
        \item The cofibrations are the injective cofibrations, i.e., the monomorphisms.
        \item The fibrant objects are the Segal spaces.
        \item The weak equivalences between Segal spaces are the levelwise weak equivalences.
        \item The model structure is left proper, combinatorial, and simplicial.
    \end{enumerate}
\end{proposition}

\begin{proof}
    Existence follows from \Cref{ch2:: thm: left Bousfield localization},
    applied to the injective model structure on simplicial spaces from the
    theorem above and the set of spine inclusions. Since cofibrations are
    unchanged under left Bousfield localization, (1) holds. By construction,
    the fibrant objects are precisely the injectively fibrant simplicial spaces
    that are local with respect to the spine inclusions, which is exactly the
    Segal condition from \Cref{ch2:: def: segal space}; this gives (2). Finally,
    (3) is the usual description of weak equivalences between fibrant objects in
    a left Bousfield localization, and (4) is inherited from the injective
    model structure.
\end{proof}

\subsection{Homotopy theory of Segal spaces}

Segal spaces, just as any good homotopical gadget, admit a few important
invariants which facilitate their study using algebraic tools. Namely:

\begin{itemize}
    \item Its set of objects.
    \item Its mapping spaces between pairs of objects.
    \item Homotopy equivalence of objects (and morphisms).
    \item Homotopy category.
\end{itemize}


\subsubsection{Set of objects and mapping space}

The best encapsulating property of what it means to be an $(\oo,1)$-category is that the collection of morphisms between any fixed pair of objects organizes as a \emph{space}, with composition being a map of spaces — thereby defining all higher invertible $k$-morphisms as homotopies, in some suitable sense.

\begin{construction}
    \label{ch2:: cons: objects mapping space}
    Let $\cF$ be a Segal space. Let us define its set of \define{objects} to be the set of $0$-simplices of $\cF_0$; and given a pair of objects $x,y\in \cF_{0,0}$, let us define the mapping space $\map_{\cF}(x,y)$ to be the pullback
    \begin{equation}
        \begin{tikzcd}
            {\Map_{\cF}(x,y)} && {\cF_{1}} \\
            \\
            {*} && {\cF_{0}\times \cF_{0}}
            \arrow[from=1-1, to=1-3]
            \arrow[from=1-1, to=3-1]
            \arrow["\lrcorner"{anchor=center, pos=0.125}, draw=none, from=1-1, to=3-3]
            \arrow["\lag 0\rag^* \times \lag 1 \rag^*", from=1-3, to=3-3]
            \arrow["{(x,y)}"', from=3-1, to=3-3]
        \end{tikzcd}
    \end{equation}
    in $\sset$. Recall our assumption that $\cF$ is injectively fibrant — it implies that the right vertical map of simplicial sets is a fibration; as observed in \Cref{ch2:: obs: reedy is injective}, the injective and Reedy model structure agree and since $\partial \Delta[1] \longhookra \Delta[1]$ is cofibration therein, we have that the following induced map is a Kan fibration
    \[\Hom_{\sset}(\Delta[1], \cF) \longra \Hom_{\sset}(\partial \Delta[1], \cF).\]
    In particular, this means that $\Map_{\cF}(x,y)$ is homotopy equivalent to the homotopy fiber of \[\cF_1 \longra \cF_0 \times \cF_0,\] which is desirable, since we want the mapping space $\Map_{\cF}(x,y)$ to be well-defined only up to homotopy, i.e., only depend on the path-component of $x,y\in \cF_0$. We will use the notation
    \begin{equation}
        f \maps x \longra y
    \end{equation}
    for a 0-simplex of $\Map_{\cF}(x,y)$.
\end{construction}

\begin{definition}
    Let $\cF$ be a Segal space. Let $f\maps x \ra y$ and $g\maps x \ra y$ be 0-simplices of the mapping space $\Map_{\cF}(x,y)$. We declare $f \sim g$ as \define{homotopic} if they are in the same component, i.e., if $([f] = [g]) \in \pi_0 \Map_{\cF}(x,y)$.
\end{definition}

\subsubsection{Homotopy equivalences and homotopy category}

Given a Segal space $\cF\maps \bDelta^{\op} \longra \sset$, in \Cref{ch2:: cons: objects mapping space} we associated to it: 1) a set of objects; 2) a mapping space for each pair of objects. Thus, a sensible question is the following: how do we compose morphisms in a Segal space? Let $(x,y,z) \in \cF_{0,0}^{\times 3}$ be a triple of objects of $\cF$. Let $f \maps x \longra y$ and $g \maps y \longra z$ be morphisms in $\Map_{\cF}(x,y)$ and $\Map_{\cF}(y,z)$ respectively. We define the pullback
\begin{equation}
    \begin{tikzcd}[ampersand replacement=\&]
            {\Map_{\cF}(x,y,z)} \&\& {\cF_{2}} \\
            \\
            \ast \&\& {\cF_0\times \cF_0\times \cF_0}
            \arrow[from=1-1, to=1-3]
            \arrow[from=1-1, to=3-1]
            \arrow["\lrcorner"{anchor=center, pos=0.125}, draw=none, from=1-1, to=3-3]
            \arrow["{\lag 0 \rag^* \times \lag 1 \rag^* \times \lag 2 \rag^*}", from=1-3, to=3-3]
            \arrow["{(x,y,z)}"', from=3-1, to=3-3]
        \end{tikzcd}
\end{equation}
to be the space of composable pairs of morphisms $(x \xra{f} y \xra{g} z)$ with fixed $x,y,z$. Using the fact that the Segal map $\cF_2 \longra \cF_1 \times_{\cF_0} \cF_1$ is a weak equivalence, we assert that the dashed map in the following diagram is a weak equivalence:
\begin{equation}
\begin{tikzcd}[ampersand replacement=\&]
	\&\&\& {\Map_{\cF}(x,y,z)} \& {\Map_{\cF}(x,y) \times \Map_{\cF}(y,z)} \\
	\\
	\&\& {\cF_{2}} \& {\cF_1 \times_{\cF_0} \cF_1} \& {*} \\
	\\
	{} \&\&\& {\cF_0\times \cF_0\times \cF_0} \& {.}
	\arrow["\simeq"', dashed, from=1-4, to=1-5]
	\arrow[from=1-4, to=3-3]
	\arrow[""{name=0, anchor=center, inner sep=0}, between={0.5}{1}, from=1-4, to=3-5]
	\arrow["\lrcorner"{anchor=center, pos=0.125, rotate=-45}, draw=none, from=1-4, to=5-4]
	\arrow[from=1-5, to=3-4]
	\arrow[from=1-5, to=3-5]
	\arrow["\lrcorner"{anchor=center, pos=0.125, rotate=-90}, draw=none, from=1-5, to=5-4]
	\arrow["\simeq"', from=3-3, to=3-4]
	\arrow[from=3-3, to=5-4]
	\arrow[from=3-4, to=5-4]
	\arrow[from=3-5, to=5-4]
	\arrow[between={0}{0.75}, no head, from=1-4, to=0]
\end{tikzcd}
\end{equation}
\smallskip

\nid Given a pair $\big(x \xra{f} y,y \xra{g} z\big)\in \Map_{\cF}(x,y)\times \Map_{\cF}(y,z)$, we can form the following diagram:
\smallskip

\begin{equation}\label{comp}
    \begin{tikzcd}[ampersand replacement=\&]
        \&\&\& {\cF_2} \&\& {\cF_1} \\
        \&\& {\Map_{\cF}(x,y,z)} \&\& {\Map_{\cF}(x,z)} \\
        \&\&\& {\cF_1\times_{\cF_0}\cF_1} \&\& {\cF_1} \\
        {*} \&\& {\Map_{\cF}(x,y)\times \Map_{\cF}(y,z)}
        \arrow["{\lag 0,2 \rag^*}", from=1-4, to=1-6]
        \arrow[between={0}{0.4}, no head, from=1-4, to=3-4]
        \arrow["{\rotatebox[origin=c]{90}{$\sim$}}"'{pos=0.75}, between={0.6}{1}, from=1-4, to=3-4]
        \arrow[equals, from=1-6, to=3-6]
        \arrow[from=2-3, to=1-4]
        \arrow[from=2-3, to=2-5]
        \arrow["{\rotatebox[origin=c]{90}{$\sim$}}",from=2-3, to=4-3]
        \arrow[from=2-5, to=1-6]
        \arrow[from=3-4, to=3-6]
        \arrow["h", dashed, from=4-1, to=2-3]
        \arrow["{(f,g)}"', from=4-1, to=4-3]
        \arrow[curve={height=24pt}, dotted, from=4-3, to=2-5]
        \arrow[""{name=0, anchor=center, inner sep=0}, from=4-3, to=3-4]
        \arrow["\lrcorner"{anchor=center, pos=0.125}, draw=none, from=2-3, to=0]
    \end{tikzcd}
\end{equation}
\smallskip

\nid where $h$ is lift of $(f,g)$. We define a composite for $(f,g)$ to be the image of $h$ under $\Map_{\cF}(x,y,z) \longra \Map_{\cF}(x,y)\times \Map_{\cF}(y,z)$. We indicate this ``composition map'' with the curved dotted arrow. We note that this arrow is only well-defined up to homotopy since it depends on the choice $h$ of a lift of $(f,g)$. 


\begin{definition}\label{ch2:: def: homotopy category segal space}
    Let $\cF$ be a Segal space. We define its \textbf{homotopy category} $\ho(\cF)$ to be the category with the following structure:
    \begin{itemize}
        \item \textbf{Objects:} The set of 0-simplices of the simplicial set $\cF_0$.
        \item \textbf{Morphisms:} Given a pair of objects $x,y \in \cF_{0,0}$ we define the set of morphisms of $x$ to $y$ to be 
        \begin{equation}
            \Hom_{\ho(\cF)}(x,y) \, \overset{\sf def}{=} \, \pi_0 \Map_{\cF}(x,y).
        \end{equation}
        \item \textbf{Composition:} Let \[([f],[g])\in \Hom_{\ho(\cF)}(x,y) \times \Hom_{\ho(\cF)}(y,z)\]
        be a pair of composable morphisms. Passing diagram (\ref{comp}) to $\pi_0$, we get a well-defined composition map given by the dashed arrow in the following commutative diagram: 
        \begin{equation}
            \begin{tikzcd}[ampersand replacement=\&]
            {*} \& {\pi_{0}\Map_{\cF}(x,y,z)} \&\& {\pi_0\Map_{\cF}(x,z)} \\
            \\
            {\pi_0\Map_{\cF}(x,y)\times \pi_0\Map_{\cF}(y,z)} \& {\pi_0\Map_{\cF}(x,y)\times \pi_0\Map_{\cF}(y,z)}
            \arrow["{([f],[g])}"{description}, from=1-1, to=3-1]
            \arrow[from=1-2, to=1-4]
            \arrow["\cong", from=1-2, to=3-2]
            \arrow["{\cong} "', from=3-1, to=3-2]
            \arrow["\circ"'{pos=0.6}, curve={height=24pt}, dashed, from=3-2, to=1-4]
        \end{tikzcd}
        \end{equation} 
        The key here is the fact that the weak equivalence became an isomorphism of sets.\footnote{The well-definedness of the composition map comes from the fact that any two choices $\lag 0,2\rag^* h$ and $\lag 0,2\rag^* h'$ of composition for a fixed pair $(f,g)$, are homotopic.} 
    \end{itemize}
\end{definition}

\begin{observation}
    For every object $x\in \ho(\cF)$, there is a 1-morphism $\lag 00\rag ^* x \in \cF_1$ which is the identity. We note that for $f\maps x \ra y$, we can contemplate the following 2-simplex $\lag 0,0,1 \rag^* f \in \cF_2$, whose image under $\lag 0,2\rag^*$ is
    \begin{equation}
        \lag 0,2 \rag ^* \circ \lag 0,0,1\rag ^* f = \left(\lag 0,0,1 \rag \circ \lag 0,2 \rag \right)^*f = \lag 0,1\rag ^{*}f = f,
    \end{equation}
    which confirms that a composite of $f$ with $\id_x$ is again $f$.
\end{observation}

\medskip

Given the fact that composition is not defined on the nose, we cannot (or should not) ask that a given morphism $f\maps x \ra y$ be an isomorphism. For that would force us to make a choice of composite and then ask that choice be the identity —  all the while ignoring all other possible composites. We avoid this by using the homotopy category.

\begin{definition}\label{ch2:: def: hom eq}
Let $\cF$ be a Segal space. We define a morphism $f\maps x \ra y$ to be a \define{homotopy equivalence}\footnote{It is also common to call homotopy equivalences in a Segal space just equivalences.} if $[f]\in \ho(\cF)$ is an isomorphism.
\end{definition}

The next step is to understand to what extent an equivalence of Segal spaces, encoded as a levelwise equivalence of simplicial spaces, captures the correct notion of equivalence between $(\oo,1)$-categories. The short answer is that it does not, but we must first articulate what we expect an equivalence of $(\oo,1)$-categories to be. To do that, we suitably emulate the usual definitions from ordinary category theory. Namely those of fully-faithful and essentially surjective functor.

\begin{definition}
    A morphism $F\maps \cF \longra \cG$ of Segal spaces (i.e., a natural transformation $\bDelta^{\op}\times [1] \longra \sset$) is \define{fully faithful} if the commutative square 
    \[\begin{tikzcd}[ampersand replacement=\&]
        {\cF_1} \&\& {\cG_1} \\
        \\
        {\cF_0\times \cF_0} \&\& {\cG_0 \times \cG_0}
        \arrow[from=1-1, to=1-3]
        \arrow[from=1-1, to=3-1]
        \arrow[from=1-3, to=3-3]
        \arrow[from=3-1, to=3-3]
    \end{tikzcd}\]
    is a homotopy pullback square.
\end{definition}

\begin{observation}
Unpacking the definition above, we see that for every $(x,y) \in \cF_0^{\times 2}$ by the universal property of the pullback, there exists the following dashed map  
\begin{equation}
    \begin{tikzcd}[ampersand replacement=\&]
        \& {\Map_{\cG}(Fx,Fy)} \&\& {\cG_1} \\
        {\Map_{\cF}(x,y)} \&\& {\cF_1} \\
        \& {*} \&\& {\cG_0 \times \cG_0} \\
        {*} \&\& {\cF_0\times \cF_0}
        \arrow[from=1-2, to=1-4]
        \arrow[""{name=0, anchor=center, inner sep=0}, between={0.6}{1}, from=1-2, to=3-2]
        \arrow["\lrcorner"{anchor=center, pos=0.125}, draw=none, from=1-2, to=3-4]
        \arrow[from=1-4, to=3-4]
        \arrow[dashed, from=2-1, to=1-2]
        \arrow[from=2-1, to=2-3]
        \arrow[from=2-1, to=4-1]
        \arrow[""{name=1, anchor=center, inner sep=0}, "\lrcorner"{anchor=center, pos=0.125}, draw=none, from=2-1, to=4-3]
        \arrow[from=2-3, to=1-4]
        \arrow["\lrcorner"{anchor=center, pos=0.125}, draw=none, from=2-3, to=3-4]
        \arrow[from=2-3, to=4-3]
        \arrow[between={0.7}{1}, from=3-2, to=3-4]
        \arrow[from=4-1, to=3-2]
        \arrow["\lrcorner"{anchor=center, pos=0.125, rotate=45}, draw=none, from=4-1, to=3-4]
        \arrow["{(x,y)}"{description}, from=4-1, to=4-3]
        \arrow[from=4-3, to=3-4]
        \arrow[between={0}{0.9}, no head, from=1-2, to=0]
        \arrow["{(Fx,Fy)}"{description, pos=0.3}, between={0.1}{0.6}, no head, from=1, to=3-4]
    \end{tikzcd}
\end{equation}
\medskip

\nid and this map is an equivalence — matching the usual notion of fully-faithfulness in ordinary category theory. In particular, the induced map between homotopy categories
\[\ho(F)\maps \ho(\cF) \longra \ho(\cG) \]
is fully-faithful. 
\end{observation}

\begin{observation}
    Let $\cF$ be a Segal space. We define an equivalence relation on the set of path components of the space of objects $\pi_0\cF_0$, by determining 
    \begin{equation}
        x \sim y, \text{ if and only if there is } f\maps x \ra y  
    \end{equation} 
    with $f$ a homotopy equivalence (\Cref{ch2:: def: hom eq}). Note the subtle distinction between being path connected in $\cF_{0}$ and being connected by an equivalence. They do not always agree, and that is at the heart of our next topic of discussion: \emph{completeness}.
\end{observation}

Before we can discuss completeness, what it means, and why we need it, let us first define what it means for a map of Segal spaces to be essentially surjective.

\begin{definition}
    Let $F \maps \cF \longra \cG$ be a map of Segal spaces. We declare $F$ to be \define{essentially surjective} every object in $\cG$ is equivalent to one in the image of $F$ or if the induced map 
    \begin{equation}
    \begin{tikzcd}[ampersand replacement=\&]
        {\cF_0} \&\& {\cG_0} \\
        {\pi_0\cF_0} \&\& {\pi_{0}\cG_{0}} \\
        {\pi_{0}\cF_{0}{/\sim}} \&\& {\pi_{0}\cG_{0}/\sim}
        \arrow[from=1-1, to=1-3]
        \arrow[from=1-1, to=2-1]
        \arrow[from=1-3, to=2-3]
        \arrow[from=2-1, to=2-3]
        \arrow[from=2-1, to=3-1]
        \arrow[from=2-3, to=3-3]
        \arrow[two heads, from=3-1, to=3-3]
    \end{tikzcd}
    \end{equation}
    in the bottom of this diagram is surjective.
\end{definition}

\begin{definition}
    We define a map $F\maps \cF \longra \cG$ to be a \define{Dwyer--Kan equivalence} if it is fully-faithful and essentially surjective.
\end{definition}

Since we want $(\oo,1)$-categories to be equivalent if and only if they are Dwyer--Kan equivalent, an immediate question is: 

\begin{enumerate}
    \item[\textbf{Q:}] To which extent do levelwise equivalences between Segal spaces agree with Dwyer--Kan equivalences? 
\end{enumerate}

If the answer to this question were affirmative we could confidently assert that
$\mathsf{SSp}$ was a correct model for the homotopy theory of
$(\oo,1)$-categories. In other words, it would serve as a model for the
$(\oo,1)$-category of $(\oo,1)$-categories. Unfortunately, the answer is no.
We construct in \Cref{ch2:: ex: need complete} a pair of Segal spaces which are
Dwyer--Kan equivalent but fail to be so levelwise. To add some context for how
these examples are constructed, let us introduce a generalization of the
simplicial nerve $\sN$, called the classifying diagram, due to Rezk
(\cite{Rezk01}).

\begin{definition}
    \label{ch2:: def: rezk nerve}
    Let $\sC$ be an ordinary category, and let $\sW \subseteq \sC$ be a wide subcategory of $\sC$. We define the following simplicial space which we call \define{Rezk's classifying diagram} of $\sC$ relative to $\sW$: 
    \begin{equation}
        \sN^{\cR}(\sC,\sW) \maps \bDelta^{\op} \longra \sset,\quad [m] \longmapsto \sN\big(\fun^{\sW}([m],\sC) \big).
    \end{equation} 
    Here $\fun^{\sW}([m],\sC)$ denotes the category whose objects are all
    functors $[m]\to \sC$, and whose morphisms are the natural transformations
    whose components lie in $\sW$.
\end{definition}

\begin{remark}
    We can describe the $n$-simplices of the $m$th space of $\sN^{\cR}(\sC,\sW)$ as an $m\times n$ commutative grid:
    \begin{equation}
       \begin{tikzcd}[ampersand replacement=\&]
        {c_{0,0}} \&\& {c_{0,1}} \&\& \cdots \&\& {c_{0,n}} \\
        {c_{1,0}} \&\& {c_{1,1}} \&\& \cdots \&\& {c_{1,n}} \\
        \vdots \&\& \vdots \&\& \vdots \&\& \vdots \\
        {c_{m,0}} \&\& {c_{m,1}} \&\& \cdots \&\& {c_{m,n}}
        \arrow[from=1-1, to=1-3]
        \arrow["\shortmid"{marking}, from=1-1, to=2-1]
        \arrow[from=1-3, to=1-5]
        \arrow["\shortmid"{marking}, from=1-3, to=2-3]
        \arrow[from=1-5, to=1-7]
        \arrow["\shortmid"{marking}, from=1-5, to=2-5]
        \arrow["\shortmid"{marking}, from=1-7, to=2-7]
        \arrow[from=2-1, to=2-3]
        \arrow["\shortmid"{marking}, from=2-1, to=3-1]
        \arrow[from=2-3, to=2-5]
        \arrow["\shortmid"{marking}, from=2-3, to=3-3]
        \arrow[from=2-5, to=2-7]
        \arrow["\shortmid"{marking}, from=2-5, to=3-5]
        \arrow["\shortmid"{marking}, from=2-7, to=3-7]
        \arrow["\shortmid"{marking}, from=3-1, to=4-1]
        \arrow["\shortmid"{marking}, from=3-3, to=4-3]
        \arrow["\shortmid"{marking}, from=3-5, to=4-5]
        \arrow["\shortmid"{marking}, from=3-7, to=4-7]
        \arrow[from=4-1, to=4-3]
        \arrow[from=4-3, to=4-5]
        \arrow[from=4-5, to=4-7]
    \end{tikzcd} 
    \end{equation}
    where each $c_{i,j} \in \sC$ for $(i,j)\in [m] \times [n]$, and the vertical marked arrows are in $\sW$, while the horizontal unmarked ones are arbitrary morphisms in $\sC$.
\end{remark}

\begin{observation}\label{ch2:: obs: nerves agree}
If $\sC_{\mathrm{id}}\subseteq \sC$ denotes the wide subcategory containing
only identity morphisms, then Rezk's classifying diagram
$\sN^{\cR}(\sC,\sC_{\mathrm{id}})$ recovers the simplicial nerve $\sN(\sC)$.
\end{observation}


\begin{definition}
    Let $\sC$ be a category, as a way of simplifying notation and naming, let us define the Rezk nerve to be \[\sN^{\cR}(\sC) \,\, \overset{\sf def}{=} \,\, \sN^{\cR}(\sC,\sC^{\simeq})\] the Rezk classifying diagram of $\sC$ relative to its maximal underlying groupoid.
\end{definition}


\begin{proposition}\label{ch2:: prop: Rezk nerve equiv disc nerve}\cite[Proposition 2.2.16]{Zha13}
Given a category $\sC$, there is a canonical inclusion of Segal spaces 
\begin{equation}
    \sN(\sC) \xlonghookra{\,\,\,\,} \sN^{\cR}(\sC)
\end{equation}    
which is a Dwyer--Kan equivalence.
\end{proposition}



\begin{example}\label{ch2:: ex: need complete}
    Let us consider $G$ a finite group. Let $\sB G$ denote the category with a single object $*$ and $\Aut_{\sB G}(*) = G$. There is a pair of Segal spaces we can construct out of this category:
\begin{enumerate}
    \item \textbf{Discrete nerve}: We can define the simplicial nerve of $\sB G$, or as noted in \Cref{ch2:: obs: nerves agree} the Rezk classifying diagram of $\sB G$ relative to $\{\id_G\}$. This is a constant simplicial space, i.e., every $m$th level is given by the discrete set \[\sN(\fun^{\{\id_* \}}([m],\sB G))\,\simeq\, G^{\times m}.\] In particular, the $0$th space is the one-point set.
    \item \textbf{Rezk nerve}: Consider the Rezk nerve $\sN^{\cR}(\sB G)$. Its $0$th space is given by $\sN(\fun^{\sB G}([0],\sB G))$, in particular 
    \begin{equation}
        \sN(\fun^{\sB G}([0],\sB G))_{\bullet} \,\, \simeq \,\, G^{\times \bullet}.
    \end{equation} 
\end{enumerate}

\nid By \Cref{ch2:: prop: Rezk nerve equiv disc nerve} we have that 
\begin{equation}
    \sN(\sB G) \hookra \sN^{\cR}(\sB G, \sB G)
\end{equation}
is a Dwyer--Kan equivalence. On the other hand, it is manifest that they are not levelwise equivalent.
\end{example}

\Cref{ch2:: ex: need complete} suggests that we should apply a left Bousfield
localization of $\mathsf{SSp}$ along the Dwyer--Kan equivalences to obtain the
correct homotopy theory of $(\oo,1)$-categories. What Rezk shows in
\cite{Rezk01}, is that this is a \emph{reflective localization}, meaning that
the localization functor admits a fully-faithful right adjoint. The image of
this right adjoint contains the Segal spaces whose space of objects captures
the correct homotopy type i.e., the homotopy type of homotopy equivalences.
These Segal spaces will be called complete, and the next sequence of
definitions gives a more precise characterization of these objects.

\begin{definition}\label{ch2:: def: complete Segal}
    Let $\cF$ be a Segal space. We say that $\cF$ is a \define{complete Segal space} if the functor
    \begin{equation}
        \Map^{h}(-, \cF) \maps \mathsf{SSp}^{\op} \longra \sset
    \end{equation}
    sends the following commutative diagram 
    \begin{equation}
        \label{eq complete}
        \begin{tikzcd}[ampersand replacement=\&]
        {\bDelta[\{0<2\}] \coprod \bDelta[\{1<3\}]} \&\& {\bDelta[\{0<1<2<3\}]} \\
        \\
        {\bDelta[0] \coprod\bDelta[0]} \&\& {\bDelta[0]}
        \arrow[from=1-1, to=1-3]
        \arrow[from=1-1, to=3-1]
        \arrow[from=1-3, to=3-3]
        \arrow[from=3-1, to=3-3]
    \end{tikzcd}
    \end{equation}
    to a homotopy pullback square in $\sset$.
\end{definition}

To make sense of \Cref{ch2:: def: complete Segal}, we must define what we mean by the space of homotopy equivalences. 

\begin{definition}
    Let $\cF$ be a Segal space. We define its \define{space of homotopy equivalences} to be
    \begin{equation}
        \cF_{1}^{\sf eq} \,\, \overset{\sf def}{=} \,\, \Map_{\mathsf{SSp}}^h(J,\cF)
    \end{equation}
    where $J$ is the following Segal space\footnote{Note that every object is cofibrant and that $\cF$ is fibrant, so there is no need for replacements to achieve the correct mapping space.} 
    \begin{equation}
        J \,\, \overset{\sf def}{=} \,\, \Delta[\{0<1<2<3\}] \coprod_{\bDelta[0]\coprod \bDelta[0]} \left( \bDelta[\{0<2\}] \sqcup \bDelta[\{1<3\}] \right),
    \end{equation}
    defined by contracting the $(0\ra 2)$ and the $(1\ra 3)$ edges of the 3-simplex. It is the pushout of diagram (\ref{eq complete}). The resulting simplicial space can be visualized as follows:
    \begin{equation}
        \begin{tikzcd}[ampersand replacement=\&]
            \&\& 2 \\
            \\
            0 \&\&\&\& 3 \&\& 0 \& 1 \& 0 \& 1 \\
            \&\& 1
            \arrow[from=1-3, to=3-5]
            \arrow[from=3-1, to=1-3]
            \arrow[between={0.6}{1}, dotted, from=3-1, to=3-5]
            \arrow[between={0}{0.4}, dotted, no head, from=3-1, to=3-5]
            \arrow[from=3-1, to=4-3]
            \arrow[between={0.2}{0.8}, squiggly, from=3-5, to=3-7]
            \arrow[from=3-7, to=3-8]
            \arrow[curve={height=-20pt}, equals, from=3-7, to=3-9]
            \arrow[from=3-8, to=3-9]
            \arrow[curve={height=20pt}, equals, from=3-8, to=3-10]
            \arrow[from=3-9, to=3-10]
            \arrow[from=4-3, to=1-3]
            \arrow[from=4-3, to=3-5]
        \end{tikzcd}
    \end{equation}
\end{definition}

\begin{observation}
    Note how giving a map $J \longra \cF$ is equivalent to specifying a morphism $f\maps x \ra y$ together with a left inverse $l$ and a right inverse $r$
    \begin{equation}
        \begin{tikzcd}[ampersand replacement=\&]
	y \& x \& y \& x.
        \arrow["l", from=1-1, to=1-2]
        \arrow[curve={height=-20pt}, equals, from=1-1, to=1-3]
        \arrow["f", from=1-2, to=1-3]
        \arrow[curve={height=20pt}, equals, from=1-2, to=1-4]
        \arrow["r", from=1-3, to=1-4]
    \end{tikzcd}
    \end{equation}
    thereby asserting that $f$ must be a homotopy equivalence. Note also that fixing $f\maps x \ra y$, the homotopy fiber of the map $ \cF^{\sf eq}_{1} \longra \cF_1$ induced by the inclusion of $\bDelta[\{1<2\}] \hookra J$:
    \begin{equation}
        \cF^{\sf eq}_{1} \overset{h}{\underset{\cF_1}{\times}} \{f \maps x \ra y\},
    \end{equation}   
    which contains all left and right homotopy inverses to $f$, might well be empty. This simplicial set can only be either contractible or empty, as one would hope. A proof of this fact can be found, for example, in \cite[Corollary 2.44]{NR18}.
\end{observation}

\begin{observation}
    Revisiting \Cref{ch2:: def: complete Segal}, we see that a complete Segal space $\cF$ is one for which the dashed arrow 
    \begin{equation}
        \begin{tikzcd}[ampersand replacement=\&]
            {\cF([0])} \\
            \& {\cF_1^{\sf eq}} \&\& {\cF([\{0<1<2<3\}])} \\
            \\
            \& {\cF([0])\times \cF([0])} \&\& {\cF([\{0<2\}])\times \cF([\{1<3\}])}
            \arrow["{\rotatebox[origin=c]{-35}{$\simeq$}}", dashed, from=1-1, to=2-2]
            \arrow[curve={height=-18pt}, from=1-1, to=2-4]
            \arrow[curve={height=18pt}, from=1-1, to=4-2]
            \arrow[from=2-2, to=2-4]
            \arrow[from=2-2, to=4-2]
            \arrow["\lrcorner"{anchor=center, pos=0.125}, draw=none, from=2-2, to=4-4]
            \arrow[from=2-4, to=4-4]
            \arrow[from=4-2, to=4-4]
        \end{tikzcd}
    \end{equation}
    is a weak equivalence of simplicial sets. Complete Segal spaces $\cF$ are thus Segal spaces for which the inclusion
    \begin{equation}
        \cF_0 \longra \cF_{1}^{\sf eq}
    \end{equation} 
    is an equivalence.
\end{observation}

Finally, we have the following result of Rezk.

\begin{proposition}\label{ch2:: prop: Rezk-model-structure}
    There is a model structure $\mathsf{CSSp}$, obtained by left Bousfield
    localizing $\mathsf{SSp}$ along the terminal map
    \[\label{local map complete} J \longra \bDelta[0],\]
    and it has the following properties:
    \begin{enumerate}
        \item The cofibrations are the monomorphisms.
        \item The fibrant objects are the complete Segal spaces.
        \item The weak equivalences among complete Segal spaces are the levelwise equivalences.  
    \end{enumerate}
    Additionally, the model structure is left proper, combinatorial and Cartesian closed.
\end{proposition}


\begin{theorem}\cite[Theorem 7.7]{Rezk01}
    Let $\cF$ and $\cG$ be Segal spaces. A morphism $F \maps \cF \longra \cG$
    is a weak equivalence in $\mathsf{CSSp}$ if and only if it is a
    Dwyer--Kan equivalence. 
\end{theorem}

In particular, this means that despite $\sN(\sB G)$ not being a complete Segal
space, therefore not fibrant in $\mathsf{CSSp}$, the Dwyer--Kan equivalence
\begin{equation}
    \sN(\sB G) \longra \sN^{\cR}(\sB G, \sB G)
\end{equation}
into its Rezk nerve is a weak equivalence, as desired.

\begin{remark}
    If necessary, we write $\cC \mathsf{SSp}\subset \mathsf{SSp}$ for the full
    subcategory of complete Segal spaces, as a way of distinguishing it from
    the localized model category $\mathsf{CSSp}$.
\end{remark}

\begin{remark}
    As a consequence of the fact that $\mathsf{CSSp}$ is Cartesian closed we
    have that given a Segal space $\cF$ and a complete Segal space $\cG$, the
    simplicial space $\fun(\cF,\cG)$ is also a complete Segal space. We will
    refer to it as the $(\oo,1)$-category of functors between $\cF$ and $\cG$.
\end{remark}





\section{Interval spaces}\label{ch2:: subsec: interval_spaces}

When working with complete Segal spaces as a model for $(\oo,1)$-categories, the combinatorial engine, the object doing the work of encoding all the higher homotopical data, is provided by the simplex category $\bDelta$. Yet, as the literature on $(\oo,1)$-categories has evolved, the community has found it useful to adopt variants, which facilitate their direct application. As we will unravel throughout the remainder of the thesis, there is a variant of $\bDelta$, called the \emph{interval category}, which provides a more flexible framework for working with $(\oo,1)$-categories of cobordisms and higher versions thereof. As such, this subsection is dedicated to introducing its mechanics.


\begin{definition}\label{ch2:: def: interval_category}
    The interval category $\mathsf{Int}$ is the category whose objects are finite disjoint unions of open intervals $I \subseteq [0,1]$ containing the endpoints $\{0,1\}$ but are \textbf{not} the whole interval $[0,1]$, and whose morphisms are the inclusions $(I \longra J):= I\subseteq J$ of open intervals. 
\end{definition}

\begin{observation}\label{ch2:: obs: poset_observation}
    The category $\mathsf{Int}$ is a poset. Manifestly, $\Int$ is a full subcategory spanned by a subset of objects of $\cU([0,1])$, the poset of opens of the unit interval. 
\end{observation}

For convenience of the reader and as a means to stipulate notation for the rest of the chapter, we introduce the following diagrammatic means of depicting the objects of $\mathsf{Int}$:
\medskip
\begin{equation}\label{ch2:: eq: interval_object}
    I \,\, = \,\, \tikzfig{figures/Ch2/INT1}{\begin{tikzpicture}
	\begin{pgfonlayer}{nodelayer}
		\node [style=none] (0) at (-7, 0) {};
		\node [style=none] (1) at (8, 0) {};
		\node [style=none] (2) at (-5, 0) {};
		\node [style=none] (3) at (-3, 0) {};
		\node [style=none] (4) at (-2.5, 0) {};
		\node [style=none] (5) at (-1.5, 0) {};
		\node [style=none] (6) at (2.5, 0) {};
		\node [style=none] (7) at (6.75, 0) {};
		\node [style=endpoint] (8) at (-7, 0) {};
		\node [style=endpoint] (9) at (8, 0) {};
		\node [style=none] (10) at (-7, 0.75) {$0$};
		\node [style=none] (11) at (8, 0.75) {$1$};
	\end{pgfonlayer}
	\begin{pgfonlayer}{edgelayer}
		\draw (0.center) to (1.center);
		\draw [style={open_int}] (0.center) to (2.center);
		\draw [style={open_int}] (3.center) to (4.center);
		\draw [style={open_int}] (5.center) to (6.center);
		\draw [style={open_int}] (7.center) to (1.center);
	\end{pgfonlayer}
\end{tikzpicture}
}
\end{equation}
\medskip

\noindent The diagram above depicts an object $I$ in $\mathsf{Int}$. It
consists of a finite disjoint union of four open intervals which contain both
endpoints but do not exhaust the whole interval. Unless we need to specify the
start- and end-points of the intervals, we will refrain from doing so, which,
given the topological nature of our applications, will prove to be a reasonable
default.

Additionally, we will often depict morphisms in this category by writing diagrams like the following one: 
\medskip
\begin{equation}\label{ch2:: eq: interval_morphism}
    \tikzfig{figures/Ch2/INT2}{\begin{tikzpicture}
	\begin{pgfonlayer}{nodelayer}
		\node [style=none] (0) at (-7.5, 3) {};
		\node [style=none] (1) at (7.5, 3) {};
		\node [style=none] (2) at (-5.5, 3) {};
		\node [style=none] (3) at (-3.5, 3) {};
		\node [style=none] (4) at (-3, 3) {};
		\node [style=none] (5) at (-2, 3) {};
		\node [style=none] (6) at (2, 3) {};
		\node [style=none] (7) at (6.25, 3) {};
		\node [style=endpoint] (8) at (-7.5, 3) {};
		\node [style=endpoint] (9) at (7.5, 3) {};
		\node [style=none] (10) at (-7.5, 3.75) {$0$};
		\node [style=none] (11) at (7.5, 3.75) {$1$};
		\node [style=none] (12) at (-7.5, -2) {};
		\node [style=none] (13) at (7.5, -2) {};
		\node [style=none] (14) at (-4.5, -2) {};
		\node [style=none] (15) at (-3.75, -2) {};
		\node [style=none] (16) at (-2.75, -2) {};
		\node [style=none] (17) at (-2, -2) {};
		\node [style=none] (18) at (4.5, -2) {};
		\node [style=none] (19) at (5.75, -2) {};
		\node [style=endpoint] (20) at (-7.5, -2) {};
		\node [style=endpoint] (21) at (7.5, -2) {};
		\node [style=none] (22) at (-7.5, -2.75) {$0$};
		\node [style=none] (23) at (7.5, -2.75) {$1$};
		\node [style=none] (24) at (0, 4) {$I$};
		\node [style=none] (25) at (0, -3) {$J$};
	\end{pgfonlayer}
	\begin{pgfonlayer}{edgelayer}
		\draw (0.center) to (1.center);
		\draw [style={open_int}] (0.center) to (2.center);
		\draw [style={open_int}] (3.center) to (4.center);
		\draw [style={open_int}] (5.center) to (6.center);
		\draw [style={open_int}] (7.center) to (1.center);
		\draw (12.center) to (13.center);
		\draw [style={open_int}] (12.center) to (14.center);
		\draw [style={open_int}] (15.center) to (16.center);
		\draw [style={open_int}] (17.center) to (18.center);
		\draw [style={open_int}] (19.center) to (13.center);
		\draw [style=ts-dashed] (8) to (20);
		\draw [style=ts-dashed] (9) to (21);
		\draw [style=ts-dashed-red] (2.center) to (14.center);
		\draw [style=ts-dashed-red] (3.center) to (15.center);
		\draw [style=ts-dashed-red] (4.center) to (16.center);
		\draw [style=ts-dashed-red] (5.center) to (17.center);
		\draw [style=ts-dashed-red] (6.center) to (18.center);
		\draw [style=ts-dashed-red] (7.center) to (19.center);
	\end{pgfonlayer}
\end{tikzpicture}
}
\end{equation}
\medskip

\noindent which codifies a unique morphism from the object $I$ to the object $J$. Every dashed edge is merely meant to indicate how these objects relate to one another and are not data associated to the objects or the morphisms themselves. In fact, if the morphism is clear we will sometimes avoid the dashed edges.

\begin{definition}\label{ch2:: def: face_maps_intervals}
    Let $I\subseteq [0,1]$ be an object of $\mathsf{Int}$. We can identify its complement with a disjoint union of finitely many closed intervals of $[0,1]$:
    \begin{equation}
        [0,1]\setminus I = \coprod_{k = 0}^{n} A_k
    \end{equation}
    where $A_k$ are closed subsets of $[0,1]$ which do not contain $0$ or $1$. Let $\sU(\{A_i\}_{i=0}^n)$ denote the open subset which is the complement of the union of closed subsets $A_i$ in $[0,1]$. We define 
    \begin{equation}
        \alpha_{i}^{I}\maps \sU(\{A_1,\cdots,A_n\}) \longrightarrow \sU(\{A_1,\cdots,\widehat{A_{i}},\cdots,A_n\})
    \end{equation} for every $i=0,\cdots,n$ as the map which \emph{forgets} the $i$-th closed interval $A_i$. We will call these maps \textbf{gluings}.


    \medskip

    \nid Let us define 
    \begin{equation}\label{ch2:: eq: degeneracy_map_interval}
        \beta_{i}^{I}\maps \sU(\{A_1,\ldots, A_n\}) \longleftarrow \sU(\{A_1,\ldots, \mathsf{Hull}(A_i,A_{i+1}),\ldots, A_n\}) 
    \end{equation}
    for every $i=1,\ldots, n-1$, as the inclusion which \emph{merges} the intervals $A_i$ and $A_{i+1}$, replacing it with its \emph{hull}. We will call these maps \textbf{refinements}.
\end{definition}

\medskip

\begin{remark}
    By $\mathsf{Hull}(A_{1},\cdots, A_n)$ we mean the smallest closed interval containing all the closed intervals $A_i$ for $i=1,\cdots,n$. In particular, if $A_i = [a_i,b_i]$ and $A_{i+1} = [a_{i+1},b_{i+1}]$, then $\mathsf{Hull}(A_i,A_{i+1}) = [\min(a_i,a_{i+1}),\max(b_i,b_{i+1})]$.
\end{remark}


\begin{example}
    Let us pictorially exhibit examples of face maps in $\mathsf{Int}$, so as to get the reader a better feel for how they relate to their analogues in $\bDelta^{\op}$. Let $I\in\mathsf{Obj}(\mathsf{Int})$ be depicted as
    \begin{equation}\label{ch2:: eq: face_map_1}
        \medskip
        \sU(\{A_0,A_1,A_2\}) \,\, = \,\, \tikzfig{figures/Ch2/INT3}{\begin{tikzpicture}
	\begin{pgfonlayer}{nodelayer}
		\node [style=none] (0) at (-7.5, 0) {};
		\node [style=none] (1) at (7.5, 0) {};
		\node [style=none] (2) at (-5.5, 0) {};
		\node [style=none] (3) at (-3.5, 0) {};
		\node [style=none] (4) at (-3, 0) {};
		\node [style=none] (5) at (-2, 0) {};
		\node [style=none] (6) at (2, 0) {};
		\node [style=none] (7) at (6.25, 0) {};
		\node [style=endpoint] (8) at (-7.5, 0) {};
		\node [style=endpoint] (9) at (7.5, 0) {};
		\node [style=none] (12) at (-4.5, 0.5) {$\scriptstyle{A_0}$};
		\node [style=none] (13) at (-2.5, 0.5) {$\scriptstyle{A_1}$};
		\node [style=none] (14) at (4.25, 0.5) {$\scriptstyle{A_2}$};
	\end{pgfonlayer}
	\begin{pgfonlayer}{edgelayer}
		\draw (0.center) to (1.center);
		\draw [style={open_int}] (0.center) to (2.center);
		\draw [style={open_int}] (3.center) to (4.center);
		\draw [style={open_int}] (5.center) to (6.center);
		\draw [style={open_int}] (7.center) to (1.center);
	\end{pgfonlayer}
\end{tikzpicture}
}.
        \medskip
    \end{equation}
    Then we have the following triple of face maps: 
    \begin{equation}
        \medskip
        \tikzfig{figures/Ch2/INT5}{\begin{tikzpicture}
	\begin{pgfonlayer}{nodelayer}
		\node [style=none] (0) at (-7.5, 6) {};
		\node [style=none] (1) at (7.5, 6) {};
		\node [style=none] (2) at (-5.5, 6) {};
		\node [style=none] (3) at (-3.5, 6) {};
		\node [style=none] (4) at (-3, 6) {};
		\node [style=none] (5) at (-2, 6) {};
		\node [style=none] (6) at (2, 6) {};
		\node [style=none] (7) at (6.25, 6) {};
		\node [style=endpoint] (8) at (-7.5, 6) {};
		\node [style=endpoint] (9) at (7.5, 6) {};
		\node [style=none] (12) at (-7.5, 1) {};
		\node [style=none] (13) at (7.5, 1) {};
		\node [style=none] (14) at (-3.5, 1) {};
		\node [style=none] (15) at (-3.5, 1) {};
		\node [style=none] (16) at (-3, 1) {};
		\node [style=none] (17) at (-2, 1) {};
		\node [style=none] (18) at (2, 1) {};
		\node [style=none] (19) at (6.25, 1) {};
		\node [style=endpoint] (20) at (-7.5, 1) {};
		\node [style=endpoint] (21) at (7.5, 1) {};
		\node [style=none] (24) at (-10.5, 6) {$\scriptstyle{\sU(\{A_0,A_1,A_2\})} \,\,\, = $};
		\node [style=none] (26) at (-4.5, 6.5) {$\scriptstyle{A_0}$};
		\node [style=none] (27) at (-2.5, 6.5) {$\scriptstyle{A_1}$};
		\node [style=none] (28) at (4.25, 6.5) {$\scriptstyle{A_2}$};
		\node [style=none] (31) at (4.5, 0.5) {$\scriptstyle{A_2}$};
		\node [style=none] (52) at (-10.5, 1) {$\scriptstyle{\sU(\{\widehat{A_0},A_1,A_2\})}\,\,\, = $};
		\node [style=none] (53) at (-10.5, 5.5) {};
		\node [style=none] (54) at (-10.5, 1.75) {};
		\node [style=none] (55) at (-11.25, 3.5) {$\alpha_0$};
		\node [style=none] (56) at (-2.5, 0.5) {$\scriptstyle{A_1}$};
	\end{pgfonlayer}
	\begin{pgfonlayer}{edgelayer}
		\draw (0.center) to (1.center);
		\draw [style={open_int}] (0.center) to (2.center);
		\draw [style={open_int}] (3.center) to (4.center);
		\draw [style={open_int}] (5.center) to (6.center);
		\draw [style={open_int}] (7.center) to (1.center);
		\draw (12.center) to (13.center);
		\draw [style={open_int}] (12.center) to (14.center);
		\draw [style={open_int}] (15.center) to (16.center);
		\draw [style={open_int}] (17.center) to (18.center);
		\draw [style={open_int}] (19.center) to (13.center);
		\draw [style=ts-dashed] (8) to (20);
		\draw [style=ts-dashed] (9) to (21);
		\draw [style=ts-dashed-red] (6.center) to (18.center);
		\draw [style=ts-dashed-red] (7.center) to (19.center);
		\draw [style=dir] (53.center) to (54.center);
		\draw [style=ts-dashed-red] (4.center) to (16.center);
		\draw [style=ts-dashed-red] (5.center) to (17.center);
	\end{pgfonlayer}
\end{tikzpicture}
}
        \medskip
    \end{equation}
    which forgets $A_0$, and
    \begin{equation}
        \medskip
        \tikzfig{figures/Ch2/INT4}{\begin{tikzpicture}
	\begin{pgfonlayer}{nodelayer}
		\node [style=none] (0) at (-7.5, 6) {};
		\node [style=none] (1) at (7.5, 6) {};
		\node [style=none] (2) at (-5.5, 6) {};
		\node [style=none] (3) at (-3.5, 6) {};
		\node [style=none] (4) at (-3, 6) {};
		\node [style=none] (5) at (-2, 6) {};
		\node [style=none] (6) at (2, 6) {};
		\node [style=none] (7) at (6.25, 6) {};
		\node [style=endpoint] (8) at (-7.5, 6) {};
		\node [style=endpoint] (9) at (7.5, 6) {};
		\node [style=none] (12) at (-7.5, 1) {};
		\node [style=none] (13) at (7.5, 1) {};
		\node [style=none] (14) at (-5.5, 1) {};
		\node [style=none] (15) at (-3.5, 1) {};
		\node [style=none] (16) at (-2.75, 1) {};
		\node [style=none] (17) at (-2.75, 1) {};
		\node [style=none] (18) at (2, 1) {};
		\node [style=none] (19) at (6.25, 1) {};
		\node [style=endpoint] (20) at (-7.5, 1) {};
		\node [style=endpoint] (21) at (7.5, 1) {};
		\node [style=none] (24) at (-10.5, 6) {$\scriptstyle{\sU(\{A_0,A_1,A_2\})} \,\,\, = $};
		\node [style=none] (26) at (-4.5, 6.5) {$\scriptstyle{A_0}$};
		\node [style=none] (27) at (-2.5, 6.5) {$\scriptstyle{A_1}$};
		\node [style=none] (28) at (4.25, 6.5) {$\scriptstyle{A_2}$};
		\node [style=none] (29) at (-4.5, 0.5) {$\scriptstyle{A_0}$};
		\node [style=none] (31) at (4.5, 0.5) {$\scriptstyle{A_2}$};
		\node [style=none] (52) at (-10.5, 1) {$\scriptstyle{\sU(\{A_0,\widehat{A_1},A_2\})}\,\,\, = $};
		\node [style=none] (53) at (-10.5, 5.5) {};
		\node [style=none] (54) at (-10.5, 1.75) {};
		\node [style=none] (55) at (-11.25, 3.5) {$\alpha_1$};
	\end{pgfonlayer}
	\begin{pgfonlayer}{edgelayer}
		\draw (0.center) to (1.center);
		\draw [style={open_int}] (0.center) to (2.center);
		\draw [style={open_int}] (3.center) to (4.center);
		\draw [style={open_int}] (5.center) to (6.center);
		\draw [style={open_int}] (7.center) to (1.center);
		\draw (12.center) to (13.center);
		\draw [style={open_int}] (12.center) to (14.center);
		\draw [style={open_int}] (15.center) to (16.center);
		\draw [style={open_int}] (17.center) to (18.center);
		\draw [style={open_int}] (19.center) to (13.center);
		\draw [style=ts-dashed] (8) to (20);
		\draw [style=ts-dashed] (9) to (21);
		\draw [style=ts-dashed-red] (2.center) to (14.center);
		\draw [style=ts-dashed-red] (3.center) to (15.center);
		\draw [style=ts-dashed-red] (6.center) to (18.center);
		\draw [style=ts-dashed-red] (7.center) to (19.center);
		\draw [style=dir] (53.center) to (54.center);
	\end{pgfonlayer}
\end{tikzpicture}
}
        \medskip
    \end{equation}
    which forgets $A_1$, and finally
    \begin{equation}
        \medskip
        \tikzfig{figures/Ch2/INT6}{\begin{tikzpicture}
	\begin{pgfonlayer}{nodelayer}
		\node [style=none] (0) at (-7.5, 6) {};
		\node [style=none] (1) at (7.5, 6) {};
		\node [style=none] (2) at (-5.5, 6) {};
		\node [style=none] (3) at (-3.5, 6) {};
		\node [style=none] (4) at (-3, 6) {};
		\node [style=none] (5) at (-2, 6) {};
		\node [style=none] (6) at (2, 6) {};
		\node [style=none] (7) at (6.25, 6) {};
		\node [style=endpoint] (8) at (-7.5, 6) {};
		\node [style=endpoint] (9) at (7.5, 6) {};
		\node [style=none] (12) at (-7.5, 1) {};
		\node [style=none] (13) at (7.5, 1) {};
		\node [style=none] (14) at (-5.5, 1) {};
		\node [style=none] (15) at (-3.5, 1) {};
		\node [style=none] (16) at (-3, 1) {};
		\node [style=none] (17) at (-2, 1) {};
		\node [style=none] (18) at (6.25, 1) {};
		\node [style=none] (19) at (6.25, 1) {};
		\node [style=endpoint] (20) at (-7.5, 1) {};
		\node [style=endpoint] (21) at (7.5, 1) {};
		\node [style=none] (24) at (-10.5, 6) {$\scriptstyle{\sU(\{A_0,A_1,A_2\})} \,\,\, = $};
		\node [style=none] (26) at (-4.5, 6.5) {$\scriptstyle{A_0}$};
		\node [style=none] (27) at (-2.5, 6.5) {$\scriptstyle{A_1}$};
		\node [style=none] (28) at (4.25, 6.5) {$\scriptstyle{A_2}$};
		\node [style=none] (52) at (-10.5, 1) {$\scriptstyle{\sU(\{A_0,A_1,\widehat{A_2}\})}\,\,\, = $};
		\node [style=none] (53) at (-10.5, 5.5) {};
		\node [style=none] (54) at (-10.5, 1.75) {};
		\node [style=none] (55) at (-11.25, 3.5) {$\alpha_2$};
		\node [style=none] (56) at (-2.5, 0.5) {$\scriptstyle{A_1}$};
		\node [style=none] (57) at (-4.5, 0.5) {$\scriptstyle{A_0}$};
	\end{pgfonlayer}
	\begin{pgfonlayer}{edgelayer}
		\draw (0.center) to (1.center);
		\draw [style={open_int}] (0.center) to (2.center);
		\draw [style={open_int}] (3.center) to (4.center);
		\draw [style={open_int}] (5.center) to (6.center);
		\draw [style={open_int}] (7.center) to (1.center);
		\draw (12.center) to (13.center);
		\draw [style={open_int}] (12.center) to (14.center);
		\draw [style={open_int}] (15.center) to (16.center);
		\draw [style={open_int}] (17.center) to (18.center);
		\draw [style={open_int}] (19.center) to (13.center);
		\draw [style=ts-dashed] (8) to (20);
		\draw [style=ts-dashed] (9) to (21);
		\draw [style=dir] (53.center) to (54.center);
		\draw [style=ts-dashed-red] (4.center) to (16.center);
		\draw [style=ts-dashed-red] (5.center) to (17.center);
		\draw [style=ts-dashed-red] (2.center) to (14.center);
		\draw [style=ts-dashed-red] (3.center) to (15.center);
	\end{pgfonlayer}
\end{tikzpicture}
}
        \medskip
    \end{equation}
    which forgets $A_2$.
\end{example}

\noindent We introduce the notation above also as a means of aiding our description of the morphisms in $\mathsf{Int}$, given how inherently geometric they are — we like to think of them as a discrete version of the  configuration category of points in the open interval —  but also to initiate the discussion of how it relates to the simplex category $\bDelta$.

\begin{notation}
Let $[i,j] \subseteq \{1,\ldots,n\}$ for $0\leq i < j\leq n$ denote the subset which contains all consecutive integers from $i$ through $j$. For $I = \sU(A_1,\cdots,A_n)$, we denote the maps 
\begin{equation}
    \alpha_{[i,j]}^{I}\maps \sU(A_{1},\ldots, A_{n}) \longra \sU(A_1, \ldots,  A_{i-1}, A_{j+1}, \ldots, A_{n}) 
\end{equation}
which forget the consecutive closed intervals $A_i,\ldots,A_j$.
\end{notation}

\begin{observation}\label{ch2:: obs: interval_category_small}
The interval category $\Int$ is a poset category, so it is small. 
\end{observation}
\begin{remark}
    The opposite $\Int^{\op}$ of the interval category is equivalent to a category with an equally geometric interpretation. Let $\cInt$ be the category whose objects are finite non-empty disjoint unions of connected closed intervals $K\subset [0,1]$ which do not intersect $\{0,1\}$. A morphism $K\longra L$ is given by inclusion — so $\cInt$ is obviously also a poset. The equivalence is defined as 
    \begin{align}
        (-)^{c}\maps \Int^{\op} &\xra{\,\,\,\cong \,\,\,} \cInt \\
        I &\longmapsto \big( I^{c} = [0,1]\setminus I\big).
    \end{align}
    where an inclusion $I\ra J$ naturally induces an inclusion $J^{c} \ra I^{c}$. This equivalence is manifest, since it is bijective poset map.
\end{remark}

\medskip

\nid When convenient, we will use this isomorphism to shift perspectives which suit a given narrative — minding, of course, the fundamental change in variance.

\begin{definition}
    Let 
    \begin{equation}
        \gamma \maps \Int \longra \bDelta^{\op}, \quad (I \subset [0,1])  \longmapsto \pi_{0}(I^{c}),
    \end{equation}
    be the functor which sends an interval $I$ to the set of connected components of its complement in $[0,1]$. Given an inclusion a morphism $I \subseteq J$, $\gamma$ maps 
    \begin{equation}
        \gamma (I \subseteq J) = \big(\,\pi_{0}(J^{c}) \longra \pi_0(I^{c})\,\big).
    \end{equation} 
\end{definition}

\begin{observation}
    We have the following canonical commutative diagram of categories:
    \begin{equation}
        \begin{tikzcd}[ampersand replacement=\&]
            \Int \&\& {\bDelta^{\op}} \\
            \\
            {\cInt^{\op}} \&\& {.}
            \arrow["\gamma", from=1-1, to=1-3]
            \arrow["{\rotatebox[origin=c]{90}{$\sim$}}"', from=1-1, to=3-1]
            \arrow["{\pi_0}"', curve={height=6pt}, dashed, from=3-1, to=1-3]
        \end{tikzcd}
    \end{equation}
\end{observation}

\begin{definition}\label{ch2:: def: updelta cat}
    Let $\sO_{\pm}$ be the following category:
    \begin{itemize}
        \item \textbf{Objects}: Linearly ordered finite sets $S$ with a distinct minimum and maximum.
        \item \textbf{Morphisms}: Order-preserving maps $f \maps S \longra T$ that preserve extrema. 
    \end{itemize}
\end{definition}

\begin{remark}
    Let $[p]\in \bDelta$, note that the set $\Hom_{\bDelta}([p],[1])$ is a linearly ordered finite set with a distinct minimum and maximum. For simplicity, we will denote an object in this set as 
    \begin{equation}
        \lag i_0, \ldots, i_p \rag,
    \end{equation}
    where $0 \leq i_{k} \leq i_{k+1} \leq 1$ for every $0\leq k < p$. We declare
    \begin{equation}
    \lag i_0, \ldots, i_p \rag \leq \lag j_0, \ldots, j_p \rag, \quad \text{ if } \quad i_k \leq j_k        
    \end{equation}
    for every $0\leq k \leq p$. This is a linear order precisely because $[p]$ is linearly ordered with the minimum and maximum elements identified respectively as
    \begin{equation}
        \lag 0, \ldots, 0\rag, \quad \lag 1,\ldots, 1 \rag. 
    \end{equation}
    Equivalently, these maps can be identified with splittings
    \begin{equation}
        [p] \,\, \simeq \,\, [q]\star [p-1-q],
    \end{equation}
    where $q = -1, 0,\ldots, p$ — by identifying $\alpha \maps [p] \ra [1]$ with $\alpha^{-1}([0])\star \alpha^{-1}([1])$.\footnote{We establish the convention that $[-1] = \emptyset$.}
\end{remark}

\begin{notation}
    Let $P$ be a poset. We define $P^{\tr}$ to be the poset obtained by adjoining $P \star [0]$ a maximal element. Analogously, we define $P^{\,\tl}$ to be the poset obtained by adjoining $[0] \star P$ a minimal element. 
\end{notation}

\begin{observation}
There is an identification in $\sO_{\pm}$
\begin{equation}
\Hom_{\bDelta}([p],[1]) \,\, \simeq \,\, [p]^{\,\tl},    
\end{equation}
defined by sending the splitting $[q] \star [p-1-q]$ to $q \in [p]^{\,\tl}$.
\end{observation}

\begin{lemma}
    The functor
    \begin{equation}
        \bDelta^{\op} \longra \sO_{\pm}, \quad \big(\,[p] \mapsto \Hom_{\bDelta}([p],[1]) \, \big), \quad \big(\, \bDelta \ni ([p] \xra{\alpha} [q]) \mapsto \Hom_{\bDelta}([q],[1]) \xra{\alpha^*} \Hom_{\bDelta}([p],[1]) \,\big)
    \end{equation}
    is an equivalence of categories.
\end{lemma}
\begin{proof}
    Surjectivity on objects follows from the fact that given a finite linearly ordered set with distinct minimum and maximum $S$, there is canonical isomorphism $[p]^{\,\tl}$ for some $p\geq 0$. Fully-faithfullness is manifest from the definition.
\end{proof}


\begin{observation}
    We have the following commutative diagram of categories:
    \smallskip
    \begin{equation}\label{diagram combinatorics}
        \begin{tikzcd}[ampersand replacement=\&]
            \&\& {\sO_{\pm}} \\
            \\
            \Int \&\& {\bDelta^{\op}} \\
            \\
            {\cInt^{\op}}
            \arrow["{\pi_{0}}", curve={height=-12pt}, from=3-1, to=1-3]
            \arrow["\gamma", from=3-1, to=3-3]
            \arrow["{\rotatebox[origin=c]{90}{$\sim$}}"', from=3-1, to=5-1]
            \arrow["{(-)^{\,\tl}}"', from=3-3, to=1-3]
            \arrow["{\rotatebox[origin=c]{90}{$\sim$}}", draw=none, from=3-3, to=1-3]
            \arrow["{\pi_0}"', curve={height=6pt}, from=5-1, to=3-3]
        \end{tikzcd}
    \end{equation}
    In particular, given a $J \in \cInt$ we have that $\pi_0(J^c) = \gamma_{0}(J^{c})^{\, \tl}$. As innocuous and simple as this fact is, it will serve us well in the upcoming sections.
\end{observation}

Understanding the yoga of~(\ref{diagram combinatorics}) in more depth and in a graphical way will pay dividends in the next sections, so let us add some elucidating remarks.

\begin{remark}\label{ch2:: rmk:: graph calc}
    The simplex category $\bDelta$ admits a particularly simple graphical calculus. We depict an object $[n]$ as a collection of $[n]$-marked points in a horizontal array:
    \begin{equation}
        \tikzfig{figures/Ch2/INT10}{\begin{tikzpicture}
	\begin{pgfonlayer}{nodelayer}
		\node [style=0 strat] (0) at (-3, 0) {};
		\node [style=0 strat] (1) at (0, 0) {};
		\node [style=0 strat] (2) at (3, 0) {};
		\node [style=none] (3) at (-3, -0.75) {$0$};
		\node [style=none] (4) at (0, -0.75) {$1$};
		\node [style=none] (5) at (3, -0.75) {$2$};
	\end{pgfonlayer}
\end{tikzpicture}
}
    \end{equation}
    \smallskip

    \nid A morphism $\alpha \maps [n] \ra [m]$ is depicted by drawing connecting 1-dimensional open intervals between the marked points. This is best understood through examples:
    
     \begin{equation}
       \tikzfig{figures/Ch2/INT9}{\begin{tikzpicture}
	\begin{pgfonlayer}{nodelayer}
		\node [style=0 strat] (0) at (-1.5, 1) {};
		\node [style=0 strat] (1) at (1.5, 1) {};
		\node [style=0 strat] (2) at (-3, -2) {};
		\node [style=0 strat] (3) at (0, -2) {};
		\node [style=0 strat] (4) at (3, -2) {};
		\node [style=none] (5) at (-1.5, 1.75) {$0$};
		\node [style=none] (6) at (1.5, 1.75) {$1$};
		\node [style=none] (7) at (-3, -2.75) {$0$};
		\node [style=none] (8) at (0, -2.75) {$1$};
		\node [style=none] (9) at (3, -2.75) {$2$};
		\node [style=none] (10) at (-6, 1) {$[1]$};
		\node [style=none] (11) at (-6, -2) {$[2]$};
		\node [style=none] (12) at (-6, 0.25) {};
		\node [style=none] (13) at (-6, -1.25) {};
		\node [style=none] (14) at (-6.75, -0.5) {$d_{1} $};
	\end{pgfonlayer}
	\begin{pgfonlayer}{edgelayer}
		\draw [style=cd-ar] (12.center) to (13.center);
		\draw (2) to (0);
		\draw (4) to (1);
	\end{pgfonlayer}
\end{tikzpicture}
}, \qquad \quad  \tikzfig{figures/Ch2/INT11}{\begin{tikzpicture}
	\begin{pgfonlayer}{nodelayer}
		\node [style=0 strat] (0) at (-1.5, -2) {};
		\node [style=0 strat] (1) at (1.5, -2) {};
		\node [style=0 strat] (2) at (-3, 1) {};
		\node [style=0 strat] (3) at (0, 1) {};
		\node [style=0 strat] (4) at (3, 1) {};
		\node [style=none] (5) at (-1.5, -2.75) {$0$};
		\node [style=none] (6) at (1.5, -2.75) {$1$};
		\node [style=none] (7) at (-3, 1.75) {$0$};
		\node [style=none] (8) at (0, 1.75) {$1$};
		\node [style=none] (9) at (3, 1.75) {$2$};
		\node [style=none] (10) at (-6, 1) {$[2]$};
		\node [style=none] (11) at (-6, -2) {$[1]$};
		\node [style=none] (12) at (-6, 0.25) {};
		\node [style=none] (13) at (-6, -1.25) {};
		\node [style=none] (14) at (-6.75, -0.5) {$s_{0}$};
	\end{pgfonlayer}
	\begin{pgfonlayer}{edgelayer}
		\draw [style=cd-ar] (12.center) to (13.center);
		\draw (2) to (0);
		\draw (3) to (0);
		\draw (4) to (1);
	\end{pgfonlayer}
\end{tikzpicture}
}
    \end{equation}
    \smallskip
    
    \nid We might call this a graphical calculus since if we allow the 1-dimensional arcs to be isotoped, while keeping them anchored at the marked points, we get equivalent representations of the same morphism. To facilitate a better visual interpretation for how the functor $(-)^{\sf \tl}\maps \bDelta^{\op} \longra \sO_{\pm}$ acts on objects and morphisms, we extend the graphical calculus of $\bDelta$ by embedding the $[n]$-marked points into the open interval $(0,1)$.\footnote{The choice of the open interval $(0,1)$ is for practical reasons — any other space homeomorphic to $\RR$ would have done the job.}
    \begin{equation}
        \tikzfig{figures/Ch2/INT10}{\begin{tikzpicture}
	\begin{pgfonlayer}{nodelayer}
		\node [style=0 strat] (0) at (-3, 0) {};
		\node [style=0 strat] (1) at (0, 0) {};
		\node [style=0 strat] (2) at (3, 0) {};
		\node [style=none] (3) at (-3, -0.75) {$0$};
		\node [style=none] (4) at (0, -0.75) {$1$};
		\node [style=none] (5) at (3, -0.75) {$2$};
	\end{pgfonlayer}
\end{tikzpicture}
} \,\,\, \xlonghookra{\,\,\,\,\,\,\,} \,\,\,\tikzfig{figures/Ch2/INT12}{\begin{tikzpicture}
	\begin{pgfonlayer}{nodelayer}
		\node [style=0 strat] (0) at (-3, 0) {};
		\node [style=0 strat] (1) at (0, 0) {};
		\node [style=0 strat] (2) at (3, 0) {};
		\node [style=none] (3) at (-3, -0.75) {$0$};
		\node [style=none] (4) at (0, -0.75) {$1$};
		\node [style=none] (5) at (3, -0.75) {$2$};
		\node [style=none] (6) at (-5.25, 0) {};
		\node [style=none] (7) at (5.5, 0) {};
	\end{pgfonlayer}
	\begin{pgfonlayer}{edgelayer}
		\draw (6.center) to (7.center);
	\end{pgfonlayer}
\end{tikzpicture}
}.
    \end{equation}
    \smallskip

    \nid Additionally, we label the gaps, the 1-dimensional open intervals, between the $[n]$-marked points with the totally ordered set $[n]^{\, \tl}$ as follows
    \begin{equation}
        \tikzfig{figures/Ch2/INT13}{\begin{tikzpicture}
	\begin{pgfonlayer}{nodelayer}
		\node [style=0 strat] (0) at (-3, 0) {};
		\node [style=0 strat] (1) at (0, 0) {};
		\node [style=0 strat] (2) at (3, 0) {};
		\node [style=none] (3) at (-3, -0.75) {$0$};
		\node [style=none] (4) at (0, -0.75) {$1$};
		\node [style=none] (5) at (3, -0.75) {$2$};
		\node [style=none] (6) at (-7, 0) {};
		\node [style=none] (7) at (7, 0) {};
		\node [style=none] (8) at (-5.25, 0.75) {$-\infty$};
		\node [style=0 strat red] (9) at (-5.25, 0) {};
		\node [style=0 strat red] (10) at (-1.5, 0) {};
		\node [style=0 strat red] (11) at (1.75, 0) {};
		\node [style=0 strat red] (12) at (5.25, 0) {};
		\node [style=none] (13) at (-1.5, 0.75) {0};
		\node [style=none] (14) at (1.75, 0.75) {1};
		\node [style=none] (15) at (5.25, 0.75) {2};
	\end{pgfonlayer}
	\begin{pgfonlayer}{edgelayer}
		\draw (6.center) to (7.center);
	\end{pgfonlayer}
\end{tikzpicture}
}
    \end{equation}
    \smallskip  
    
    \nid This makes clear how the functor $(-)^{\, \tl}$ acts on morphisms. The simple idea is that it tracks the induced mapping on the 1-dimensional open intervals, as follows:
    \begin{equation}
         \tikzfig{figures/Ch2/INT14}{\begin{tikzpicture}
	\begin{pgfonlayer}{nodelayer}
		\node [style=0 strat] (0) at (0, 1) {};
		\node [style=0 strat] (1) at (3, 1) {};
		\node [style=0 strat] (2) at (-1.5, -2) {};
		\node [style=0 strat] (3) at (1.5, -2) {};
		\node [style=0 strat] (4) at (4.5, -2) {};
		\node [style=none] (5) at (0, 1.75) {$0$};
		\node [style=none] (6) at (3, 1.75) {$1$};
		\node [style=none] (7) at (-1.5, -2.75) {$0$};
		\node [style=none] (8) at (1.5, -2.75) {$1$};
		\node [style=none] (9) at (4.5, -2.75) {$2$};
		\node [style=none] (10) at (-6, 1) {$[1]^{\, \tl}$};
		\node [style=none] (11) at (-6, -2) {$[2]^{\, \tl}$};
		\node [style=none] (12) at (-6, -1.25) {};
		\node [style=none] (13) at (-6, 0.25) {};
		\node [style=none] (14) at (-6.75, -0.5) {$d_{1}^{\, \tl} $};
		\node [style=none] (15) at (-3, 1) {};
		\node [style=none] (16) at (6, 1) {};
		\node [style=none] (17) at (-3, -2) {};
		\node [style=none] (18) at (6, -2) {};
		\node [style=none] (19) at (-1.75, 1) {};
		\node [style=none] (20) at (1.5, 1) {};
		\node [style=none] (21) at (4.75, 1) {};
		\node [style=none] (22) at (-2.5, -2) {};
		\node [style=0 strat red] (23) at (-1.75, 1) {};
		\node [style=0 strat red] (24) at (1.5, 1) {};
		\node [style=0 strat red] (25) at (4.75, 1) {};
		\node [style=0 strat red] (26) at (-2.25, -2) {};
		\node [style=0 strat red] (27) at (0, -2) {};
		\node [style=0 strat red] (28) at (3, -2) {};
		\node [style=0 strat red] (29) at (5.25, -2) {};
	\end{pgfonlayer}
	\begin{pgfonlayer}{edgelayer}
		\draw [style=cd-ar, in=-90, out=90] (12.center) to (13.center);
		\draw (2) to (0);
		\draw (4) to (1);
		\draw (15.center) to (16.center);
		\draw (17.center) to (18.center);
		\draw [style=cd-ar red] (26) to (23);
		\draw [style=cd-ar red] (27) to (24);
		\draw [style=cd-ar red] (28) to (24);
		\draw [style=cd-ar red] (29) to (25);
	\end{pgfonlayer}
\end{tikzpicture}
} \qquad \qquad \tikzfig{figures/Ch2/INT15}{\begin{tikzpicture}
	\begin{pgfonlayer}{nodelayer}
		\node [style=0 strat] (0) at (-1.5, -2) {};
		\node [style=0 strat] (1) at (1.5, -2) {};
		\node [style=0 strat] (2) at (-3, 1) {};
		\node [style=0 strat] (3) at (0, 1) {};
		\node [style=0 strat] (4) at (3, 1) {};
		\node [style=none] (5) at (-1.5, -2.75) {$0$};
		\node [style=none] (6) at (1.5, -2.75) {$1$};
		\node [style=none] (7) at (-3, 1.75) {$0$};
		\node [style=none] (8) at (0, 1.75) {$1$};
		\node [style=none] (9) at (3, 1.75) {$2$};
		\node [style=none] (10) at (-6, 1) {$[2]^{\, \tl}$};
		\node [style=none] (11) at (-6, -2) {$[1]^{\, \tl}$};
		\node [style=none] (12) at (-6, -1.25) {};
		\node [style=none] (13) at (-6, 0.25) {};
		\node [style=none] (14) at (-6.75, -0.5) {$s_{0}^{\, \tl}$};
		\node [style=none] (15) at (-4, 1) {};
		\node [style=none] (16) at (4, 1) {};
		\node [style=none] (17) at (-4, -2) {};
		\node [style=none] (18) at (4, -2) {};
		\node [style=0 strat red] (19) at (-3.5, 1) {};
		\node [style=0 strat red] (20) at (-1.5, 1) {};
		\node [style=0 strat red] (21) at (1.5, 1) {};
		\node [style=0 strat red] (22) at (3.5, 1) {};
		\node [style=0 strat red] (23) at (-3, -2) {};
		\node [style=0 strat red] (24) at (0.25, -2) {};
		\node [style=0 strat red] (25) at (3, -2) {};
		\node [style=0 strat red] (26) at (-3.5, 1) {};
	\end{pgfonlayer}
	\begin{pgfonlayer}{edgelayer}
		\draw [style=cd-ar] (12.center) to (13.center);
		\draw (2) to (0);
		\draw (3) to (0);
		\draw (4) to (1);
		\draw (15.center) to (16.center);
		\draw (17.center) to (18.center);
		\draw [style=cd-ar red] (23) to (26);
		\draw [style=cd-ar red] (24) to (21);
		\draw [style=cd-ar red] (25) to (22);
	\end{pgfonlayer}
\end{tikzpicture}
}
    \end{equation} 
\end{remark}












\subsection{Interval spaces}

When working with complete Segal spaces as a model for $(\oo,1)$-categories, all the combinatorics is encoded in the simplex category $\bDelta$. In this chapter we instead use the interval category $\Int$ as our indexing category. One should think of $\Int$ as keeping track of \emph{where} composition of morphisms is allowed to happen (via unions of open subintervals of $[0,1]$), while the functor
\begin{equation}
  \gamma \colon \Int \longra \bDelta^{\op}
\end{equation}
forgets this extra geometric information and remembers only the number and order of “slots” for composition (the number of inner connected components).

\begin{definition}
A \define{pre-interval space} is a functor
\begin{equation}
  \cF \colon \Int \longra \sset.
\end{equation}
\end{definition}

The fibres of $\gamma$ consist of intervals which have the same underlying combinatorics in $\bDelta^{\op}$ but differ by refining or sliding the closed pieces in their complements. We single out the class of maps
\begin{equation}
  \cW_{\gamma} \subset \Fun(\Int,\sset)
\end{equation}
consisting of morphisms of co-representables
\begin{equation}
  \Hom_{\Int}(J,-) \xra{\,f_*\,} \Hom_{\Int}(I,-)
\end{equation}
induced by morphisms $f\colon I \to J$ in $\Int$ such that $\gamma(f)=\id_{[n]}$ for some $n$. Equivalently, these are the maps coming from morphisms which induce an isomorphism on path-components. A pre-interval space which sends all such $f$ to weak equivalences is, heuristically, “constant along the
fibres of $\gamma$”.

\begin{definition}
A pre-interval space $\cF\colon \Int \to \sset$ is called an \define{interval space} if for every morphism $f\colon I \to J$ in $\Int$ with $\gamma(f)=\id$ the induced map
\begin{equation}
  \cF(I) \xra{\ \cF(f)\ } \cF(J)
\end{equation}
is a weak equivalence of simplicial sets.
\end{definition}

Thus, an interval space is a pre-interval space whose homotopy type only depends on the image of $I$ in $\bDelta^{\op}$. In §\ref{ch2:: subsec: interval_spaces_model_infinity_cats} we will make this precise by
left Bousfield localizing $\Fun(\Int,\sset)$ along $\cW_{\gamma}$ and showing that the resulting model category is Quillen equivalent to Rezk's model structure on simplicial spaces.

The next definition is the analogue of the Segal condition for simplicial spaces, but phrased geometrically in terms of cutting an interval into smaller pieces.

\begin{definition}\label{ch2:: def: interval Segal condition}
Let $\cF\colon \Int \to \sset$ be a pre-interval space. We say that $\cF$ satisfies the \define{interval Segal condition} if for every $n \geq 2$ and every $\sU(A_1,\ldots,A_n)\in\Int$ the induced
commuting square
\begin{equation}
    \begin{tikzcd}[ampersand replacement=\&,column sep=2.25em]
        {\cF(\sU(A_1,\ldots, A_n))} \&\& {\cF(\sU(A_1,\ldots, A_k))} \\
        \\
        {\cF(\sU(A_k,\ldots, A_n))} \&\& {\cF(\sU(A_k))}
        \arrow["{(\alpha_{[1,k]})_*}", from=1-1, to=1-3]
        \arrow["{(\alpha_{[k,n]})_*}"', from=1-1, to=3-1]
        \arrow["{(\alpha_{[k]})_*}", from=1-3, to=3-3]
        \arrow["{(\alpha_{[k]})_*}"', from=3-1, to=3-3]
    \end{tikzcd}
\end{equation}
is a homotopy pullback square in $\sset$, for every $1 \le k \le n-1$.
\end{definition}

Informally, one should think of $\sU(A_1,\ldots,A_n)$ as encoding a string of $n$ composable 1-morphisms, with the closed intervals $A_i$ marking where the topology is allowed to change. The interval Segal condition says that gluing along any intermediate cut $A_k$ recovers $\cF(\sU(A_1,\ldots,A_n))$
as the appropriate homotopy pullback of the two shorter strings.

\begin{definition}[Complete interval Segal space]
An \define{interval Segal space} is a pre-interval space which is both an interval space and satisfies the interval Segal condition. We will say that an interval Segal space $\cF$ is \define{complete} if its associated
Segal space
\begin{equation}
  \Rrkan_{\gamma}(\cF)\colon \bDelta^{\op} \longra \sset
\end{equation}
is a complete Segal space in the sense of Rezk.   
\end{definition}

We will justify this definition and construct a model structure whose fibrant objects are complete interval Segal spaces in §\ref{ch2:: subsec: interval_spaces_model_infinity_cats}, where we prove that this model structure is Quillen equivalent to the Rezk model structure on complete Segal spaces.


\section{Interval Segal spaces as a model for \texorpdfstring{$(\oo,1)$}{(infty,1)}-categories}
\label{ch2:: subsec: interval_spaces_model_infinity_cats}

In this subsection we construct a model structure on the category
$\fun(\Int,\sset)$ of pre-interval spaces and prove that it is Quillen
equivalent to the Rezk model structure \cite{Rezk01} on complete Segal spaces
described in \Cref{ch2:: prop: Rezk-model-structure}. In particular, we obtain
a left proper, combinatorial model structure on the category of pre-interval
spaces whose fibrant objects are complete interval Segal spaces and whose weak
equivalences are the Dwyer--Kan equivalences.

\medskip

\nid The structure of this subsection can be understood through progressively more refined model structures on the category of pre-interval spaces. We begin by localizing pre-interval spaces along $\pi_0$-isomorphisms, which yields a model category Quillen equivalent to simplicial spaces; we then further localize along the derived image of the spine inclusions to recover a Segal condition for interval spaces.

\subsection{Equivalence with simplicial spaces}

\begin{definition}\label{ch2:: def: pi_0 equiv}
    Let $\cW_{\pi_0}$ be the set of morphisms in $\fun(\Int,\sset)$ which are the image of $\pi_0$-isomorphisms under the co-Yoneda embedding 
    \begin{equation}
        \Int^{\op} \hookrightarrow \fun(\Int,\sset).
    \end{equation}
\end{definition}

\medskip

Recall the discussion in
\Cref{ch2:: subsubsec: homotopy-limits-bousfield} about left Bousfield
localizations. We are interested in studying pre-interval spaces
$\cF\colon \Int \to \sset$ which send every $\pi_0$-isomorphism $f\colon I \to
J$ to a weak equivalence $\cF(I) \xra{\simeq} \cF(J)$ in $\sset_{\sf QK}$, and
we would like to design a model structure for which these are precisely the
fibrant objects.

\begin{observation}
    A morphism in $\cW_{\pi_0}$ is a map of co-representables
    \begin{equation}
        \Hom_{\Int}(J,-) \xra{\,f_*\,} \Hom_{\Int}(I,-)
    \end{equation}
    such that 
    \[
      \gamma(I) \xra{\gamma(f)} \gamma(J)
    \]
    is an isomorphism, i.e., the identity since $\bDelta^\op$ is gaunt.
\end{observation}

\begin{proposition}\label{ch2:: prop: model structure interval}
    There exists a model structure on the category of pre-interval spaces $\fun(\Int,\sset)$, which we will denote $\fun(\Int,\sset)_{\mathsf{inj},h}$, with the following properties:
    \begin{enumerate}
        \item The weak equivalences are the $\cW_{\pi_0}$-local equivalences.
        \item The cofibrations are the injective cofibrations.
        \item The fibrant objects are the $\cW_{\pi_0}$-local objects, i.e.\ injectively fibrant functors $\cF$ such that $\cF(I) \xra{\cF(f)} \cF(J)$ is a level-wise weak equivalence of simplicial sets for every $f\colon I \to J$ in $\cW_{\pi_0}$.
        \item The model structure is left proper, combinatorial, and simplicial.
    \end{enumerate} 
\end{proposition}

\begin{proof}
We construct $\fun(\Int,\sset)_{\mathsf{inj},h}$ as the left Bousfield
localization of $\fun(\Int,\sset)_{\mathsf{inj}}$ with respect to the set
$\cW_{\pi_0}$ of inclusions inducing an isomorphism on $\pi_0$.

To see that this localization exists, note first that
$\sset_{\mathsf{QK}}$ is a cofibrantly generated, left proper, simplicial
model category. Since $\Int$ is small, the general existence theorem for
injective model structures on functor categories (see, for instance,
\Cref{ch2:: thm: functor model structures}) implies that
$\fun(\Int,\sset)_{\mathsf{inj}}$ exists and is again cofibrantly generated,
left proper, combinatorial, and simplicial. In particular, $\fun(\Int,\sset)$
is locally presentable. The set $\cW_{\pi_0}$ consists of maps between
co-representables, so it is indeed a \emph{set} of morphisms in
$\fun(\Int,\sset)$. We may therefore apply the general existence theorem for
left Bousfield localizations (see \Cref{ch2:: thm: left Bousfield localization})
to obtain the localized model structure
$\fun(\Int,\sset)_{\mathsf{inj},h}$ with the claimed weak equivalences. This
establishes (1) and (4).

Assertion (2) is straightforward from the properties of left Bousfield
localization: the class of cofibrations is unchanged.

For (3), recall that $\cW_{\pi_0}$-local objects are precisely those
injectively fibrant pre-interval spaces $\cF$ such that, for every morphism
$\eta\colon \cG\to\cH$ in $\cW_{\pi_0}$, the induced map
\[
  \Hom_{\fun(\Int,\sset)}(\cH,\cF)
  \xlongrightarrow{\ \eta^*\ }
  \Hom_{\fun(\Int,\sset)}(\cG,\cF)
\]
is a weak equivalence of simplicial sets. By definition of $\cW_{\pi_0}$, any
$\eta$ is of the form
\[
  \Hom_{\Int}(J,-)\xrightarrow{\,f_*\,}\Hom_{\Int}(I,-)
\]
for some $f\colon I\to J$ in $\Int$ inducing an isomorphism on $\pi_0$. In
this case the Yoneda lemma identifies the induced map above with
\[
  \cF(J)\xrightarrow{\;\cF(f)\;}\cF(I).
\]
Thus $\cF$ is $\cW_{\pi_0}$-local if and only if it is injectively fibrant and
$\cF(f)$ is a weak equivalence for every $f\in\cW_{\pi_0}$, which is exactly
the description in (3).
\end{proof}

Having defined a model structure on $\fun(\Int, \sset)$ with the fibrant objects having the homotopical properties we originally desired, we now construct a Quillen pair between the injective model structure on simplicial spaces $\fun(\bDelta^{\op},\sset)_{\sf inj}$ and $\fun(\Int, \sset)_{\mathsf{inj},h}$. We will use the injective model structure to obtain a naive Quillen pair from which we will build the desired one between $\fun(\bDelta^{\op},\sset)_{\mathsf{inj}}$ and $\fun(\Int,\sset)_{\mathsf{inj},h}$.

\begin{proposition}\label{ch2:: prop: Quillen_pair_interval_simplicial_naive}
There is a Quillen pair
\begin{equation}
    \begin{tikzcd}[ampersand replacement=\&]
	{\fun(\bDelta^{\op},\sset)_{\mathsf{inj}}} \&\& {\fun(\Int,\sset)_{\mathsf{inj}}}
	\arrow[""{name=0, anchor=center, inner sep=0}, "{\gamma^*}", bend left=10, from=1-1, to=1-3]
	\arrow[""{name=1, anchor=center, inner sep=0}, "{\rkan_{\gamma}}", bend left=10, from=1-3, to=1-1]
	\arrow["\dashv"{anchor=center, rotate=-90}, draw=none, from=1, to=0]
    \end{tikzcd}
\end{equation}
induced by the functor $\gamma\colon \Int \to \bDelta^{\op}$. 
\end{proposition}

\begin{proof}
Precomposition with $\gamma$ preserves levelwise cofibrations and levelwise
weak equivalences. Since these are precisely the cofibrations and weak
equivalences in the injective model structures, $\gamma^*$ is a left Quillen
functor. Its right adjoint is the right Kan extension $\rkan_\gamma$.
\end{proof}

We want to show that this Quillen pair extends to a Quillen pair between
$\fun(\bDelta^{\op},\sset)_{\mathsf{inj}}$ and
$\fun(\Int,\sset)_{\mathsf{inj},h}$. Using
\Cref{ch2:: obs: interval_category_small} together with the general existence
results for injective model structures on functor categories (again
\Cref{ch2:: thm: functor model structures}), we know that
$\fun(\Int,\sset)_{\mathsf{inj}}$ is left proper and combinatorial, so its left
Bousfield localization at $\cW_{\pi_0}$ exists.

\begin{proposition}
    There is a Quillen pair 
    \begin{equation}
    \begin{tikzcd}[ampersand replacement=\&]
	{\fun(\bDelta^{\op},\sset)_{\mathsf{inj}}} \&\& {\fun(\Int,\sset)_{\mathsf{inj}}[\cW_{\pi_0}^{-1}]}
	\arrow[""{name=0, anchor=center, inner sep=0}, "{\gamma^*}", bend left=10, from=1-1, to=1-3]
	\arrow[""{name=1, anchor=center, inner sep=0}, "{\rkan_{\gamma}}", bend left=10, from=1-3, to=1-1]
	\arrow["\dashv"{anchor=center, rotate=-90}, draw=none, from=1, to=0]
    \end{tikzcd}
    \end{equation}
\end{proposition}

\begin{proof}
    Building on the Quillen pair from \Cref{ch2:: prop: Quillen_pair_interval_simplicial_naive}, this is a formal consequence of left Bousfield localization: the localization leaves the class of cofibrations unchanged and enlarges the class of weak equivalences. In particular, $\gamma^*$ continues to send level-wise cofibrations to injective cofibrations, and level-wise weak equivalences to level-wise weak equivalences, hence in particular to weak equivalences in the localized model structure. Therefore the adjunction descends to the localized model structure.
\end{proof}

Having shown that the adjoint pair $\gamma^{*} \dashv \rkan_{\gamma}$ extends to the localized model structure, we now proceed to show that it is in fact a Quillen equivalence. Before we delve into the details of the proof, it is helpful to offer a brief blueprint of the argument:
\begin{enumerate}
    \item[\textbf{\color{brown}{Step \ding{192}}}] We first check that if $\cF$ is fibrant, then the right Kan extension agrees with the homotopy right Kan extension:
    \begin{equation}
        \rkan_{\gamma}(\cF) \xra{\,\simeq\,} \Rrkan_{\gamma}(\cF).
    \end{equation}
    \item[\textbf{\color{brown}{Step \ding{193}}}] The homotopy right Kan extension is computed pointwise by 
    \begin{equation}
        \Rrkan_{\gamma}(\cF)([n]) \,\simeq\, \underset{I \in ([n]/ \Int)}{\mathsf{holim}} \, \cF(I).
    \end{equation}
    \item[\textbf{\color{brown}{Step \ding{194}}}] The full subcategory $\cJ_n$ of $([n]/ \Int)$ on the objects $\big([n] \xra{\id} \gamma(I)\big)$ is shown to be homotopy initial (see \Cref{ch2:: prop: hom initial} below).
    \item[\textbf{\color{brown}{Step \ding{195}}}] The previous observations imply that we can compute the homotopy right Kan extension as a homotopy limit over the category $\cJ_n$.
    \item[\textbf{\color{brown}{Step \ding{196}}}] Since $\cF$ is $\cW_{\pi_0}$-local, every morphism in $\cJ_n$ is mapped to a weak equivalence. Together with the fact that the nerve of $\cJ_n$ is weakly contractible, this implies that $\Rrkan_{\gamma}(\cF)([n]) \to \cF(I)$ is a weak equivalence for every $I\in \Int$ (cf.\ \Cref{ch2:: prop: homotopy-limit-homotopically-constant}).
\end{enumerate}

\subsubsection{Digression on configuration spaces}

Before we move on to proving that $\gamma^* \dashv \rkan_\gamma$ is a Quillen equivalence, let us first provide a brief recollection on configuration spaces of $\RR^{n}$. Let $r \geq 0$ be an integer, and let $\underline{r} = \{1,\ldots,r\}$. We define the space
\begin{equation}
    \mathsf{Conf}_{r}(\RR^{n}) \,\overset{\sf def}{=}\, \big\{\underline{r} \xlonghookra{\phi} \RR^{n}\, : \, \phi \text{ is an injection} \big\}
\end{equation}
of injections from the finite set $\underline{r}$ into $\RR^{n}$, topologized by removing all the closed diagonal subspaces 
\begin{equation}
    (\RR^{n})^{\underline{r}} \setminus \Big(\coprod_{0\leq i\neq j\leq r} \Delta_{i,j}\Big).
\end{equation}
We call this the \emph{ordered configuration space} of $r$ particles in $\RR^{n}$. The ordered configuration space of $r$ points in $\RR^{n}$ admits an action by $\Sigma_r$, the symmetric group on $r$ letters, defined by
\begin{align}
 \Sigma_r \times \mathsf{Conf}_r(\RR^{n}) &\longra \mathsf{Conf}_r(\RR^{n}) \\
 \big(\sigma,(x_1,\ldots,x_r) \big) &\longmapsto (x_{\sigma(1)},\ldots, x_{\sigma(r)}).
\end{align}

\begin{lemma}
    The action of the symmetric group $\Sigma_r \curvearrowright \mathsf{Conf}_r(\RR^{n})$ is free and properly discontinuous. As a result, the quotient map 
    \begin{equation}
        q \colon \mathsf{Conf}_{r}(\RR^{n}) \longra \mathsf{UConf}_{r}(\RR^{n})
    \end{equation} 
    is a regular covering map (of degree $r!$).
\end{lemma}

\begin{definition}
    We call the space $\mathsf{UConf}_{r}(\RR^{n})$ the \emph{unordered configuration space of $r$ points in $\RR^{n}$}.
\end{definition}

In general, configuration spaces of $\RR^{n}$ for arbitrary $n$ can have a complicated and interesting homotopy type, but not when $n=1$.

\begin{lemma}
    The space $\mathsf{UConf}_{r}(\RR) \cong \RR^{r}$. In particular, the unordered configuration space of $r$ points in $\RR$ is contractible.
\end{lemma}

The point of the preceding digression is that in dimension one the unordered
configuration spaces we will encounter are \emph{contractible}: they are open,
convex subsets of a Euclidean space. We now explain how these configuration
spaces appear naturally in the proof \textbf{\color{brown}{Step \ding{194}}}, stating that 
\begin{equation}
  \iota_n \maps \cJ_n \xlonghookra{\,\,\,\,\,\,\,} ([n]/\Int)
\end{equation}
is homotopy initial, which is the key combinatorial input in our computation of
the homotopy right Kan extension along $\gamma$.

Roughly speaking, fixing an object
\begin{equation}
  \alpha \maps [n] \longra \gamma(J)
\end{equation}

in the slice category $([n]/\Int)$, we will associate to each object of the
comma category $(\alpha / \cJ_n)$ an open region in a suitable configuration
space of points in $J^c \subset (0,1)$. The Lurie--Seifert--van Kampen theorem
then allows us to identify the nerve of $(\alpha / \cJ_n)$ with this
configuration space. Since the latter is built out of unordered configuration
spaces of open intervals in~$\RR$, the contractibility results above imply that
each $(\alpha / \cJ_n)$ has contractible nerve, hence that $\iota_n$ is
homotopy initial.

\begin{definition}
Let $\cJ_n$ denote the full subcategory of $([n]/\Int)$ on the identity arrows $\big([n] \xra{\id} \gamma(I)\big)\in \bDelta^{\op}$, for $I\in \Int$.    
\end{definition}

\begin{observation}
    The objects of $\cJ_n$ are parameterized by the intervals $I\in \Int$ with $n$ interior connected components. Moreover, given a fixed pair of objects 
    \[
      [n] \xra{\id} \gamma(I)
      \quad\text{and}\quad
      [n] \xra{\id} \gamma(I')
    \]
    in $\cJ_n$, the set of morphisms 
    \[
        \Hom_{\cJ_n}\big([n] \xra{\id} \gamma(I), [n] \xra{\id} \gamma(I')\big)
    \]
    is naturally identified with
    \[
        \{\,f\in\Hom_{\Int}(I,I') \mid \gamma(f) = \id_{[n]} \,\},
    \]
    i.e.\ with the inclusions $I \hookrightarrow I'$ which induce isomorphisms on $\pi_0$.
\end{observation}

The next two results are key propositions. As mentioned in the outline of the proof above, they ensure that the structure maps of the homotopy limit are weak equivalences — the proof presented here is a generalization of a similar kind in the work of Karlsson--Scheimbauer--Walde \cite{KSW25}. 

\begin{proposition}\label{ch2:: prop: hom initial}
    The inclusion functor 
    \begin{equation} 
        \iota_n\maps \cJ_n \xlonghookra{\,\,\,\,\,} [n]/\Int 
    \end{equation} 
    is homotopy initial. 
\end{proposition} 
\begin{proof}
     It is useful to start by outlining the strategy of the proof, that will greatly simplify exposition of its technical part. The first disclaimer is that we are going to show its opposite: we will prove the theorem using the $\cInt$ perspective. What we mean by that, is that we will instead prove the functor 
     \begin{equation} 
        \iota_{n}^{\op} \maps \cJ_{n}^{\op} \xlonghookra{\,\,\,\,} ([n]/\Int)^{\op} = \cInt/[n] 
    \end{equation} is homotopy final. For simplification, we will just call $\cJ_n^{\op} = \cJ_n$ for the remainder of the proof — assuring the reader that this is only local notation. 
    
    
    Secondly, the way we will prove that $\left|\alpha / \cJ_n\right|$ is contractible is through the usage of Lurie--Seifert-van-Kampen theorem \Cref{ap:: thm: lurie Seifert}. Roughly, proofs of contractibility using this theorem work as follows: one shows that the category $\cD$ is weakly equivalent to a space $C$, which is known to be contractible, by constructing a functor 
    \begin{equation} 
        f\maps \cD \longra \cU(C) 
    \end{equation} 
    from $\cD$ into the poset of open subsets of $C$ and satisfying two additional properties: 
    \begin{enumerate} 
        \item For every $d\in \cD$ the open subset $f(d) \subset C$ is contractible. 
        \item For every point $x\in C$, the subcategory $\cD_{x} \hookra \cD$ containing the objects $d\in \cD$ for which $x \in f(d)$ is weakly contractible — in particular it is non-empty. 
    \end{enumerate} 
    The contractibility of $\cD$ then follows from the following sequence of equivalences: 
    \begin{align*} 
        \left| \cD \right| \,\,\, &\overset{(1)}{\simeq} \,\,\, \hoColimit[d \in \cD] * \\ 
        &\overset{(2)}{\simeq} \,\,\, \hoColimit[d \in \cD] f(d) \\
        &\overset{(3)}{\simeq} \,\,\, C \\ 
        &\overset{(4)}{\simeq} \,\,\, *. 
    \end{align*} 
    Equivalence (1) follows from ..., (2) results from the contractibility of $f(d)$ and the homotopy invariance of homotopy colimits, (3) follows from the Lurie--Seifert-van-Kampen theorem, if assuming contractibility of $\cD_x$ for every $x$, and finally (4) follows from the contractibility of $C$. \smallskip 
    
    We begin by constructing the space $C$, in the outline above. Fix $\alpha \maps \pi_{0}(J) \ra [n]$ a morphism in $\bDelta$, i.e., an object of $\cInt/[n]$. We define: 
    \begin{equation} 
        \mathsf{Conf}(J;\alpha)\,\, \overset{\sf def}{=} \,\, \left\{ 
            \begin{tikzcd}[ampersand replacement=\&] 
                {[n]^{\, \tl} } \& {(0,1)} \\ \& {J^{c}} \arrow[hook, from=1-1, to=1-2] \arrow["p"', curve={height=6pt}, dashed, from=1-1, to=2-2] \arrow[hook, from=2-2, to=1-2] \end{tikzcd} \text{ s.t. } p \text{ is injective and } \begin{tikzcd}[ampersand replacement=\&] {\pi_{0}([n]^{\, \tl})} \& {\pi_0(J^{c})} \\ { [n]^{\, \tl}} \& {\gamma(J^{c})^{\, \tl} } 
                    \arrow["{\pi_{0}(p)}", from=1-1, to=1-2] 
                    \arrow[equals, from=1-1, to=2-1] 
                    \arrow[equals, from=1-2, to=2-2] 
                    \arrow["{(\alpha^{\op})^{\, \tl}}"', from=2-1, to=2-2] 
                \end{tikzcd} \right\} \end{equation} This is the ordered space of $(n+2) = \left| [n]^{\, \tl} \right|$ points in $(0,1)$ satisfying two additional conditions: 
                
                \begin{enumerate} 
                    \item The injection factors through the disjoint union of open intervals $J^c$. 
                    \item Upon passage to $\pi_0$, agrees with $(\alpha^\op)^{\, \tl}$ in $\sO_{\pm}$. 
                \end{enumerate} 
                
                \nid We let $\mathsf{UConf}(J;\alpha)$ be its space of $\Sigma_{n+2}$-coinvariants, i.e., the unordered analogue. This is an open and convex subspace of $\mathsf{UConf}_{n+2}((0,1))$; therefore contractible. Before continuing with the proof, it is useful to indicate an example of a point in this space. Let $J$ be the following object of $\cInt$ 
                \smallskip 
                \begin{equation} 
                    \tikzfig{figures/Ch2/INT16}{\begin{tikzpicture}
	\begin{pgfonlayer}{nodelayer}
		\node [style=none] (6) at (-7, 0) {};
		\node [style=none] (7) at (7, 0) {};
		\node [style=none] (8) at (-5, 0) {};
		\node [style=none] (9) at (-3, 0) {};
		\node [style=none] (10) at (2, 0) {};
		\node [style=none] (11) at (4, 0) {};
		\node [style=none] (12) at (-5, 0.2) {};
		\node [style=none] (13) at (-5, -0.2) {};
		\node [style=none] (14) at (-3, 0.2) {};
		\node [style=none] (15) at (-3, -0.2) {};
		\node [style=none] (16) at (2, 0.2) {};
		\node [style=none] (17) at (2, -0.2) {};
		\node [style=none] (18) at (4, 0.2) {};
		\node [style=none] (19) at (4, -0.2) {};
	\end{pgfonlayer}
	\begin{pgfonlayer}{edgelayer}
		\draw (6.center) to (7.center);
		\draw [style={open_int_blue}] (8.center) to (9.center);
		\draw [style={open_int_blue}] (10.center) to (11.center);
		\draw [style={open_int_blue}] (12.center) to (13.center);
		\draw [style={open_int_blue}] (14.center) to (15.center);
		\draw [style={open_int_blue}] (16.center) to (17.center);
		\draw [style={open_int_blue}] (18.center) to (19.center);
	\end{pgfonlayer}
\end{tikzpicture}
}, 
                \end{equation} 
                \smallskip 
                
                \nid and let $\alpha \maps \pi_{0}(J) \simeq [1] \xra{\lag 0,2 \rag} [2]$. A point $p \maps [2]^{\, \tl} \ra J^{c}$ in $\mathsf{UConf}(J;\alpha)$ is thus an unordered configuration of $4$ points in $J^{c} = \sqcup_{i} U_i$ — which is the disjoint union of three ordered open intervals — satisfying an additional condition: 
                \begin{enumerate} 
                    \item[$(*)$] The number of points in each connected open interval is determined by the cardinality of $\pi_{0}(p)^{-1}([U_i])$. 
                \end{enumerate} 
                Which in this case is 1 for $U_1$, 2 for $U_2$, and $1$ for $U_3$. Therefore the configuration on the left is a point in $\mathsf{UConf}(J;\alpha)$ and the one on the right is not 
                \begin{equation} 
                    \tikzfig{figures/Ch2/INT18}{\begin{tikzpicture}
	\begin{pgfonlayer}{nodelayer}
		\node [style=none] (6) at (-7, 0) {};
		\node [style=none] (7) at (7, 0) {};
		\node [style=none] (8) at (-5, 0) {};
		\node [style=none] (9) at (-3, 0) {};
		\node [style=none] (10) at (2, 0) {};
		\node [style=none] (11) at (4, 0) {};
		\node [style=none] (12) at (-5, 0.2) {};
		\node [style=none] (13) at (-5, -0.2) {};
		\node [style=none] (14) at (-3, 0.2) {};
		\node [style=none] (15) at (-3, -0.2) {};
		\node [style=none] (16) at (2, 0.2) {};
		\node [style=none] (17) at (2, -0.2) {};
		\node [style=none] (18) at (4, 0.2) {};
		\node [style=none] (19) at (4, -0.2) {};
		\node [style=none] (61) at (-8.5, 0) {$J$};
		\node [style=none] (72) at (-6, -0.5) {$U_1$};
		\node [style=none] (73) at (0, -0.5) {$U_2$};
		\node [style=none] (74) at (5.5, -0.5) {$U_3$};
		\node [style=0 strat] (76) at (-6.5, 0) {};
		\node [style=0 strat] (77) at (-2.25, 0) {};
		\node [style=0 strat] (78) at (1.25, 0) {};
		\node [style=0 strat] (79) at (4.5, 0) {};
	\end{pgfonlayer}
	\begin{pgfonlayer}{edgelayer}
		\draw (6.center) to (7.center);
		\draw [style={open_int_blue}] (8.center) to (9.center);
		\draw [style={open_int_blue}] (10.center) to (11.center);
		\draw [style={open_int_blue}] (12.center) to (13.center);
		\draw [style={open_int_blue}] (14.center) to (15.center);
		\draw [style={open_int_blue}] (16.center) to (17.center);
		\draw [style={open_int_blue}] (18.center) to (19.center);
	\end{pgfonlayer}
\end{tikzpicture}
}, \quad \quad \tikzfig{figures/Ch2/INT19}{\begin{tikzpicture}
	\begin{pgfonlayer}{nodelayer}
		\node [style=none] (6) at (-7, 0) {};
		\node [style=none] (7) at (7, 0) {};
		\node [style=none] (8) at (-5, 0) {};
		\node [style=none] (9) at (-3, 0) {};
		\node [style=none] (10) at (2, 0) {};
		\node [style=none] (11) at (4, 0) {};
		\node [style=none] (12) at (-5, 0.2) {};
		\node [style=none] (13) at (-5, -0.2) {};
		\node [style=none] (14) at (-3, 0.2) {};
		\node [style=none] (15) at (-3, -0.2) {};
		\node [style=none] (16) at (2, 0.2) {};
		\node [style=none] (17) at (2, -0.2) {};
		\node [style=none] (18) at (4, 0.2) {};
		\node [style=none] (19) at (4, -0.2) {};
		\node [style=none] (61) at (-8.5, 0) {$J$};
		\node [style=none] (72) at (-6, -0.5) {$U_1$};
		\node [style=none] (73) at (0, -0.5) {$U_2$};
		\node [style=none] (74) at (5.5, -0.5) {$U_3$};
		\node [style=0 strat] (76) at (-6.5, 0) {};
		\node [style=0 strat] (77) at (-2.25, 0) {};
		\node [style=0 strat] (78) at (5.25, 0) {};
		\node [style=0 strat] (79) at (4.5, 0) {};
	\end{pgfonlayer}
	\begin{pgfonlayer}{edgelayer}
		\draw (6.center) to (7.center);
		\draw [style={open_int_blue}] (8.center) to (9.center);
		\draw [style={open_int_blue}] (10.center) to (11.center);
		\draw [style={open_int_blue}] (12.center) to (13.center);
		\draw [style={open_int_blue}] (14.center) to (15.center);
		\draw [style={open_int_blue}] (16.center) to (17.center);
		\draw [style={open_int_blue}] (18.center) to (19.center);
	\end{pgfonlayer}
\end{tikzpicture}
} \, . 
                \end{equation} 
                \smallskip 
                
                \nid We indicate the combinatorics at play with the following diagram 
                \begin{equation} 
                    \tikzfig{figures/Ch2/INT17}{\begin{tikzpicture}
	\begin{pgfonlayer}{nodelayer}
		\node [style=none] (6) at (-7, -3) {};
		\node [style=none] (7) at (7, -3) {};
		\node [style=none] (8) at (-5, -3) {};
		\node [style=none] (9) at (-3, -3) {};
		\node [style=none] (10) at (2, -3) {};
		\node [style=none] (11) at (4, -3) {};
		\node [style=none] (12) at (-5, -2.8) {};
		\node [style=none] (13) at (-5, -3.2) {};
		\node [style=none] (14) at (-3, -2.8) {};
		\node [style=none] (15) at (-3, -3.2) {};
		\node [style=none] (16) at (2, -2.8) {};
		\node [style=none] (17) at (2, -3.2) {};
		\node [style=none] (18) at (4, -2.8) {};
		\node [style=none] (19) at (4, -3.2) {};
		\node [style=none] (20) at (-7, 0) {};
		\node [style=none] (21) at (7, 0) {};
		\node [style=none] (22) at (-5, 0) {};
		\node [style=none] (23) at (-3, 0) {};
		\node [style=none] (24) at (2, 0) {};
		\node [style=none] (25) at (4, 0) {};
		\node [style=none] (26) at (-5, 0.2) {};
		\node [style=none] (27) at (-5, -0.2) {};
		\node [style=none] (28) at (-3, 0.2) {};
		\node [style=none] (29) at (-3, -0.2) {};
		\node [style=none] (30) at (2, 0.2) {};
		\node [style=none] (31) at (2, -0.2) {};
		\node [style=none] (32) at (4, 0.2) {};
		\node [style=none] (33) at (4, -0.2) {};
		\node [style=none] (34) at (0, -2.5) {};
		\node [style=none] (35) at (0, -0.5) {};
		\node [style=0 strat] (36) at (-4, 0) {};
		\node [style=0 strat] (37) at (3, 0) {};
		\node [style=none] (38) at (-7, 2.75) {};
		\node [style=none] (39) at (7, 2.75) {};
		\node [style=none] (40) at (-5, 2.75) {};
		\node [style=none] (41) at (-3, 2.75) {};
		\node [style=0 strat] (52) at (-4, 2.75) {};
		\node [style=0 strat] (53) at (3, 2.75) {};
		\node [style=0 strat] (54) at (0, 2.75) {};
		\node [style=none] (55) at (-4, -0.75) {$0$};
		\node [style=none] (56) at (3, -0.75) {$1$};
		\node [style=none] (57) at (-4, 3.5) {$0$};
		\node [style=none] (58) at (0, 3.5) {$1$};
		\node [style=none] (59) at (3, 3.5) {$2$};
		\node [style=none] (60) at (0.75, -1.75) {$\pi_0$};
		\node [style=none] (61) at (-8.5, -3) {$J$};
		\node [style=none] (62) at (-8.5, 0) {$\pi_{0}(J)$};
		\node [style=none] (63) at (-8.5, 2.75) { $[2]$};
		\node [style=none] (64) at (-6, 2.75) {};
		\node [style=0 strat red] (65) at (-6, 2.75) {};
		\node [style=0 strat red] (66) at (-1.5, 2.75) {};
		\node [style=0 strat red] (67) at (1.25, 2.75) {};
		\node [style=0 strat red] (68) at (5.5, 2.75) {};
		\node [style=0 strat red] (69) at (-6, 0) {};
		\node [style=0 strat red] (70) at (0, 0) {};
		\node [style=0 strat red] (71) at (5.5, 0) {};
		\node [style=none] (72) at (-6, -3.5) {$U_1$};
		\node [style=none] (73) at (0, -3.5) {$U_2$};
		\node [style=none] (74) at (5.5, -3.5) {$U_3$};
		\node [style=none] (75) at (8.5, -3) {$J^{c}$};
		\node [style=none] (76) at (8.5, 0) {$\gamma (J^{c})^{\, \tl}$};
		\node [style=none] (77) at (8.5, 2.75) { $[2]^{\, \tl}$};
	\end{pgfonlayer}
	\begin{pgfonlayer}{edgelayer}
		\draw (6.center) to (7.center);
		\draw [style={open_int_blue}] (8.center) to (9.center);
		\draw [style={open_int_blue}] (10.center) to (11.center);
		\draw [style={open_int_blue}] (12.center) to (13.center);
		\draw [style={open_int_blue}] (14.center) to (15.center);
		\draw [style={open_int_blue}] (16.center) to (17.center);
		\draw [style={open_int_blue}] (18.center) to (19.center);
		\draw (20.center) to (21.center);
		\draw [style={open_int_blue}] (22.center) to (23.center);
		\draw [style={open_int_blue}] (24.center) to (25.center);
		\draw [style={open_int_blue}] (26.center) to (27.center);
		\draw [style={open_int_blue}] (28.center) to (29.center);
		\draw [style={open_int_blue}] (30.center) to (31.center);
		\draw [style={open_int_blue}] (32.center) to (33.center);
		\draw [style=cd-ar] (34.center) to (35.center);
		\draw (38.center) to (39.center);
		\draw [style=cd-ar] (36) to (52);
		\draw [style=cd-ar] (37) to (53);
		\draw [style=cd-ar red] (65) to (69);
		\draw [style=cd-ar red] (66) to (70);
		\draw [style=cd-ar red] (67) to (70);
		\draw [style=cd-ar red] (68) to (71);
	\end{pgfonlayer}
\end{tikzpicture}
} 
                \end{equation} 
                \smallskip 
                
                \nid This example also clarifies how $\mathsf{UConf}(J;\alpha)$ is homeomorphic to a product of unordered configuration spaces of open intervals with fixed numbers of particles \begin{equation} \mathsf{UConf}_{k_1}(U_1) \times \mathsf{UConf}_{k_2}(I_2) \times \mathsf{UConf}_{k_3}(U_3), \end{equation} where $k_i = \pi_{0}(p)^{-1}([U_i])$, for $i \in \pi_{0}(J^{c}) = [- < 0 < 1]$. 
                \medskip 
                
                \nid Next, we want to define the subspaces $f(d)$ in the outline. In particular, for every object 

                \[\big(\pi_0(J) \ra \pi_0(I) \xra{\id} [n]\big) = \alpha,\]
                in $\alpha /\cJ_{n}$, we seek to construct an open subspace $\mathsf{UConf}(J;\alpha; I) \subset \mathsf{UConf}(J;\alpha)$. We do so as follows: 
                
                \begin{equation} 
                    \mathsf{UConf}(J;\alpha;I)\,\, \overset{\sf def}{=} \,\, \left\{ 
                        \begin{tikzcd}[ampersand replacement=\&]
                         {[n]^{\, \tl} } \& {(0,1)} \\ {I^{c}} \& {J^{c}} \arrow["\phi", hook, from=1-1, to=1-2] \arrow["p"', dashed, from=1-1, to=2-1] \arrow[hook, from=2-1, to=2-2] \arrow[hook, from=2-2, to=1-2] 
                        \end{tikzcd} 
                    \text{ s.t. } p\text{ is injective and } 
                    \begin{tikzcd}[ampersand replacement=\&] {\pi_{0}([n]^{\, \tl})} \& {\pi_0(I^{c})} \& {\pi_0(J^{c})} \\ { [n]^{\, \tl}} \& {\gamma(I^{c})^{\, \tl} } \& {\gamma(J^{c})^{\, \tl} } 
                    \arrow["{{\pi_{0}(p)}}", from=1-1, to=1-2] 
                    \arrow[equals, from=1-1, to=2-1] 
                    \arrow[from=1-2, to=1-3] 
                    \arrow[equals, from=1-2, to=2-2] 
                    \arrow[equals, from=1-3, to=2-3] 
                    \arrow[from=2-1, to=2-2] 
                    \arrow["{\alpha^\op}"', curve={height=18pt}, from=2-1, to=2-3] 
                    \arrow[from=2-2, to=2-3] 
                \end{tikzcd} \right\} 
                \end{equation} 
                \smallskip 
                
                \nid Note that for $(p\maps [2]^{\, \tl} \ra I^{c})$ to be an object in $\mathsf{UConf}(J;\alpha;I)$ we need only to ensure that $\pi_{0}(p)=\id$. In particular, this means that we can interpret points in this open subspace as configurations of $n+2$ points in $I^{c}$ which have a unique particle per connected component. We again depict an example of a point in this space (implicitly making a choice of $I\in \cInt$): 
                \begin{equation} 
                    \tikzfig{figures/Ch2/INT20}{\begin{tikzpicture}
	\begin{pgfonlayer}{nodelayer}
		\node [style=none] (6) at (-7, 0) {};
		\node [style=none] (7) at (7, 0) {};
		\node [style=none] (8) at (-5.5, 0) {};
		\node [style=none] (9) at (-2.5, 0) {};
		\node [style=none] (10) at (1.5, 0) {};
		\node [style=none] (11) at (4.75, 0) {};
		\node [style=none] (12) at (-5.5, 0.2) {};
		\node [style=none] (13) at (-5.5, -0.2) {};
		\node [style=none] (14) at (-2.5, 0.2) {};
		\node [style=none] (15) at (-2.5, -0.2) {};
		\node [style=none] (16) at (1.5, 0.2) {};
		\node [style=none] (17) at (1.5, -0.2) {};
		\node [style=none] (18) at (4.75, 0.2) {};
		\node [style=none] (19) at (4.75, -0.2) {};
		\node [style=none] (61) at (-8.5, 0) {$I$};
		\node [style=none] (80) at (-7, -2.5) {};
		\node [style=none] (81) at (7, -2.5) {};
		\node [style=none] (82) at (-5, -2.5) {};
		\node [style=none] (83) at (-3, -2.5) {};
		\node [style=none] (84) at (2, -2.5) {};
		\node [style=none] (85) at (4, -2.5) {};
		\node [style=none] (86) at (-5, -2.3) {};
		\node [style=none] (87) at (-5, -2.7) {};
		\node [style=none] (88) at (-3, -2.3) {};
		\node [style=none] (89) at (-3, -2.7) {};
		\node [style=none] (90) at (2, -2.3) {};
		\node [style=none] (91) at (2, -2.7) {};
		\node [style=none] (92) at (4, -2.3) {};
		\node [style=none] (93) at (4, -2.7) {};
		\node [style=none] (94) at (-8.5, -2.5) {$J$};
		\node [style=none] (95) at (-1.5, 0) {};
		\node [style=none] (96) at (0.25, 0) {};
		\node [style=none] (97) at (-1.5, 0) {};
		\node [style=none] (98) at (0.25, 0) {};
		\node [style=none] (99) at (-1.5, 0.2) {};
		\node [style=none] (100) at (-1.5, -0.2) {};
		\node [style=none] (101) at (0.25, 0.2) {};
		\node [style=none] (102) at (0.25, -0.2) {};
		\node [style=none] (103) at (-8.5, -1.75) {};
		\node [style=none] (104) at (-8.5, -0.75) {};
		\node [style=0 strat] (105) at (-6.5, 0) {};
		\node [style=0 strat] (106) at (-2, 0) {};
		\node [style=0 strat] (107) at (1.25, 0) {};
		\node [style=0 strat] (108) at (6.5, 0) {};
	\end{pgfonlayer}
	\begin{pgfonlayer}{edgelayer}
		\draw (6.center) to (7.center);
		\draw [style={open_int_blue}] (8.center) to (9.center);
		\draw [style={open_int_blue}] (10.center) to (11.center);
		\draw [style={open_int_blue}] (12.center) to (13.center);
		\draw [style={open_int_blue}] (14.center) to (15.center);
		\draw [style={open_int_blue}] (16.center) to (17.center);
		\draw [style={open_int_blue}] (18.center) to (19.center);
		\draw (80.center) to (81.center);
		\draw [style={open_int_blue}] (82.center) to (83.center);
		\draw [style={open_int_blue}] (84.center) to (85.center);
		\draw [style={open_int_blue}] (86.center) to (87.center);
		\draw [style={open_int_blue}] (88.center) to (89.center);
		\draw [style={open_int_blue}] (90.center) to (91.center);
		\draw [style={open_int_blue}] (92.center) to (93.center);
		\draw [style={open_int_blue}] (97.center) to (98.center);
		\draw [style={open_int_blue}] (99.center) to (100.center);
		\draw [style={open_int_blue}] (101.center) to (102.center);
		\draw [style=cd-ar] (103.center) to (104.center);
	\end{pgfonlayer}
\end{tikzpicture}
} 
                \end{equation} 
                \smallskip 
                
                \nid This subspace is, again, convex for similar reasons and therefore contractible. All that is left showing to finish the proof is that, given an arbitrary point in $\mathsf{UConf}(J;\alpha)$: 
                \[([n]^{\, \tl} \xra{p} J^{c} \hookra (0,1)) \in \mathsf{UConf}(J;\alpha).\] 
                the subposet 
                \begin{equation}
                     \left\{(I\subseteq J) \in \alpha/\cJ_n \, : (p \maps [n]^{\, \tl} \ra J^{c}) \in \mathsf{UConf}(J;\alpha; I) \right\} \,\, \subseteq \,\, \cJ_{n} 
                \end{equation}
                is contractible. To show that is non-empty we offer a simple procedure which produces an object $I\in \cInt$ such that $p\in \mathsf{UConf}(J;\alpha;I)$. For each connected open interval \begin{equation} 
                    J^{c} = U_1 \sqcup \cdots \sqcup U_{k} 
                \end{equation} 
                you do one of two things: if it contains no points, an impossibility for the extremal components, you close it by taking the hull of the two adjacent closed intervals; if it contains a single point, you do nothing; if it contains $m$ points $(p_i < \cdots < p_{i+m})$ with $m\geq 1$, you add $m-1$ closed connected intervals in $(p_{j},p_{j+1})$ for $j=i,\ldots,m+i-1$. We essentially used the quantity of particles on a given open edge of $J^c$, as a way of determining how many connected open intervals an object $I$ should have in that region, in order for it to receive an inclusion from $J$ in $\cInt$. 
                
                Finally, contractibility of this subposet comes from the fact that it is isomorphic to a product of posets, each given by the poset of connected closed intervals in $(p_i, p_{i+1})$ for $i=0,\ldots, n+2$ — which is filtered and therefore contractible. \end{proof}

\begin{proposition}\label{ch2:: prop: weak contractibility of pi_0}
    The nerve of the full subcategory $\iota_n \colon \cJ_n \hookrightarrow ([n]/\Int)$ is weakly contractible for every $n\geq 0$.
\end{proposition}

\begin{proof}
   This is the special case of \Cref{ch2:: prop: hom initial} obtained by taking $\alpha = \id$.
\end{proof}

\begin{proposition}\label{ch2:: prop: structure_maps_weak_equivalences}
    Let $\cF\colon \Int \to \sset$ be a $\cW_{\pi_{0}}$-local object. Then for every $n\geq 0$ and every object $I\in \cJ_{n}$, the structure map
    \begin{equation}
        \underset{J\in ([n]/\Int)}{\mathsf{holim}} \,\, \cF(J) \xlongrightarrow{\ \sim\ } \cF(I)
    \end{equation}
    is a weak equivalence.
\end{proposition}

\begin{proof}
    Since $\cF$ is $\cW_{\pi_0}$-local, in particular it is injectively fibrant and sends all morphisms in $\cJ_n$ to weak equivalences, so the restriction $\cF|_{\cJ_n}$ is a homotopically constant diagram of Kan complexes. By \Cref{ch2:: prop: hom initial} the inclusion $\iota_n\colon \cJ_n \hookrightarrow ([n]/\Int)$ is homotopy initial, hence the canonical map
    \[
      \holim_{([n]/\Int)} \cF \;\xrightarrow{\ \sim\;}\; \holim_{\cJ_n} \cF
    \]
    is a weak equivalence. On the other hand, by \Cref{ch2:: prop: weak contractibility of pi_0} the nerve of $\cJ_n$ is weakly contractible, so we can apply the general result on homotopy limits of homotopically constant diagrams over contractible indexing categories (see \Cref{ch2:: prop: homotopy-limit-homotopically-constant}). This yields, for every object $I\in \cJ_n$, a weak equivalence
    \[
      \holim_{\cJ_n} \cF \;\xrightarrow{\ \sim\;}\; \cF(I).
    \]
    Composing these two weak equivalences gives the desired statement.
\end{proof}

\begin{theorem}
    The Quillen pair 
    \begin{equation}
    \begin{tikzcd}[ampersand replacement=\&]
	{\fun(\bDelta^{\op},\sset)_{\mathsf{inj}}} \&\& {\fun(\Int,\sset)_{\mathsf{inj},h}}
	\arrow[""{name=0, anchor=center, inner sep=0}, "{\gamma^*}", bend left=10, from=1-1, to=1-3]
	\arrow[""{name=1, anchor=center, inner sep=0}, "{\rkan_{\gamma}}", bend left=10, from=1-3, to=1-1]
	\arrow["\dashv"{anchor=center, rotate=-90}, draw=none, from=1, to=0]
    \end{tikzcd}
    \end{equation}
    is a Quillen equivalence.
\end{theorem}

\begin{proof}
    We use the standard criterion for Quillen equivalences which says that it is enough to show:
    \begin{enumerate}
        \item \label{ch2:: item: Quillen_equivalence_conditions_1} The right adjoint $\rkan_\gamma$ reflects weak equivalences between fibrant objects.
        \item \label{ch2:: item: Quillen_equivalence_conditions_2} For every cofibrant object $X$ in $\fun(\bDelta^{\op},\sset)_{\mathsf{inj}}$, the derived unit 
        \begin{equation}
            X \longra \rkan_{\gamma}(\bR \gamma^{*}X)
        \end{equation}
        is a weak equivalence.
    \end{enumerate} 

    To check \eqref{ch2:: item: Quillen_equivalence_conditions_1}, let $\eta\colon \cF \to \cG$ be a natural transformation between fibrant objects in $\fun(\Int,\sset)_{\mathsf{inj}, h}$ and assume that
    \begin{equation}
        \rkan_\gamma(\cF) \xra{\rkan_\gamma (\eta)} \rkan_\gamma(\cG)
    \end{equation}
    is a weak equivalence in $\fun(\bDelta^{\op},\sset)_{\mathsf{inj}}$. Since weak equivalences in $\fun(\bDelta^{\op}, \sset)_{\mathsf{inj}}$ are level-wise, for every $[n]\in \bDelta^{\op}$ the map
    \begin{equation}
        \rkan_\gamma(\cF)([n]) \longra \rkan_\gamma(\cG)([n])
    \end{equation}
    is a weak equivalence in $\sset$.

    Fix $I\in \Int$ and write $[n]=\gamma(I)$. By \Cref{ch2:: prop: structure_maps_weak_equivalences}, the canonical maps
    \[
      \rkan_\gamma(\cF)([n]) \longra \cF(I),
      \qquad
      \rkan_\gamma(\cG)([n]) \longra \cG(I)
    \]
    are weak equivalences. Since the square
    \[
        \begin{tikzcd}[ampersand replacement=\&]
            {\rkan_\gamma(\cF)([n])} \&\& {\rkan_\gamma(\cG)([n])} \\
            \\
            {\cF(I)} \&\& {\cG(I)}
            \arrow[from=1-1, to=1-3]
            \arrow["\sim"', from=1-1, to=3-1]
            \arrow["\sim", from=1-3, to=3-3]
            \arrow[from=3-1, to=3-3]
        \end{tikzcd}
    \]
    commutes, the map $\eta_I\colon \cF(I)\to\cG(I)$ is a weak equivalence by the $2$-out-of-$3$ property. This holds for every $I\in\Int$, so $\eta$ is a level-wise weak equivalence, hence a weak equivalence in $\fun(\Int,\sset)_{\mathsf{inj},h}$.

    To verify \eqref{ch2:: item: Quillen_equivalence_conditions_2}, let $X$ be cofibrant in $\fun(\bDelta^{\op},\sset)_{\mathsf{inj}}$ and choose an injective fibrant replacement
    \[
      \gamma^*X \longra \cF
    \]
    in $\fun(\Int,\sset)_{\mathsf{inj}}$. Since $\gamma^*X$ sends every morphism in $\cW_{\pi_0}$ to an isomorphism, the fibrant object $\cF$ is automatically $\cW_{\pi_0}$-local, hence fibrant in $\fun(\Int,\sset)_{\mathsf{inj},h}$. The derived unit is represented by the adjoint map
    \[
      u_X\colon X \longra \rkan_\gamma(\cF).
    \]
    Fix $[n]\in \bDelta^{\op}$ and choose $I\in\Int$ with $\gamma(I)=[n]$. Then the composite
    \[
      X([n]) \xra{\,u_X([n])\,} \rkan_\gamma(\cF)([n]) \longra \cF(I)
    \]
    is exactly the component at $I$ of the fibrant replacement map $\gamma^*X\to\cF$, hence is a weak equivalence in $\sset$. By \Cref{ch2:: prop: structure_maps_weak_equivalences}, the second map is also a weak equivalence. Therefore $u_X([n])$ is a weak equivalence by the $2$-out-of-$3$ property. Since this holds for every $[n]$, the derived unit is a level-wise weak equivalence, as required.
\end{proof}




\subsection{Imposing the Segal condition}
\smallskip

The next step in constructing a model structure on $\Fun(\Int,\sset)$ which
models the homotopy theory of $(\oo,1)$-categories is to encode the Segal
condition. Concretely, we want to further localize the model structure
$\Fun(\Int,\sset)_{\mathsf{inj},h}$ so that the fibrant objects not only send
$\pi_0$-isomorphisms to weak equivalences, but also satisfy the interval Segal
condition introduced in \Cref{ch2:: subsec: interval_spaces}. The heavy lifting
here is provided by \Cref{ap:: thm: transfer of localizations}.

First, let us recall what we want from an \emph{interval Segal condition}. If an
interval space $\cF\colon \Int \to \sset$ is to model an $(\oo,1)$-category,
then its space of morphisms should be encoded by $\cF\big(\sU(A,B)\big)$, with
the evident maps to $\cF\big(\sU(B)\big)$ and $\cF\big(\sU(A)\big)$ playing the
role of target and source, respectively. Under this interpretation, we expect
the Segal-type condition to express that $\cF\big(\sU(A,B,C)\big)$ captures the
data of composable pairs of 1-morphisms. More explicitly, the Segal-type
condition requires that the canonical map
\begin{equation}
    \cF\big(\sU(A,B,C)\big)
      \longra
      \cF\big(\sU(A,B)\big)
        \underset{\cF\big(\sU(B)\big)}{\overset{h}{\times}}
      \cF\big(\sU(B,C)\big)
\end{equation}
is a levelwise weak equivalence of simplicial sets.

\begin{definition}
    For each $n\geq 2$, let
    \begin{equation}
      \varphi_n\colon G(n) \hookrightarrow \Delta[n]
    \end{equation}
    denote the $n$-dimensional spine inclusion of simplicial spaces, and let
    \begin{equation}
      \mathbb{L}\gamma^{*}(\varphi_n)
    \end{equation}
    be its image under the left derived functor of the left Quillen functor
    $\gamma^{*}$. We write
    \begin{equation}
      \varphi := \{\varphi_n\}_{n\geq 2}
    \end{equation}
    for the resulting family of maps, and similarly
    $\mathbb{L}\gamma^{*}(\varphi)$ for their images.
\end{definition}

\begin{observation}
    Since we are working with the injective model structure on the category of
    simplicial spaces, cofibrations are precisely monomorphisms. In particular,
    each spine inclusion $\varphi_n\colon G(n)\hookrightarrow\Delta[n]$ is a
    cofibration, so its left derived image satisfies
    \begin{equation}
        \mathbb{L}\gamma^*(\varphi_n)
        \simeq
        \gamma^*\big(G(n)\big) \longra \gamma^*\big(\Delta[n]\big).
    \end{equation}
    In other words, the cofibrant replacement is trivial. The same remark
    applies to finite coproducts of spine inclusions, since these are again
    monomorphisms.
\end{observation}

The next result is the key observation of this subsection: it shows that we can
recover the interval Segal condition for interval spaces by localizing
simultaneously on the $\Int$-side and on the simplicial side.

\begin{lemma}\label{ch2:: lem: fibrant Segal condition}
Let $\cF\in \fun(\Int,\sset)_{\mathsf{inj},h}$. Then $\cF$ is
$\mathbb{L}\gamma^{*}(\varphi)$-local if and only if $\cF$ is fibrant and
satisfies the interval Segal condition.
\end{lemma}

\begin{proof}
    Suppose first that $\cF\colon \Int \to \sset$ is
    $\mathbb{L}\gamma^{*}(\varphi)$-local. By definition of locality, $\cF$ is
    fibrant in $\Fun(\Int,\sset)_{\mathsf{inj},h}$, hence injectively fibrant
    and homotopically constant along $\cW_{\pi_0}$. It remains to check that it
    satisfies the interval Segal condition, i.e.\ that for every
    $\sU(A_1,\ldots,A_n)$ with $n\geq 2$, the canonical map
    \begin{equation}\label{eq segal}
        \cF\big(\sU(A_1,\ldots, A_n)\big)
          \longra
        \cF\big(\sU(A_1,A_2)\big)
          \times^{h}_{\cF\big(\sU(A_2)\big)}
        \cdots
          \times_{\cF\big(\sU(A_{n-1})\big)}^{h}
        \cF\big(\sU(A_{n-1},A_n)\big)
    \end{equation}
    is a weak equivalence.

    By $\mathbb{L}\gamma^*(\varphi)$-locality, for every $n\geq 2$ the map
    \begin{equation}
      \map^{h}\big(\gamma^* \Delta[n], \cF\big)
        \xrightarrow{\ \simeq\ }
      \map^{h}\big(\gamma^* G(n), \cF\big)
    \end{equation}
    is a weak equivalence of simplicial sets. Since every object is cofibrant
    and $\cF$ is fibrant, the derived mapping spaces agree with the ordinary
    ones, and we can identify
    \begin{equation}
      \map^{h}(-,-)\;\simeq\;
      \Hom_{\fun(\Int,\sset)}(-,-).
    \end{equation}
    Using the adjunction $\gamma^\ast \dashv \rkan_{\gamma}$, this yields a
    commutative diagram
    \begin{equation}
        \begin{tikzcd}[ampersand replacement=\&]
            {\Hom_{\fun(\Int,\sset)}\big(\gamma^* \Delta[n], \cF\big)} \&\&
            {\Hom_{\fun(\Int,\sset)}\big(\gamma^* G(n), \cF\big)} \\
            \\
            {\Hom_{\fun(\bDelta^{\op},\sset)}\big(\Delta[n], \rkan_{\gamma}(\cF)\big)} \&\&
            {\Hom_{\fun(\bDelta^{\op},\sset)}\big(G(n), \rkan_{\gamma}(\cF)\big)} \\
            \\
            {\rkan_{\gamma}(\cF)([n])} \&\&
            {\rkan_{\gamma}(\cF)([1])\underset{\rkan_{\gamma}(\cF)([0])}{\times}\cdots
              \underset{\rkan_{\gamma}(\cF)([0])}{\times}\rkan_{\gamma}(\cF)([1])}
            \arrow["\sim", from=1-1, to=1-3]
            \arrow["\cong"', from=1-1, to=3-1]
            \arrow["\cong", from=1-3, to=3-3]
            \arrow["\sim",  from=3-1, to=3-3]
            \arrow["\cong"', from=3-1, to=5-1]
            \arrow["\cong", from=3-3, to=5-3]
            \arrow["\sim",from=5-1, to=5-3]
        \end{tikzcd}
    \end{equation}
    where:
    \begin{itemize}
      \item the top horizontal arrow is a weak equivalence by
            $\mathbb{L}\gamma^*(\varphi)$-locality;
      \item the first pair of vertical arrows are the natural isomorphisms of
            the adjunction;
      \item the bottom vertical arrows are the identifications
            $\Hom(\bDelta[k],X)\cong X_k$ coming from the (enriched) Yoneda
            lemma;
      \item the middle horizontal arrow is therefore the Segal map for the
            simplicial space $\rkan_\gamma(\cF)$ at $[n]$, and is a weak
            equivalence by the $2$-out-of-$3$ property.
    \end{itemize}

    Recall that the $n$-dimensional spine $G(n)$ is obtained as the colimit of
    the diagram
    \begin{equation}
        \begin{tikzcd}[ampersand replacement=\&,column sep=small]
            {\Delta[\{0<1\}]} \&\& {\Delta[\{1<2\}]} \&\& \cdots \&\&
            {\Delta[\{n-1<n\}]} \\
            \\
            \& {\Delta[\{1\}]} \&\& {\Delta[\{2\}]} \&\& {\,\Delta[\{n-1\}]}
            \arrow[from=1-1, to=3-2]
            \arrow[from=1-3, to=3-2]
            \arrow[from=1-3, to=3-4]
            \arrow[from=1-5, to=3-4]
            \arrow[from=1-5, to=3-6]
            \arrow[from=1-7, to=3-6]
        \end{tikzcd}
    \end{equation}
    in $\Fun(\bDelta^{\op},\sset)$, so that
    $\Hom_{\fun(\bDelta^{\op},\sset)}\big(G(n), \rkan_{\gamma}(\cF)\big)$ can be
    identified with the iterated fibre product in the bottom right corner of
    the previous diagram.

    Now fix an object $\sU(A_0,\ldots,A_n)\in \Int$ with $n$ interior connected
    components. By \Cref{ch2:: prop: structure_maps_weak_equivalences}, the
    structure maps
    \begin{equation}
      \rkan_\gamma(\cF)([k])
        \longrightarrow
      \cF\big(\text{corresponding interval with $k$ ‘‘slots’’}\big)
    \end{equation}
    are weak equivalences for $k=0,1,n$ and compatible with the simplicial
    structure maps. This yields a commutative diagram of spans
    \begin{equation}
    \begin{tikzcd}[ampersand replacement=\&]
        {\rkan_{\gamma}(\cF)([1])} \& {\rkan_{\gamma}(\cF)([0])} \&
        {\rkan_{\gamma}(\cF)([1])} \& \cdots \&
        {\rkan_{\gamma}(\cF)([0])} \& {\rkan_{\gamma}(\cF)([1])} \\
        \\
        {\cF\big(\sU(A_0,A_1)\big)} \& {\cF\big(\sU(A_1)\big)} \&
        {\cF\big(\sU(A_1,A_2)\big)} \& \cdots \&
        {\cF\big(\sU(A_{n-1})\big)} \&
        {\cF\big(\sU(A_{n-1},A_{n})\big)}
        \arrow[from=1-1, to=1-2]
        \arrow[from=1-1, to=3-1]
        \arrow[from=1-2, to=3-2]
        \arrow[from=1-3, to=1-2]
        \arrow[from=1-3, to=3-3]
        \arrow[from=1-4, to=1-3]
        \arrow[from=1-4, to=1-5]
        \arrow[from=1-5, to=3-5]
        \arrow[from=1-6, to=1-5]
        \arrow[from=1-6, to=3-6]
        \arrow[from=3-1, to=3-2]
        \arrow[from=3-3, to=3-2]
        \arrow[from=3-4, to=3-3]
        \arrow[from=3-4, to=3-5]
        \arrow[from=3-6, to=3-5]
    \end{tikzcd}
    \end{equation}
    in which all vertical arrows are weak equivalences. Passing to the
    corresponding homotopy pullbacks, we obtain a weak equivalence between the
    homotopy pullback computing the Segal map for $\rkan_\gamma(\cF)$ and the
    homotopy pullback on the right-hand side of \eqref{eq segal}. Since the
    former is a weak equivalence by the previous argument, the latter is as
    well. Thus $\cF$ satisfies the interval Segal condition.

    Conversely, if $\cF$ is fibrant and satisfies the interval Segal condition,
    the same reasoning can be run backwards to show that the Segal maps of
    $\rkan_\gamma(\cF)$ are weak equivalences, which is equivalent to
    $\mathbb{L}\gamma^*(\varphi)$-locality.
\end{proof}

\begin{proposition}\label{ch2:: prop: model structure interval Segal space}
    There exists a model structure on the category of pre-interval spaces
    $\fun(\Int,\sset)$, which we will denote by $\mathsf{ISSp}$, with the
    following properties:
    \begin{enumerate}
        \item The cofibrations are the injective cofibrations.
        \item The fibrant objects are interval Segal spaces, i.e.\ functors
              $\cF\colon \Int \to \sset$ satisfying the interval Segal condition \Cref{ch2:: def: interval Segal condition}.
        \item The weak equivalences are the $\mathbb{L}\gamma^{*}(\varphi)$-local
              equivalences. In particular, a weak equivalence between interval
              Segal spaces is a levelwise weak equivalence.
        \item The model structure is left proper, combinatorial, and simplicial.
    \end{enumerate}
    We call this the model category of \textbf{interval Segal spaces}. 
\end{proposition}

\begin{proof}
    Existence of the left Bousfield localization of
    $\Fun(\Int,\sset)_{\mathsf{inj},h}$ at $\mathbb{L}\gamma^{*}(\varphi)$ gives
    (1) and (4), since cofibrations are unchanged and left properness,
    combinatoriality, and simpliciality are preserved. Assertion (2) follows
    from \Cref{ch2:: lem: fibrant Segal condition}, which identifies the
    fibrant objects with interval Segal spaces. Finally, (3) is the general
    description of weak equivalences between fibrant objects in a left
    Bousfield localization: between interval Segal spaces,
    $\mathbb{L}\gamma^{*}(\varphi)$-local equivalences coincide with levelwise
    weak equivalences.
\end{proof}

\begin{theorem}\label{ch2:: thm: Quillen-equivalence-SSp-ISSp}
    The Quillen equivalence previously established between simplicial spaces and
    interval spaces
    \begin{equation}
        \begin{tikzcd}[ampersand replacement=\&]
            {\fun(\bDelta^{\op},\sset)_{\mathsf{inj}}} \&\&
            {\fun(\Int,\sset)_{\mathsf{inj},h}}
            \arrow[""{name=0, anchor=center, inner sep=0}, "{\gamma^*}",
                   curve={height=-16pt}, from=1-1, to=1-3]
            \arrow[""{name=1, anchor=center, inner sep=0}, "{\rkan_{\gamma}}",
                   curve={height=-16pt}, from=1-3, to=1-1]
            \arrow["\dashv"{anchor=center, rotate=90}, draw=none, from=1, to=0]
        \end{tikzcd}
    \end{equation}
    induces a Quillen equivalence
    \begin{equation}
        \begin{tikzcd}[ampersand replacement=\&]
            {\mathsf{SSp}} \&\& {\mathsf{ISSp}}
            \arrow[""{name=0, anchor=center, inner sep=0}, "{\gamma^*}",
                   curve={height=-18pt}, from=1-1, to=1-3]
            \arrow[""{name=1, anchor=center, inner sep=0}, "{\rkan_{\gamma}}",
                   curve={height=-18pt}, from=1-3, to=1-1]
            \arrow["\dashv"{anchor=center, rotate=90}, draw=none, from=1, to=0]
        \end{tikzcd}
    \end{equation}
    between the model category of Segal spaces and the model category of
    interval Segal spaces, upon further localizing along $\varphi$ and
    $\mathbb{L}\gamma^*(\varphi)$, respectively. 
\end{theorem}

\begin{proof}
    This follows directly from \Cref{ch2:: prop: Segal-model-structure} and
    \Cref{ch2:: prop: model structure interval Segal space}, together with
    \Cref{ap:: thm: transfer of localizations}, applied to the Quillen
    equivalence $\gamma^* \dashv \rkan_{\gamma}$ between
    $\Fun(\Delta^{\op},\sset)_{\mathsf{inj}}$ and $\Fun(\Int,\sset)_{\mathsf{inj},h}$.
\end{proof}

\begin{observation}\label{ch2:: obs: complete Segal localization}
If we left Bousfield localize this Quillen equivalence along
$\{J \longra \bDelta[0]\}$ in $\mathsf{SSp}$ and along
$\LL \gamma^* (J \longra \bDelta[0])$ in $\mathsf{ISSp}$, we get a new pair of
model structures whose fibrant objects are complete Segal spaces in
$\mathsf{SSp}$ and $\LL \gamma^* (J \longra \bDelta[0])$-local objects in
$\mathsf{ISSp}$. By the derived adjunction, we have that if $\cF$ is
$\LL \gamma^* (J \longra \bDelta[0])$-local then the induced map
\begin{equation}
    \Map^h(\LL\gamma^* \bDelta[0], \cF) \xra{\,\,\,\simeq \,\,\,} \Map^{h}(\LL\gamma^* J, \cF)
\end{equation} 
is a weak equivalence, which translates to $\Rrkan_\gamma(\cF)$ being a
complete Segal space. We call the resulting model structure $\mathsf{CISSp}$.
From \Cref{ap:: thm: transfer of localizations} we have the following Quillen
equivalence 
\begin{equation}
    \begin{tikzcd}[ampersand replacement=\&]
        {\mathsf{CSSp}} \&\& {\mathsf{CISSp}}
        \arrow[""{name=0, anchor=center, inner sep=0}, "{\gamma^*}", curve={height=-12pt}, from=1-1, to=1-3]
        \arrow[""{name=1, anchor=center, inner sep=0}, "{\rkan_{\gamma}}", curve={height=-12pt}, from=1-3, to=1-1]
        \arrow["\dashv"{anchor=center, rotate=90}, draw=none, from=1, to=0]
    \end{tikzcd}
\end{equation}
\end{observation}


\begin{theorem}[Complete interval Segal spaces]\label{ch2:: thm: complete interval model}
    There exists a model structure on the category of pre-interval spaces
    $\fun(\Int,\sset)$, which we will denote by $\mathsf{CISSp}$, with the
    following properties:
    \begin{enumerate}
        \item The cofibrations are the injective cofibrations.
        \item The fibrant objects are complete interval Segal spaces.
        \item The weak equivalences between complete interval
              Segal spaces are the levelwise weak equivalences.
        \item The model structure is left proper, combinatorial, and simplicial.
    \end{enumerate}
    We call this the model category of \textbf{complete interval Segal spaces}.
    Moreover, restriction and right Kan extension along $\gamma$ induce a
    Quillen equivalence
    \begin{equation}
        \mathsf{CSSp}
        \ \underset{\rkan_{\gamma}}{\overset{\gamma^{*}}{\rightleftarrows}}\
        \mathsf{CISSp}.
    \end{equation}
    In particular, $\mathsf{CISSp}$ presents the $(\oo,1)$-category of
    $(\oo,1)$-categories.
\end{theorem}
\begin{proof}
    The theorem follows from \Cref{ch2:: obs: complete Segal localization}.
\end{proof}

\begin{remark}
    One of the most important results in Rezk's original paper \cite{Rezk01} is the fact that $\mathsf{CSSp}$ is compatible with its Cartesian closed structure, which ensures that for a pair of complete Segal spaces $X$ and $Y$, the internal-hom $X^Y$ is again a complete Segal space. In order to ensure that $\mathsf{CISSp}$ satisfies a similar property, a little more work would be needed, since this is not a property transferred through a Quillen equivalence.
\end{remark}

\begin{notation}
    We denote the full subcategory of $\mathsf{ISSp}$ spanned by the \emph{complete} Segal interval spaces as $\cC\mathsf{ISSp}$.
\end{notation}




\section{Homotopy theory of interval Segal spaces}
\label{ch2:: sec: homotopy theory int}

The theory of $(\oo,1)$-categories' primary goal is to smoothly accommodate category theory with homotopy theory, so before we end this chapter we should articulate the sense in which interval Segal spaces achieve this. To do this, we construct a fundamental invariant of every $(\oo,1)$-category, its \emph{homotopy category}. 

We begin with the canonical definition, obtained by passing from an interval Segal space to its associated Segal space via the homotopy right Kan extension. We then unpack this definition in terms of explicit homotopy limits over $\Int$.

\begin{observation}\label{ch2:: obs: RRkan-gamma-pointwise}
Let \(\cF \colon \Int \longra \sset\) be fibrant in
\(\Fun(\Int,\sset)_{\inj,h}\).
By the construction of the homotopy right Kan extension, for every
\([n]\in \bDelta\) we have a natural weak equivalence
\begin{equation}
  \Rrkan_{\gamma} \cF([n])
  \;\simeq\;
  \hoLimit[([n]\xra{\alpha}\gamma(I))\in [n]/\Int] \cF(I).
\end{equation}
Moreover, by \Cref{ch2:: prop: hom initial} the inclusion
\(\cJ_n \hookrightarrow [n]/\Int\) is homotopy initial, so there is a
natural weak equivalence
\begin{equation}
  \Rrkan_{\gamma} \cF([n])
  \;\simeq\;
  \hoLimit[([n]\xra{\id}\gamma(I))\in \cJ_n] \cF(I).
\end{equation}
For \(n=0,1\) this expresses the \(0\)- and \(1\)-simplices of the
associated Segal space in terms of explicit homotopy limits over \(\Int\).
\end{observation}


\begin{definition}\label{ch2:: def: homotopy category ISSP}
  Let \(\cF\colon \Int \longra \sset\) be an interval Segal space, i.e.\ a
  fibrant object of \(\mathsf{ISSp}\).
  We define its associated \define{homotopy category}
  \begin{equation}
    \ho(\cF) \;\overset{\mathsf{def}}{=}\;
    \ho\big(\Rrkan_{\gamma}\cF\big)
  \end{equation}
  to be the homotopy category of the induced Segal space
  \(\Rrkan_{\gamma}\cF\colon \bDelta^{\op}\longra \sset\), as in
  \Cref{ch2:: def: homotopy category segal space}.
\end{definition}


\begin{remark}\label{ch2:: rem: objects-morphisms-as-holims}
Let \(\cF\colon \Int \longra \sset\) be an interval Segal space.
By Observation~\ref{ch2:: obs: RRkan-gamma-pointwise}, the simplicial set
of objects of the associated Segal space is
\begin{equation}
  \big(\Rrkan_{\gamma}\cF\big)_0
  \;\simeq\;
  \hoLimit[I \in \cJ_0] \cF(I),
\end{equation}
and the simplicial set of \(1\)-simplices is
\begin{equation}
  \big(\Rrkan_{\gamma}\cF\big)_1
  \;\simeq\;
  \hoLimit[J \in \cJ_1] \cF(J).
\end{equation}
The source and target maps
\begin{equation}
  s,t \colon
  \big(\Rrkan_{\gamma}\cF\big)_1 \longra
  \big(\Rrkan_{\gamma}\cF\big)_0
\end{equation}
are induced by precomposition with the maps
\(\langle 0\rangle,\langle 1\rangle \colon [0]\ra [1]\) in~\(\bDelta\), and they can be understood as taking a morphism over $\sU(A,B)$ and then restricting to $\sU(A)$ and $\sU(B)$, respectively — where both maps are given by gluing two of the opens or by covering one of the closed intervals. 


\nid Given \(x,y \in \Ob\big(\ho(\cF)\big)\), the mapping space
\(\Map_{\Rrkan_{\gamma}\cF}(x,y)\) is the homotopy fibre of
\begin{equation}
  \Map^{h}_{\Rrkan_{\gamma}\cF}(x,y) \,\, \overset{\sf def}{=}\,\, \mathsf{hofiber}_{(x,y)}\left(
  \big(\Rrkan_{\gamma}\cF\big)_1 \longra
  \big(\Rrkan_{\gamma}\cF\big)_0 \times
  \big(\Rrkan_{\gamma}\cF\big)_0\right)
\end{equation}
over the point \((x,y)\), exactly as in \Cref{ch2:: cons: objects mapping space}.
\end{remark}




\section{Symmetric monoidal structures}
\label{ch2:: sec: symmetric monoidal structure}

Symmetric monoidal structures are key to the study of field theories and it is thus imperative that we discuss how they will be incorporated into our constructions. There is an entire taxonomy of ways to encode symmetric monoidal structures into homotopical or higher categorical constructions — a profound study of this topic can be found in Lurie's book \cite{LurieHA}. The two main ones, dual via straightening-unstraightening, are as algebras for the $\cE_{\oo}$-operad and as a Segal $\Gamma$-object — we will, given the future applications we have in mind, work only with the latter. 

\subsection{As \texorpdfstring{$\Gamma$}{Gamma}-objects}

Lurie defined what it means to be a commutative monoid in an arbitrary $(\oo,1)$-category, based on ideas dating back to Segal. We expand on that definition in this section and we construct a model structure on a certain category of diagrams in $\cC\mathsf{ISSp}$ whose fibrant objects are the symmetric monoidal ones. Our sources for this section are \cite{CS19}, \cite[Appendix~A]{JFS15}, and of course \cite{LurieHA}.

\begin{definition}
    Let $\Gamma$ be a category with the following data and structure:
    \begin{enumerate}
        \item \textbf{Objects:} Finite sets 
        \begin{equation}
            \langle m \rangle = \{0\,\cdots, m\}, \quad m\geq 0
        \end{equation}
        which we consider to be \emph{based} at 0. 
        \item \textbf{Morphisms:} Morphisms are basepoint-preserving functions 
        \begin{equation}
            f\maps \langle m \rangle \longra \langle k \rangle 
        \end{equation} 
        such that $f(0) = 0$.
    \end{enumerate}
    This category is called the \define{Segal category}. For each $n\geq 0$, let
      \begin{equation}
        i_k \colon \langle n\rangle \longra \langle 1\rangle,
        \qquad 1\leq k\leq n,
      \end{equation}
      be the maps picking out the $k$-th non-basepoint element. 
\end{definition}

\begin{observation}
    $\Gamma$ is the skeleton of $\Fin_*$, the category of based finite sets and basepoint-preserving maps among them.
\end{observation}

In \cite{CS19}, Calaque--Scheimbauer construct a model for the
$(\oo,1)$-category of symmetric monoidal $(\oo,1)$-categories using
$\Gamma$-objects in complete Segal spaces. Concretely, they consider the
projective model structure on
\begin{equation}
  \fun(\Gamma,\mathsf{CSSp})_{\mathsf{proj}}
\end{equation}
and perform a left Bousfield localization at a suitable set of maps
\begin{equation}
  S_{\Gamma} \subset \fun(\Gamma,\mathsf{CSSp})_{\mathsf{proj}}
\end{equation}
encoding the Segal and unit conditions in arity $m\geq 0$; see
\cite[Definition~3.6]{CS19} for an explicit description in terms of the
(co)Yoneda embedding. The fibrant objects in the localized model
structure
\begin{equation}
  L_{S_{\Gamma}}\fun(\Gamma,\mathsf{CSSp})_{\mathsf{proj}}
\end{equation}
are precisely the \emph{symmetric monoidal complete Segal spaces} in the
sense of \cite{CS19}, and the underlying $(\oo,1)$-category of this
model category is the $(\oo,1)$-category of symmetric monoidal
$(\oo,1)$-categories and symmetric monoidal functors.

\smallskip

On the other hand, \Cref{ch2:: thm: Quillen-equivalence-SSp-ISSp}
constructed a simplicial Quillen equivalence
\begin{equation}
  \begin{tikzcd}[ampersand replacement=\&]
    {\mathsf{CSSp}} \&\& {\mathsf{CISSp}}
    \arrow[""{name=0, anchor=center, inner sep=0}, "{\gamma^*}", curve={height=-12pt}, from=1-1, to=1-3]
    \arrow[""{name=1, anchor=center, inner sep=0}, "{\rkan_{\gamma}}", curve={height=-12pt}, from=1-3, to=1-1]
    \arrow["\dashv"{anchor=center, rotate=90}, draw=none, from=1, to=0]
  \end{tikzcd}
\end{equation}
between complete Segal spaces and complete interval Segal spaces. By
\cite[Proposition~A.3.3.6]{HTT}, this induces a simplicial Quillen
equivalence on $\Gamma$-diagrams
\begin{equation}\label{ch2:: eq: Gamma-diagram-Quillen}
  \begin{tikzcd}[ampersand replacement=\&]
    {\fun(\Gamma, \mathsf{CSSp})_{\mathsf{proj}}} \&\& {\fun(\Gamma, \mathsf{CISSp})_{\mathsf{proj}}}
    \arrow[""{name=0, anchor=center, inner sep=0}, "{(\gamma^*)^{\Gamma}}", curve={height=-12pt}, from=1-1, to=1-3]
    \arrow[""{name=1, anchor=center, inner sep=0}, "{\rkan_{\gamma}^{\Gamma}}", curve={height=-12pt}, from=1-3, to=1-1]
    \arrow["\dashv"{anchor=center, rotate=90}, draw=none, from=1, to=0]
  \end{tikzcd}
\end{equation}

Let $\widetilde{S}_{\Gamma} \subset \fun(\Gamma,\mathsf{CISSp})_{\mathsf{proj}}$ be any
set of maps weakly equivalent to the image of $S_{\Gamma}$ under
$(\gamma^*)^{\Gamma}$. By the general
existence theorem for left Bousfield localizations
(Theorem~\ref{ap:: thm: transfer of localizations}) and the invariance
of left Bousfield localization under Quillen equivalence, the adjunction
\eqref{ch2:: eq: Gamma-diagram-Quillen} descends to a simplicial Quillen
equivalence
\begin{equation}\label{ch2:: eq: Gamma-localized-Quillen}
  \begin{tikzcd}[ampersand replacement=\&]
    {L_{S_{\Gamma}}\fun(\Gamma, \mathsf{CSSp})_{\mathsf{proj}}} \&\& {L_{\widetilde{S}_\Gamma}\fun(\Gamma, \mathsf{CISSp})_{\mathsf{proj}}}
    \arrow[""{name=0, anchor=center, inner sep=0}, "{(\gamma^*)^{\Gamma}}", curve={height=-12pt}, from=1-1, to=1-3]
    \arrow[""{name=1, anchor=center, inner sep=0}, "{\rkan_{\gamma}^{\Gamma}}", curve={height=-12pt}, from=1-3, to=1-1]
    \arrow["\dashv"{anchor=center, rotate=90}, draw=none, from=1, to=0]
  \end{tikzcd}.
\end{equation}

Thus the localized model category
$L_{\widetilde{S}_\Gamma}\fun(\Gamma, \mathsf{CISSp})_{\mathsf{proj}}$ presents a
Quillen-equivalent
$(\oo,1)$-category of symmetric monoidal $(\oo,1)$-categories as
$L_{S_{\Gamma}}\fun(\Gamma,\mathsf{CSSp})_{\mathsf{proj}}$.

\begin{proposition}\label{ch2:: prop: fibrant-sym-monoidal-CISSp}
  A diagram
  \begin{equation}
    \cF^{\langle \bullet \rangle} \colon \Gamma \longra \mathsf{CISSp}
  \end{equation}
  is fibrant in the localized model category
  $L_{\widetilde{S}_\Gamma}\fun(\Gamma, \mathsf{CISSp})_{\mathsf{proj}}$ if and only if
  the following conditions hold:
  \begin{enumerate}
    \item For every $m\geq 0$, the value
      $\cF^{\langle m\rangle}$ is a complete interval Segal space (i.e.\
      fibrant in $\mathsf{CISSp}$).
    \item For every $m\geq 0$, the induced map
      \begin{equation}\label{ch2:: eq: Gamma-Segal-map}
        \cF^{\langle m\rangle}
        \longra
        \big(\cF^{\langle 1\rangle}\big)^{\times m}
      \end{equation}
      induced by the product of $\cF^{\lag \bullet \rag}$ evaluated on the inert maps, is a weak equivalence in $\mathsf{CISSp}$.
    \item The value at $\langle 0\rangle$ is essentially terminal: the
      terminal map
      \begin{equation}
        \cF^{\langle 0\rangle} \longra \ast
      \end{equation}
      is a weak equivalence in $\mathsf{CISSp}$.
  \end{enumerate}
\end{proposition}

\begin{proof}
  On the $\mathsf{CSSp}$-side this is exactly the characterization of
  fibrant objects in $L_{S_{\Gamma}}\fun(\Gamma,\mathsf{CSSp})_{\mathsf{proj}}$
  given in \cite[Definition~3.6]{CS19}. The Quillen equivalence
  \eqref{ch2:: eq: Gamma-localized-Quillen} transports this description
  to the $\mathsf{CISSp}$-side.
\end{proof}

We package this as our working definition.

\begin{definition}\label{ch2:: def: symmetric-monoidal-CISSp}
  A \textbf{symmetric monoidal (complete) interval Segal space} is a
  fibrant object of the localized model category
  $L_{\widetilde{S}_\Gamma}\fun(\Gamma, \mathsf{CISSp})_{\mathsf{proj}}$, equivalently a
  functor
  \begin{equation}
    \cF^{\langle \bullet \rangle} \colon \Gamma \longra \mathsf{CISSp}
  \end{equation}
  satisfying the conditions of
  Proposition~\textup{\ref{ch2:: prop: fibrant-sym-monoidal-CISSp}}. We denote this model structure on $\fun(\Gamma, \mathsf{CISSp})_{\mathsf{proj}}$ as $\mathsf{CISSp}^{\otimes}$.
\end{definition}

\begin{remark}
  Via the Quillen equivalence
  \eqref{ch2:: eq: Gamma-localized-Quillen}, the homotopy theory of
  symmetric monoidal (complete) interval Segal spaces is equivalent to
  the homotopy theory of symmetric monoidal complete Segal spaces of
  \cite{CS19}, and hence presents the $(\oo,1)$-category of symmetric
  monoidal $(\oo,1)$-categories. When no confusion can arise, we will
  simply speak of \emph{symmetric monoidal $(\oo,1)$-categories} and
  implicitly model them by symmetric monoidal (complete) interval Segal
  spaces.
\end{remark}


\subsection{Other models for \texorpdfstring{$(\oo,1)$}{(infty,1)}-categories}

It is important to make a note of a handful of different models for
$(\oo,1)$-categories we will be using throughout this thesis.  This entire
chapter is dedicated to proposing a variant model of $(\oo,1)$-categories that
is better suited to construct differential-topological objects --- yet,
throughout the rest of this thesis, other models will be employed and
comparison results between different models will be assumed.  To keep the text
as self-contained as possible, we briefly recall here some of the models we
use and how they are related.

\subsubsection{Quasicategories}

As we have discussed in the introduction to this chapter, an
$(\oo,1)$-category can be understood as a category enriched in
$(\oo,0)$-categories, or $\oo$-groupoids.  By the \emph{Homotopy Hypothesis},
the theory of $\oo$-groupoids, taken up to weak homotopy equivalence, agrees
with the theory of Kan complexes up to homotopy equivalence --- which, by
\Cref{ch2:: thm: homotopy hyp}, is equivalent to the ordinary homotopy theory
of topological spaces.

\begin{definition}\label{ch2:: def: quasicategory}
  A \textbf{quasicategory} $X$ is a simplicial set
  $X_{\bullet} \colon \bDelta^{\op} \to \set$ which satisfies the following
  lifting condition: for every $n \geq 1$, for every $0 < p < n$, and for
  every $v \colon \Lambda^{n}_{p} \to X$ there exists (not necessarily
  uniquely) a \emph{filler} for the commutative diagram
  \[
  \begin{tikzcd}[ampersand replacement=\&]
  	{\Lambda^{n}_{p}} \&\& X \\
  	\\
  	{\Delta^{n}} \&\& {.}
  	\arrow["v", from=1-1, to=1-3]
  	\arrow[hook, from=1-1, to=3-1]
  	\arrow[dashed, from=3-1, to=1-3]
  \end{tikzcd}
  \]
  In other words, $X$ satisfies the \emph{inner Kan condition}.
\end{definition}

\begin{observation}
  A quasicategory $X$ encodes the homotopy theory of its mapping spaces in its
  higher simplices.  Given a pair $x,y \in X_0$ of objects, one defines the
  \emph{mapping space} $\Map_{X}(x,y)$ as the homotopy pullback
  \[
  \begin{tikzcd}[ampersand replacement=\&]
    {\Map_{X}(x,y)} \&\& {X^{\Delta^1}} \\
    \\
    {*} \&\& {X^{\Delta^0}\times X^{\Delta^0} \quad ,}
    \arrow[from=1-1, to=1-3]
    \arrow[from=1-1, to=3-1]
    \arrow["{(d_1,d_0)}", from=1-3, to=3-3]
    \arrow["{(x,y)}"', from=3-1, to=3-3]
  \end{tikzcd}
  \]
  where $X^{K} \coloneqq \Hom_{\sset}(K,X)$ and $d_0,d_1\colon X_1\to X_0$
  are the face maps.  It is a theorem of Lurie that $\Map_X(x,y)$ is always a
  Kan complex; see, for instance, \cite[§2.2.5]{HTT}.
\end{observation}

\begin{observation}
  The theory of quasicategories is especially useful because it offers a very
  compact way of presenting an $(\oo,1)$-category: one constructs a simplicial
  set and checks a collection of lifting conditions.  Depending on the
  $(\oo,1)$-categories one wants to use, this is often tractable in practice.
  Another strong motivation to work with quasicategories is that they admit a
  good notion of $(\oo,1)$-category of functors: for quasicategories $X$ and
  $Y$ there is an internal hom $Y^{X}$, again a quasicategory, which models
  the $(\oo,1)$-category of functors $X \to Y$; see \cite[§1.2.7]{HTT}.
\end{observation}

\begin{remark}
  For every (ordinary) category $\cC$, the (ordinary) simplicial nerve
  $\sN(\cC)$ is a quasicategory.  In this sense, ordinary 1-categories embed
  fully faithfully into the $(\oo,1)$-category of all quasicategories.  We
  will almost always tacitly conflate an ordinary category with its nerve.
\end{remark}

\begin{warning}
  In some parts of the literature this ambiguity is avoided at great cost,
  with authors carefully notating whenever a category is viewed through its
  embedding into quasicategories.  We will not follow this practice except
  when it is genuinely helpful to avoid confusion.
\end{warning}

\medskip

The theory of quasicategories as avatars for $(\oo,1)$-categories is developed
systematically in \cite{HTT}.  In particular, $\oo$-versions of foundational
constructions and theorems of category theory such as the Yoneda lemma,
adjoint functor theorems, and Kan extensions are defined and suitably
generalized in that framework.

\subsubsection{Topological categories}

\begin{definition}
  A \textbf{topological category} is a category $\cC$ enriched in (say)
  compactly generated weak Hausdorff topological spaces.  Thus $\Ob(\cC)$ is
  a set, and for every pair of objects $x,y \in \Ob(\cC)$ there is a mapping
  space
  \[
    \Map_{\cC}(x,y) \in \Top,
  \]
  together with continuous composition maps and units satisfying the usual
  associativity and unitality axioms.
\end{definition}

Starting from a topological category $\cC$, one obtains a simplicial category
by applying the singular complex functor levelwise to the mapping spaces:
\[
  \Sing(\cC)(x,y) \coloneqq \Sing\bigl(\Map_{\cC}(x,y)\bigr) \in \sset.
\]
Lurie defines the \emph{topological nerve} $\sN^{t}(\cC)$ as the coherent
simplicial nerve of this simplicial category:
\[
  \sN^{t}(\cC) \,\coloneqq\, \sN\bigl(\Sing(\cC)\bigr) \;\in\; \sset.
\]

\begin{proposition}
  Let $\cC$ be a (small) topological category.  Then the topological nerve
  $\sN^{t}(\cC)$ is a quasicategory.
\end{proposition}

\begin{proof}[Sketch of proof]
  The singular construction $\Sing(\cC)$ is a simplicial category whose
  mapping spaces are Kan complexes.  The coherent nerve of such a simplicial
  category is a quasicategory; see \cite[§1.1.5]{HTT}.  By definition
  $\sN^t(\cC)$ is exactly this coherent nerve.
\end{proof}

\begin{example}
  Let $\cC$ be the topological category whose objects are CW-complexes and
  whose mapping spaces are the compact-open mapping spaces
  $\Map_{\Top}(X,Y)$.  Then $\sN^{t}(\cC)$ is an $(\oo,1)$-category of spaces,
  equivalent (via a suitable chain of Quillen equivalences) to the
  $(\oo,1)$-category presented by the Kan--Quillen model structure
  $\sset_{\mathrm{QK}}$.
\end{example}


\subsubsection{Simplicially-enriched categories}

\begin{definition}
  A \textbf{simplicial category} is a category $\cC$ enriched in the cartesian
  monoidal category of simplicial sets.  Concretely, this means that for each
  pair of objects $x,y \in \Ob(\cC)$ there is a simplicial set
  $\Map_{\cC}(x,y) \in \sset$, together with simplicial composition maps
  \[
    \Map_{\cC}(y,z) \times \Map_{\cC}(x,y)
    \longrightarrow
    \Map_{\cC}(x,z)
  \]
  and units $*\to\Map_{\cC}(x,x)$ satisfying the usual associativity and unit
  axioms.  A functor between simplicial categories
  $F \colon \cC \to \cD$ is \emph{simplicial} if it induces morphisms of
  simplicial sets
  \[
    \Map_{\cC}(x,y) \longrightarrow \Map_{\cD}\bigl(F(x),F(y)\bigr)
  \]
  compatible with composition and identities.  We write $\mathsf{SC}$ for the
  (ordinary) category of simplicial categories and simplicial functors.
\end{definition}

\begin{remark}
  Applying the singular complex functor to the mapping spaces of a topological
  category $\cC$ produces a simplicial category $\Sing(\cC)$, and conversely
  geometric realization of the mapping spaces of a simplicial category
  produces a topological category.  These two constructions are inverse up to
  homotopy, and they induce a Quillen equivalence between suitable model
  structures on topological and simplicial categories; see, for instance,
  \cite[Ch.~2]{Ber18}.  In practice, we will freely pass between topological
  and simplicial categories without comment.
\end{remark}

A key notion of weak equivalence between simplicial categories is the
following.

\begin{definition}
  A simplicial functor $F \colon \cC \to \cD$ is a
  \textbf{Dwyer--Kan equivalence} if:
  \begin{enumerate}
    \item for every pair of objects $x,y \in \Ob(\cC)$, the induced map
      \[
        \Map_{\cC}(x,y) \longrightarrow \Map_{\cD}\bigl(F(x),F(y)\bigr)
      \]
      is a weak equivalence of simplicial sets; and
    \item the induced functor on homotopy categories
      \[
        \ho(\cC) \longrightarrow \ho(\cD),
        \qquad
        \ho(\cC)(x,y) \coloneqq \pi_0\Map_{\cC}(x,y),
      \]
      is essentially surjective.
  \end{enumerate}
\end{definition}

\begin{theorem}[Bergner]
  There is a cofibrantly generated model structure
  $\mathsf{SC}_{\kan}$ on $\mathsf{SC}$ whose weak equivalences are the Dwyer--Kan
  equivalences.  In this model structure, fibrant objects are precisely the
  simplicial categories whose mapping spaces are Kan complexes (i.e.~categories
  enriched in Kan complexes), and we refer to $\mathsf{SC}_{\kan}$ as the
  \emph{$\kan$-model structure} on simplicial categories.
\end{theorem}

There is a close relationship between this model structure and the Joyal model
structure on quasicategories.

\begin{theorem}
  There exists a Quillen equivalence
  \[
    \mathfrak{C} \colon \sset_{\mathrm{Joyal}}
    \rightleftarrows
    \mathsf{SC}_{\kan} \colon \sN
  \]
  between the Joyal model structure on simplicial sets and the Bergner model
  structure on simplicial categories, where the right adjoint $\sN$ is the
  \emph{coherent nerve} functor.  In particular, both model categories
  present the same underlying $(\oo,1)$-category of $(\oo,1)$-categories.
\end{theorem}

\begin{remark}
  Via Rezk's model structure on complete Segal spaces 
   $\mathsf{CSSp}$ recalled earlier, there is in fact a zigzag
  of Quillen equivalences
  \[
    \sset_{\mathrm{Joyal}}
    \;\simeq_Q\;
    \mathsf{CSSp}
    \;\simeq_Q\;
    \mathsf{SC}_{\kan}
  \]
  presenting the same homotopy theory of $(\oo,1)$-categories; see for
  instance \cite[Ch.~2]{Ber18}.  We will use these comparison results
  implicitly when moving between quasicategories, complete Segal spaces,
  simplicial categories, and the interval Segal spaces introduced in this
  chapter.
\end{remark}

\begin{remark}[Notational convention]
  We will sometimes write a map $\cC \to \cD$ from a topological category
  $\cC$ directly into an $(\oo,1)$-category modelled as a quasicategory or a
  complete Segal space.  By this we always mean the composite
  \[
    \sN\bigl(\Sing(\cC)\bigr) \longrightarrow \cD,
  \]
  where $\Sing$ takes each topological mapping space $\Map_{\cC}(x,y)$ to its
  singular complex and $\sN$ is the coherent simplicial nerve.  This
  convention allows us to suppress a number of intermediate functors from the
  notation.
\end{remark}


\begin{comment}
\subsection{Other models for \texorpdfstring{$(\oo,1)$}{(infty,1)}-categories}

It is important to make a note of a handful of different models for $(\oo,1)$-categories we will be using throughout this thesis. This entire chapter is dedicated to proposing a variant model of $(\oo,1)$-categories that is better suited to construct differential-topological objects — yet, throughout the rest of this thesis, other models will be employed and comparison results between different models will be assumed. To make this as self-contained as possible, we will define here some of the models used and how they relate.


\subsubsection{Quasicategories}

As we have discussed in the introduction to this chapter, $(\oo,1)$-categories can be understood as categories enriched in $\oo$-groupoids or $(\oo,0)$-categories. By the \emph{Homotopy Hypothesis}, the theory of $\oo$-groupoids, when taken up to weak homotopy equivalence, agrees with the theory of Kan complexes, up to homotopy equivalence — which from \Cref{ch2:: thm: homotopy hyp}, is equivalent to the ordinary homotopy theory of topological spaces.


\begin{definition}\label{ch2:: def: quasicategory}
A \textbf{quasicategory} $X$ is a simplicial set $X_{\bullet}\maps \bDelta^{\op} \ra \set$ which satisfies the following lifting condition: for every $n \geq 1$, for every $0 < p < n$, and for every $v\maps \Lambda^{n}_{p} \longra X$ there exists a (non-necessarily unqiue) filler for the following commutative diagram:
\begin{equation}
\begin{tikzcd}[ampersand replacement=\&]
	{\Lambda^{n}_{p}} \&\& X \\
	\\
	{\Delta^{n}} \&\& {.}
	\arrow["v", from=1-1, to=1-3]
	\arrow[hook, from=1-1, to=3-1]
	\arrow[dashed, from=3-1, to=1-3]
\end{tikzcd}
\end{equation}
\end{definition}

\begin{observation}
    A quasicategory $X_\bullet$ encodes the homotopy theory of its mapping spaces in its higher simplices. Given a pair $x,y\in X_0$ of objects, one defines $\Hom_{X}(x,y)$ as the homotopy pullback square 
    \begin{equation}
        \begin{tikzcd}[ampersand replacement=\&]
            {\Hom_{X}(x,y)} \&\& {\Hom_{\sset}(\Delta^1\times \Delta^\bullet,X)} \\
            \\
            {*} \&\& {X_0\times X_0 \quad .} 
            \arrow[from=1-1, to=1-3]
            \arrow[from=1-1, to=3-1]
            \arrow[from=1-3, to=3-3]
            \arrow["{(x,y)}"', from=3-1, to=3-3]
        \end{tikzcd}
    \end{equation}
    It is then a result in Lurie \cite[Theorem 4.2.1.8]{HTT} that this simplicial set is a Kan complex.
\end{observation}

\begin{observation}
The theory of quasicategories is especially useful because it offers a very contained way of presenting an $(\oo,1)$-category: construct a simplicial set, then show it satisfies a collection of lifting conditions; depending on the kinds of $(\oo,1)$-categories one is trying to use, this is often very approachable. Another strong motive for pursuing quasicategories is the fact that it offers an almost instant definition of $(\oo,1)$-category of functors — a trait not shared with all other models.
\end{observation}

\begin{remark}
    For every (ordinary) category $\cC$, the (ordinary) simplicial nerve $\sN(\cC)$ is a quasicategory, and almost always we will tacitly conflate the category with its nerve.   
\end{remark}

\begin{warning}
In some parts of the literature this ambiguity is avoided at great cost, with authors carefully noting whenever a category is seen through its embedding into quasicategories. 
\end{warning}

\medskip 

The theory of quasicategories as the avatars for $(\oo,1)$-categories is fully developed in \cite{HTT}. Namely, $\oo$-versions of foundational constructions and theorems of category theory such as the Yoneda lemma, adjoint functor theorems, and Kan extensions are defined and suitably generalized. 

\subsubsection{Topological categories}

\begin{proposition}
    Let $\cC$ be a topological category. Then the topological nerve $\sN^{t}(\cC)$ is a quasicateogry.
\end{proposition}


\subsubsection{Simplicially-enriched categories}

\begin{definition}
    A \textbf{simplicial category} is a simplicially-enriched category $\cC$, i.e., a category enriched in the category of simplicial sets together with the Cartesian product. Given objects $x,y\in \mathsf{Ob}(\cC)$, there is a simplicial set $\mathsf{Map}_{\cC}(x,y)\in \sset$. We will say a functor between the categories $\cC$ and $\cD$ is simplicial if there is an induced morphism of simplicial sets 
    \begin{equation}
        \mathsf{Map}_\cC(x,y) \longra \mathsf{Map}_{\cC}(F(x), F(y))
    \end{equation}
    compatible with composition and identities. We call the category of all simplicially-enriched categories and simplicial functors $\mathsf{SC}$.
\end{definition}

\begin{remark}
    We want a remark here that talks about?
\end{remark}

\begin{definition}
    I introduce a strange model structure called $\kan$ on simplicial categories, the bottom line is that its fibrant objects are the $\kan$-enriched categories and the weak equivalences are the Dwyer--Kan equivalences.
\end{definition}

We might sometimes make the conscious decision to write a map $\cC \longra \cD$ from a topological category to an $(\oo,1)$-category, by which we always mean $\mathsf{Sing}\cC \longra \cD$ where $\mathsf{Sing}$ passes all the topological mapping spaces $\mathsf{Map}_{\cC}(x,y)$ to their singular complexes and then, if we are modelling $\cD$ as a quasicategory or a complete Segal space, finally applying the (coherent) simplicial nerve $\sN$.  
\end{comment}


